% ============================================================
%  H. Brøchner, BENEDICT SPINOZA. EN MONOGRAPHIE.
%  Kjøbenhavn: P. G. Philipsens Forlag (Thieles Bogtrykkeri), 1857.  188 pp.
%  "Benedict Spinoza. A Monograph" -- TRANSLATION (English).
%
%  GENERATED by ebook-method/brochner-spinoza/tr_parts.py from the parts in
%  texts/brochner/spinoza/.parts/tr/ (six, cut at chapter heads, in the print's order);
%  edit the parts and rebuild, or edit this file only after the last build.
%  Translated from transcription.tex (Danish) in this directory, the text of record
%  (fully collated 2026-09-23; see its header), as it stood at commit db09964
%  (2026-10-06).  If the transcription is repaired later, re-check the English at
%  every changed reading (TRANSLATION-PLAYBOOK.md §6a).
%
%  DRAFT COMPLETE (2026-10-07): the whole book -- Contents, Preface (pp. 1-14), the
%  half-title (15-16), Introduction and Chapters 1-10 (17-188) -- with every \opage of
%  the Danish once and in order, and the footnotes and counter resets page by page as
%  in the Danish (tr_parts.py checks both).  Made by six translator agents working
%  from ebook-method/brochner-spinoza/TRANSLATOR-BRIEF.md, one part each, cut at
%  chapter heads.  The translators' 30 odd spots in the Danish were read at the scan
%  (spotsheets.py): one transcription error found (p. 169 `(V., 36. schol.)', print
%  `(V, 36. schol.)'; the English follows the print), the rest the print's own.  The
%  transcription was repaired accordingly on 2026-10-07 (its header §10): that one
%  correction and 17 `as printed' comments; the English needed no change.
%  Not yet read by an independent reviewer.
%
%  CONVENTIONS (beyond TRANSLATION-PLAYBOOK.md §2):
%   * \opage{N} at the English word where printed page N begins; footnotes numbered
%     per printed page as in the print (1), 2), ...), hence \setcounter{footnote}{0}
%     at the marker of every page that carries a note, exactly as in the Danish.
%   * Letterspacing in the Danish (\emph) -> \emph.  LATIN IS ROMAN, as in the print
%     (an all-antiqua book, which sets Latin, French and German in the body type; see
%     the transcription's header §3); it is not italicized here either.  \textit only
%     where the print has a real italic sort (the Greek of p. 63; pp. 69, 70, 171).
%   * Quotations Brøchner gives in Latin, German or French stay as printed; what he
%     gives in Danish is translated from his Danish.  »...«, „...“ -> ``...''.
%     "ɔ:" / "d. v. s." -> "i.e."; Ethics and letter citations as printed (Eth. II, 11.
%     Coroll.; Ep. 6; Tract. theol. pol.); "Jvfr." / "Smlgn." -> "Cf."
%   * Misprints the transcription records as printed: the sense is translated and the
%     print is named in a comment at the site.
%   * The edition sets the chapter-closing rules single throughout (transcription
%     header §8); so does this.
%  TERMS (Spinoza's vocabulary as Brøchner renders it): Substants = substance;
%  Attribut = attribute; Modus/Modi = mode(s); Aand = mind where it is Spinoza's
%  mens (Chapter 5, "Den menneskelige Aand" = "The Human Mind"), spirit where it is
%  Hegel's or Christianity's Geist; Erkjendelse = knowledge (cognition where the act
%  is meant); Forestilling = representation; Idee = idea; Begreb = concept; Væsen =
%  essence (being, where an entity is meant); Tilværelse = existence; Udstrækning =
%  extension; Tænkning = thought; Affect = affect; Lidelse = passion (passio);
%  Trældom = bondage; Begjær(lighed) = desire; Forstand = intellect (intellectus);
%  Fornuft = reason; Indbildningskraft = imagination; Magt = power; Dyd = virtue;
%  Naturret = natural right; Selvopholdelse = self-preservation.
% ============================================================
\documentclass[12pt,a4paper,openany]{book}
\usepackage[utf8]{inputenc}
\usepackage[T1]{fontenc}
\usepackage{libertinus}
\usepackage{libertinust1math}
\usepackage{textalpha}   % the Greek of p. 63, typed directly
\usepackage[english]{babel}
\usepackage[protrusion=true,expansion=true]{microtype}
\usepackage{enumitem}
\usepackage{setspace}
\usepackage[margin=1.2in]{geometry}
\usepackage{fancyhdr}
\setlength{\headheight}{28pt}
\addtolength{\topmargin}{-13.5pt}
\usepackage{hyperref}
\hypersetup{pdftitle={Benedict Spinoza. A Monograph (1857)}, pdfauthor={H. Brøchner}, hidelinks}
% The print numbers its notes 1), 2), ... restarting on every printed page.
\renewcommand{\thefootnote}{\arabic{footnote})}
\pagestyle{fancy}
\fancyhf{}
\fancyhead[LE,RO]{\thepage}
\fancyhead[RE]{\textit{Benedict Spinoza. A Monograph}}
\fancyhead[LO]{\textit{\leftmark}}
\fancypagestyle{plain}{\fancyhf{}\fancyhead[LE,RO]{\thepage}}
\setstretch{1.3}
\usepackage{marginnote}
\newcommand{\opage}[1]{\marginnote{\footnotesize\textit{[#1]}}}
\begin{document}
\begin{titlepage}\centering\vspace*{2cm}
{\Huge BENEDICT SPINOZA}\\[2ex]\rule{3em}{0.4pt}\\[2ex]{\Large A Monograph}\\[3ex]
{\large by Dr.\ H.\ Brøchner}\par\vfill
{\small Translated from the Danish}\\[0.6em]
{\small \textit{Benedict Spinoza. En Monographie.} Kjøbenhavn: P.\ G.\ Philipsens Forlag, 1857}
\end{titlepage}
\chapter*{Contents}
\addcontentsline{toc}{chapter}{Contents}
\markboth{Contents}{Contents}

\begingroup
\setstretch{1.15}
\begin{list}{}{%
  \setlength{\leftmargin}{3.6em}\setlength{\itemindent}{-3.6em}%
  \setlength{\itemsep}{2pt}\setlength{\parsep}{0pt}}

\item[] \hfill Page

\item \textbf{Preface} \dotfill 1---14

\item[] \vspace{2pt}

\item \textbf{Benedict Spinoza} \dotfill 15---188
\item \textbf{Introduction} \dotfill 17---26
\item \textbf{Chap.\ 1.\ Spinoza's Life and Writings} \dotfill 27---43
\item \textbf{Chap.\ 2.\ The Fundamental Concepts of Spinozism} \dotfill 44---64
\item \textbf{Chap.\ 3.\ The Immediate Consequences of Spinoza's Conception of the
  Concept of Substance} \dotfill 65---71
\item \textbf{Chap.\ 4.\ Spinoza's Concept of Nature} \dotfill 71---82
\item \textbf{Chap.\ 5.\ The Human Mind} \dotfill 82---108
\item \textbf{Chap.\ 6.\ Man as naturæ pars. The Affects. Man's
  Bondage} \dotfill 108---132
\item \textbf{Chap.\ 7.\ Natural Right. Religion} \dotfill 132---161
\item \textbf{Chap.\ 8.\ Man's Freedom and Eternity} \dotfill 162---171
\item \textbf{Chap.\ 9.\ Characterization and Critique of Spinozism} \dotfill 171---181
\item \textbf{Chap.\ 10.\ Spinoza's Relation to Later Philosophy} \dotfill 181---188

\end{list}
\endgroup

\newpage

\chapter*{Preface}
\addcontentsline{toc}{chapter}{Preface}
\markboth{Preface}{Preface}

% --- p. 1 ---
\opage{1}The presentation of the Spinozist philosophy which I have given in
this work is intended to begin a series of monographs which, while
each of them forms a whole by itself, shall at the same time, taken
together, give a survey of the whole history of philosophy. I have not
wished to attempt to treat that history with equal thoroughness in its
entire connection; I believe that the powers of a single person are
insufficient for such a work, when it is to be grounded on a real study
of the sources, on a reliable autopsy; at least I feel that my own
powers would be. I also assume that at the present moment, after the many
attempts that have been made in this direction in the last 30 years,
with greater and lesser skill and success, it would be less needed,
and that what matters now is rather, through a new and careful treatment
of the individual main points, to procure a more reliable material for
a renewed comprehensive treatment of the whole subject. A fragmentary
treatment which nevertheless, as from the commanding heights, gives an
overview of the whole domain will perhaps also, in that it concentrates
the gaze on what is essential, on the historically unfolded fundamental
forms of thought, which are not something merely past, belonging to
memory and to learned research, but something that within its limit is
constantly justified and abiding, that retains its significance for all
times, and in that it presents the systems of the individual thinkers
in such
% --- p. 2 ---
\opage{2}detail that a deeper insight can be opened into
their inner mainsprings, so that they are seen as the unfolding of a
\emph{living} thought, be better fitted than the connected presentation
to pave the way, not only for a more thorough understanding, but also,
through it, in some, for some interest in a subject to which the mood and
striving of the moment are not favorable. No one who has made
philosophy the occupation of his life can well have escaped doubt and
despondency at seeing the science which a few decades ago was the
object of an enthusiastic, hopeful study become, through the reaction
of disappointment, an object of indifference, and almost hidden beneath
alien, often idea-less interests. But as the self-reflection of the
human spirit, philosophy is its inalienable possession, which, once
won, will never be given up; it is like those rivers which over a
stretch of their course disappear under the earth, but only to come
into view again, mightier and richer after their union with invisible
springs.

In making it my task, through a fragmentary presentation of the history
of philosophy, to give a survey of it in its whole extent, my choice has
had to fall on those systems which mark high points, from which there
opens at once a view back over a longer, completed development, and
forward over a new series, proceeding from the system presented or
substantially influenced by it. With this consideration \emph{this}
coincides: to select such systems as by their principle and direction
represent main forms of human thought; for only such stand in that
relation to past and posterity. Greek philosophy will be represented by
\emph{Plato}, \emph{Aristotle} and \emph{Plotinus}, the Renaissance by the
\emph{Platonic Academy in Florence} and \emph{Giordano Bruno}; modern
philosophy, in its one line of development by \emph{Bacon} and the
\emph{Système de la nature}, in its other by \emph{Spinoza} and
\emph{Leibnitz}, and at its last stage of development by \emph{Kant}
and \emph{Hegel}. By means of introductions and
% --- p. 3 ---
\opage{3}concluding essays to the
individual presentations the connection between the fragmentary parts
will be established.

The complete justification of the choice of the representative systems
must be given by the presentation of them itself. Yet I must here,
briefly, and more by way of hint than of elaboration, set out the
conception of the course of the history of philosophy that underlies the
choice, and this all the more since I begin with one of the systems in
the middle of the series, so that in the presentation of it there will
be no natural place for such a development.

In my view the whole history of philosophy, from its first beginning
until our own day, divides into only \emph{two} great, \emph{completed}
periods: Greek philosophy and that of modern times. They are separated
by a great interval; perhaps a similar one will come to lie between the
philosophy of our time and that of the future, which will rest on a
richer development of the individual sciences, especially the natural
sciences. \emph{Greek} philosophy I regard as the beginning of
philosophy; for the Orient knew only the infinity of fantasy, not that
of thought. The philosophical work of Greece was, with the nature-life
of the Orient as its presupposition, to discover the \emph{idea}, to
grasp the universe in intuition as the infinite, living, harmonious
\emph{work of art}, which received its true reality through form, the
idea, the eternally existing plastic thought. The thought of
\emph{beauty} was the fundamental thought of Greece; it governed Greek
thought as it governed Greek life, life in the state and the moral life
of the individual. The dissolution of Greek thought is grounded in the
inner contradiction in beauty's \emph{immediate} union of its elements,
form and material, in the impossibility of making the material, matter,
be wholly absorbed in form, of making ``the non-being'' vanish over
against being, which has all reality in itself but constantly needs the
material in order to realize this reality. It is this inner
contradiction of which Greek thought after Aristotle becomes conscious
and which it vainly seeks to resolve within its own domain, until in
Neoplatonism it raises itself above its
% --- p. 4 ---
\opage{4}foundation, nature, in order to find in a higher region a
reconciliation which, given its peculiar character, it could not find
there, but carries the contradiction with it over into the new element,
and, in the service of a new spirit, exhausts its last strength in a
fruitless attempt. For this development the natural representatives are
to be sought in \emph{Plato}, who gathers up in himself the whole older
Greek philosophy, first utters the fundamental thought of Greece, and
first draws the universe under the dominion of the idea; --- in
\emph{Aristotle}, who determines the last genuinely Greek formula for the
unity of form and matter, who by means of it supplements Plato, who,
while Alexander subjugates a world for Greece, wins the riches of a world
for Greek thought, and with them gives nourishment to a millennium; ---
and finally in \emph{Plotinus}, the greatest Greek thinker in the time
when Greece had become a stranger to itself, when a new principle of
life was already moving the world, and the last energetic attempt of
Greek thought to assert itself only showed its impotence and heralded
the victory of the new age.

The whole period from the fall of Greece to the rebirth of the sciences
occupies no place in the history of philosophy. I believe that the more
one returns, unswayed by unclear tendencies, to the study of this time,
the more decidedly one will come to the insight that all the
\emph{thought-}content on which the Christian Middle Ages fed is owed
to Greek philosophy, and that one can even extend this proposition so
that in everything essential it holds for the Renaissance as well.
Christianity could not and would not develop a philosophy out of itself;
it had denied the powers from which such a philosophy would have to
develop, abandoned the ground from which it would have to spring up. In
that Christianity wishes to accomplish what Greece had not been able to
accomplish because it rested in the thought of beauty: to tear spirit
loose from nature as the one true and real, to posit nature as the
essenceless, the evil, to demand death as the condition of life, to move
reality over into a kingdom that ``is not of this world'', it places
itself in a merely negative relation
% --- p. 5 ---
\opage{5}to the forms and powers of human and natural life; its demand
remains always only this: to deny them. There can therefore be no talk,
in the \emph{strict sense}, of a Christian philosophy, any more than of a
Christian art, a Christian state, a Christian ethics. It is a misuse of
words to speak so; one does it only by overlooking, intentionally or
unintentionally, that Christianity did not come into a \emph{new} world,
but that it came into a world in which a great period of culture was
just coming to its close, that it found there the whole rich inheritance
of antiquity, and that, when the consistent hope of the first Christian
age of an imminent annihilation of natural reality was disappointed, it
had to accommodate itself to the \emph{alien} reality, speak its language
and use its riches, yet without ever being able or willing to enter into
a unity with it; so that everything in the way of thought, of art, of
forms for an ethical human life, that has developed \emph{simultaneously
with} Christianity has not proceeded \emph{from it}, but only
\emph{under its influence}, frequently in struggle against it, from the
natural ground of human life, fertilized by the thought of antiquity.

It was only in the fifteenth century that thought, which had gained
strength in its bondage, ventured, inspired by the new revelation of
antiquity, and borne up by a richly unfolded secular life, and by a
nature no longer alien to spirit, to develop a philosophy independent of
the authority of the Church. It was still a period of transition; thought
still had to struggle with fantasy for dominion, and did not have the
strength to develop a new content out of itself alone; but in the
philosophy of antiquity it found a content in which the human spirit
recognized itself in its reality, and with enthusiasm it seized this
content and formed it anew. It was the bacchantic period of thought:
with an enormous confidence, with an unshakable \emph{faith in man}, it
reveled in the new riches, and formed new systems, in adherence to a
\emph{self-chosen} model. The thought-content of these systems is that of
antiquity, as it had been
% --- p. 6 ---
\opage{6}prepared through the last Greek philosophers: it was a grandiose
pantheistic view of the divine universe, in which man held an exalted
place as the point of union between the spheres, as a microcosm that
mirrored the universe in itself. The character of the systems was still
predominantly theological, but different from that of scholasticism.
Nature exerted a stronger attraction on spirit, and the view of nature
pushed itself, consciously or unconsciously, more and more in under the
view of the divine. The universe became a harmonious unfolding of the
life of the Godhead, and faith in the God-man transformed itself into
faith in the divine, microcosmic human spirit.

But as thought was thus drawn toward nature, which throughout the Middle
Ages had been banished, and at most had been only an object for
fantasy, the eye had, on the one hand, to be opened to the natural
conditions under which the human spirit stands; on the other hand, the
struggle between the new direction and the inherited one, which still
kept a secret power, had to lead to a dualistic separation between the
natural and the supernatural, between body and spirit, between a
sensualist and an idealist theory of knowledge, between a materialist
and a spiritualist conception of the principles of existence. Through
this dualism, which was bound to bring dissatisfaction and skepticism
with it, attention was directed to the necessity of a firm starting point
for thought, a secure method for its development, and at this point
begins \emph{the philosophy of modern times}, which breaks off the
tradition of antiquity and seeks, independently and self-consciously, to
unfold a new content of thought.

As an inheritance from the Middle Ages the new philosophy had received
the presupposition of a dualism between spirit and nature, a dualism
which in principle was given with Christianity. For thought, which
independently sought after truth, there lay two separate spheres, % print: efter, Sandhed (stray comma)
as it were two worlds; but in the unity of self-consciousness lies a
protest against every dualism, a demand for the unity and
% --- p. 7 ---
\opage{7}coherence of knowledge. But the dualistic presupposition must
determine this unity in such a way that one of the opposites swallows up
the other, that in the union the peculiarity of the one is given up. This
will prove to be the case in the systems that mark the end points of
this philosophical epoch. And the consequence that is drawn in them is
the conscious or unconscious mainspring of the earlier systems, which
rest on the same fundamental presupposition; by the final goal one
understands their striving.

The first form under which dualism appears is that philosophical
knowledge takes a double starting point, from nature or from spirit, and
this double starting point, by which the character and essential object
of knowledge are determined in essentially different ways, constitutes
for modern philosophy in its progressive development a double series,
which, however, by the nature of the case, was bound at once to acquire
points of contact, since both proceeded from the presupposition of the
power of spirit to know, since nature retained its mighty power of
attraction for both, and since both likewise began by borrowing their
method from the natural sciences and mathematics.

\emph{Bacon} and \emph{Descartes} mark the double starting point for
thought. \emph{Bacon} posits \emph{experience}, directed toward nature and
its real fullness, as the principle of knowledge; he makes nature the
essential object of thought, the source of truth, and natural philosophy
the mother of the sciences. But he will not let experience lose itself in
the particular; he wants to find the creative forms, the unities of
nature, the universal in the concrete, and he acknowledges the activity
of thought and the demand of idealism. But in the principle of
experience lies the danger of the consequence that the I behaves
passively, that, as merely receiving through the senses, it receives only
the finite, that this finite, conceived strictly as finite, is conceived
mechanically, and that mechanism is carried over to spirit as well. The
premises for all these consequences are already given in \emph{Hobbes};
in part the consequences are also drawn, and the
% --- p. 8 ---
\opage{8}final goal of this direction of knowledge hinted at; but a
remnant of a rationalist element in him, and the fact that his interest
turned essentially toward the practical, prevented him from becoming for
the development of philosophy what \emph{Locke} later became.

While empiricism thus quickly developed its consequences, the opposite
direction, which seeks to develop the world of thought out of the
inwardness of the thinking I, had taken definite shape in significant
systems. \emph{Descartes} sought, through doubt about everything, through
the separation of everything heterogeneous, back to the unshakable
presupposition of doubt, the I, which was determined as one with thought
and as sundered from matter and extension. But in order to obtain a
content he must leave this point of unity; he must take refuge in that
from which the I has separated itself as its opposite, and which it
cannot even \emph{think} as its opposite without cancelling the abstract
relation of opposition. Through God thought comes back to the material;
the soul acknowledges, as if under compulsion, its connection with the
body, even though this connection is reduced to a minimum of space and to
the most abstract form of influence. But the connection of mind and body
is again conceived as a \emph{composition} of two heterogeneous things;
the thinking and the extended substance remain separate, and
\emph{Geulincx} and \emph{Malebranche} draw the consequences of
occasionalism. The divine substance in Descartes remains standing outside
reality, separated from the two substances that fill reality, and akin
only to the thinking substance. Thought's striving for the unity of
knowledge is thus hampered at two points. But \emph{Spinoza} removes the
ambiguity in the concept of substance (cf. the Introduction): the
infinite unity of the divine substance fills everything, and thereby the
determinate substances are posited as its attributes. Thought and
extension are posited in the unity of parallelism as equally justified,
and the preponderance is even on the point of passing over from thought
to matter, God is on the point of coinciding with nature; but the
attribute of thought then reaches out again beyond extension, nature is
again conceived through an
% --- p. 9 ---
\opage{9}abstraction of thought, and thought asserts the significance that
it had as starting point. --- In opposition to Spinoza's pantheistic
absorption of finitude, \emph{Leibnitz}, within the same fundamental
direction, holds fast to the interest of the age, individuality, and by
his doctrine of monads asserts the subsistence of the individual, while
against empiricism he asserts the right of idealism, the validity of the
a priori; and the theological background of his system preserves the
element of infinity and strives toward unity. In \emph{Wolf} his
philosophy was reduced from speculative depth to the clarity of the
understanding, was systematized and diluted, and thereby made fit to
pass into the common consciousness in the form in which we find it again
in the \emph{German popular philosophy}, the apotheosis of the
understanding, where the finite, spiritless I, the parody of Leibnitz's
monad, became the center of the universe, indirectly and unconsciously
made into an absolute, while the absolute vanished into the relative, the
infinite into abstractions, so that finitude fills the whole territory
of philosophy.

At this stage of development the idealist direction meets, at the
boundary of the unphilosophical, with French materialism, which was the
consistent development of the Lockean philosophy.

In \emph{Locke} philosophy had returned to self-consciousness; but not,
as in Descartes, in order to develop out of the I a content that lay
originally in it, but in order to demonstrate the passivity of the I in
relation to the object, in order to approach the unity of thought and
being by reducing the former to the passive copy of the latter. The
inconsistencies that still remained in the Lockean philosophy were
removed as it passed over to France. \emph{Condillac} drew from it the
consequence of sensualism, and in the French materialists the concept
of spirit vanished, and the difference between thought and the material
object of thought became only a difference of degree within mechanically
moved matter. With this development the last consequences of empiricism
had been drawn; the opposites had been reduced to unity, but in such a
way that the one was annihilated. A corresponding
% --- p. 10 ---
\opage{10}development had not been reached by the contemporary idealism.
The idealism of \emph{Berkeley} was only formal; under another name it
kept the whole content of realism. The \emph{Scottish} philosophy rested
on a postulated \emph{experience}; the \emph{German} was idealism only
in abstracto, and aimed unconsciously at the goal of the opposite
direction. \emph{Rousseau}, finally, did indeed point to a definite
starting point for a more vigorous idealism; but without clarity and
without carrying it through.

At the close of the eighteenth century everything was prepared for a
decisive turn in philosophy. Materialism had drawn its ultimate
consequence; idealism was languishing, but was not exhausted; in its
abstract presupposition there lay an incitement to a new form, and
dissatisfaction with its content was bound, in one equal to the task, to
call forth the consciousness of what lay at its ground and from that to
draw out a development. The way out that still stood open to thought was
to make all actual being into \emph{thought being}, to posit thought as
the source of all reality, and out of the I, which is one with the
absolute and with the true essence of nature, to develop a system that
appears at once as the system of thought and of reality. It is this
direction that German philosophy took at the end of the last century,
and which culminated in \emph{Hegel}. It is, despite all its difference
from, indeed its opposition to, Christianity, the last philosophical
consequence of Christianity's conception of the relation between the
spiritual and the natural. \emph{Not Berkeley, but Fichte and Hegel are
the opposites and the parallels of abstract materialism.}

In \emph{Kant}, in whose doctrine of \emph{transcendental apperception}
this direction is given as in a germ, it still appears limited and bound
by its opposite. Kant does not rest on a single one of the contemporary
directions, but on both; in accordance with the whole tendency of his
time, he proceeds from a critique of consciousness, but in consciousness
he acknowledges both the active and the passive moment, and thus asserts
the right of idealism while acknowledging the
% --- p. 11 ---
\opage{11}significance of empiricism. He does not yet put everything
objective into self-consciousness, teaches as yet no monism of thought;
on the contrary, it is precisely with him that the opposition between
thought and being is made the object of philosophy; but objectivity is
thinned out for him into an abstract \emph{Ding an sich}, which at every
moment threatens to coincide with the most abstract form of the I. This
consequence \emph{Fichte} drew, in that the I, as practical, was posited
as producing the not-I; and \emph{Schelling} removed the ambiguity that
lay in the determination ``the I'', by conceiving it as the absolute I,
identical with nature intuited under the forms of self-consciousness,
whose identity in the ideal and the real series of potencies it became
the task of philosophy to grasp. The method that was to realize this
demand was set forth by \emph{Hegel}. For him pure logical thought
becomes all reality, its eternal movement within itself the source of
all reality, nature its stepping out of itself, spirit its return to
itself, the method the inner movement of thought itself, and therefore
one with the system of philosophy, and this system equal to reality.

With Hegel a whole series of development stands completed; what is later
than he is only endeavor that does not yet belong to history. \emph{His
system on the one side, French materialism on the other, represent the
end points of the doubly developed philosophy whose presupposition is
the dualism of the Middle Ages.}

The choice of the systems that are to represent the philosophy of modern
times will have found its justification in what has been set out.
\emph{Bacon} and the \emph{Système de la nature} mark the beginning and
the end of a consistently progressing series that goes from empiricism
to materialism. \emph{Spinoza} represents the beginning of the opposite
series because he was the first to state consistently what Descartes only
hinted at. \emph{Leibnitz}, as the one who philosophically integrates
Spinoza and contains in himself the fruitful germ of a subsequent
development, represents the conclusion of idealism's
% --- p. 12 ---
\opage{12}first main stage (as dogmatism). \emph{Kant} and \emph{Hegel}
enclose the most recent philosophy.

For the period of transition the choice has been more difficult, and the
isolation of its systems, the manifold and incoherent character of their
tendencies, will always make the preponderance of any single system
conditional. But \emph{the Platonic Academy} in Florence represents, as a
collective individual, the first enthusiastic adherence to the
rediscovered antiquity, while \emph{Giordano Bruno} most vigorously
develops and most fully sums up the thought-content of his time, in a
form that bears the stamp of a spirit moved by the passion of thought,
in which poetry and philosophy, the great powers of that time, contend
for dominion. ---

In presenting the individual systems my position will be essentially
that of historical characterization, and critical only insofar as
inconsistencies in the carrying out show themselves in the course of the
reproducing development, especially in failed attempts to take up what is
excluded by the one-sidedness or deficiency of the principles. It is thus
essentially \emph{the execution of the fundamental thought}, not the
fundamental thought itself, that becomes the object of criticism. The
criticism of the fundamental thought itself is, for all the earlier
systems, given in the historical development; the task with the
individual system is only to bring that thought out in its whole
peculiarity and determinacy, and thereby to demonstrate its historical
justification as a link in the development of thought. That criticism
would be unjustified which approached the systems with a demand for the
solution of definite problems that the systems have not set themselves.
If a system could succeed in closing itself up into a real whole without % print: selv, Kunde (comma for full stop)
particular problems coming forward, that would be a sign that the
problems were falsely posed, that the task would be, not to solve them,
but to dismiss them.

In order to bring out the fundamental thought completely, it will be
necessary to include everything in the detailed execution that gives an
insight into the inner motives of the system, or throws a more definite
light back on the principle. Even what
% --- p. 13 ---
\opage{13}\setcounter{footnote}{0}is connected with the principle only by
looser threads, but has acquired a peculiar historical significance, may
not be missing from the presentation.

I have begun the series of monographic presentations with Spinoza, partly
for the external reason that the material here lay ready in the greatest
completeness, partly for the internal one that is given in the
importance of his system and in the great influence it has had on the
development of philosophy in this century. I have had all the less
hesitation in beginning in the middle of the series since each individual
monograph constitutes a whole by itself, which will have its
completeness even without the others.

In the presentation of Spinoza I have seldom mentioned the earlier
literature. Those who are familiar with it will soon see that account
has been taken of everything in it that is even moderately significant,
including what belongs to the most recent time\footnote{Besides the
well-known works on the history of philosophy, partly in its whole
extent, partly for shorter periods, by \emph{Brucker}, \emph{Tennemann},
\emph{Tiedemann}, \emph{Buhle}, \emph{Hegel}, \emph{Schleiermacher},
\emph{Ritter}, \emph{Feuerbach}, \emph{Erdmann}, K.\ \emph{Fischer}
(Gesch.\ der neueren Philosophie Bd.\ I.), \emph{Cousin} (fragments de
philosophie Cartésienne), and \emph{Bouillier} (histoire de la philos.\
Cartésienne. Paris 1854.\ 2.\ Vol.), as well as the monographs by
\emph{Sigwart} (der Spinozismus historisch und philosophisch erläutert
mit Beziehung auf ältere und neuere Ansichten), \emph{Schaarschmidt}
(Descartes und Spinoza), \emph{Helfferich} (Spinoza und Leibnitz),
\emph{Montbeillard} (de l'Ethique de Spinoza), \emph{Erdmann}
(Vermischte Aufsätze), and \emph{Trendelenburg} (Ueber Spinozas
Grundgedanken und dessen Erfolg), I have used what is to be found of
critical and polemical developments of Spinozism in \emph{Leibnitz}
(especially in the work edited by \emph{Foucher de Careil}: réfutation
inédite de Spinoza par Leibnitz), \emph{Wolf}, \emph{Jacobi},
\emph{Herder}, \emph{Mendelssohn}, \emph{Maimon}, \emph{Schelling},
\emph{Hegel}, \emph{Baader}, \emph{Herbart} and \emph{Trendelenburg}.
The writings of Spinoza's contemporary theological and philosophical
opponents have been accessible to me only in part. They yield profit
only through their misunderstandings.}. I have thought that the right
way of using that literature, and the one best serving the unity of the
historical presentation, would be, after having used it as an aid to
% --- p. 14 ---
\opage{14}avoid misunderstandings and an incomplete conception, then to
let the connected reproduction of Spinoza's thoughts itself, by its
unity, prove itself the right one, and at the same time to reject the
false conceptions in an indirect way. If one succeeds in rendering a
system in its wholeness from its living center, then this organic
wholeness itself will prove the correctness of the conception in the
whole and in the particular, and will of itself remove the false or
incomplete conception.

With this preface, whose length must find its excuse in the fact that
much has had to seek its place in it which in a connected history of
philosophy would have found a place in the book itself, I now hand this
work over to the not numerous public that books of this content can
count on finding in Denmark. The remaining monographs will be published
as quickly as circumstances allow me. Next after this one will follow
the two treatises on \emph{Francis Bacon of Verulam} and the
\emph{Système de la nature}, which, on account of their relation to each
other, will appear simultaneously.

\bigskip
\begin{flushleft}
\hspace{2em}{\small\emph{Copenhagen}, 6 April 1856.}
\end{flushleft}

\begin{flushright}
{\small\textbf{H.\ Brøchner.}}\hspace{4em}
\end{flushright}

\bigskip
\begin{center}\rule{2.5cm}{0.4pt}\end{center}

\newpage

% --- p. 15 ---
\chapter*{Benedict Spinoza}
\addcontentsline{toc}{chapter}{Benedict Spinoza}
\markboth{Benedict Spinoza}{Benedict Spinoza}
\opage{15}

\begin{center}
{\Large\textbf{\emph{Benedict Spinoza.}}}
\end{center}

\vspace{2\baselineskip}

\begingroup
\setstretch{1.15}
\begin{quote}
\small
Quicquid est, \emph{in Deo} est; et nihil sine Deo esse neque concipi
potest. --- Eth.\ I.\ prop.\ 15.

\medskip

Æternum illud et infinitum ens, quod Deum seu naturam appellamus, eadem,
qua existit, \emph{necessitate} agit. --- Eth.\ IV.\ præf.

\medskip

Summum mentis bonum est \emph{Dei cognitio}, et summa mentis virtus Deum
cognoscere. --- Eth.\ IV.\ prop.\ 28.
\end{quote}
\endgroup

% --- p. 16 ---
\opage{16}

\newpage

\chapter*{Introduction}
\addcontentsline{toc}{chapter}{Introduction}
\markboth{Introduction}{Introduction}

\begin{center}
\textbf{The Historical Genesis of Spinozism.}
\end{center}

% --- p. 17 ---
\opage{17}If, in order to understand a philosophical system in its historical
significance in the widest sense, in order to grasp the moving forces
which, often only dimly sensed by the system's originator, work through
it, in order to see its justification and the limit of that
justification, it is necessary to cast a glance over the whole course of
development of human thought; then, in order to grasp a system in its
determinate form, in order to follow its fundamental thought into its
consequences, and in order to discover its historical and psychological
genesis, it is sufficient to go back to the nearest great division in
history, where the new beginning forms a gulf between itself and the
past and comes forward with the justified claim of what is peculiar to
independence. We also find that the systems which follow immediately
after such epochs see in them the limits of their historical
connection, and, often with a misjudging disdain for the past, take
them only as the starting point for a consistent development of thought.
Thus Spinoza's historical parallel, Fichte, in his first appearance goes
back only to Kant, and Spinoza himself only to Descartes.

With Descartes (cf. the Preface, p. 8) the one line of development of
modern philosophy began, with the demand to develop the content of
philosophy out of self-consciousness, from the thinking I, according to
a fixed method.
% --- p. 18 ---
\opage{18}\setcounter{footnote}{0}Indications of this direction had already
been given in earlier thinkers. The \emph{French skeptics} had pointed
back to self-consciousness, \emph{Campanella} had again recalled
Augustine's inference from the I's thinking to its being; but these
indications had stood as isolated; they had not been pursued into their
consequences. They come to be so only when Descartes, casting off
tradition and turning polemically not so much against the already
retreating scholastic philosophy as against the revived thought of
antiquity, leads back through doubt about everything to the I as the
secure point of attachment, and, as doubt becomes a separating-out of
everything that does not belong to the essence of the I, determines that
essence as thought, and now, in an immediate intuition (simplici mentis
intuitu), not through an inference --- for cogito ergo sum is no
inference --- lets the thinking I confirm itself as being (sum
cogitans). In the starting point there is thus given thought's immediate
certainty of itself, its immediate self-confirmation, its active
separation from and opposition to that whose essence is not determined
as thought, and which, consistently, ought to have been posited as a
non-being. But in this separation and opposition thought loses its
content\footnote{This consequence holds for Descartes, not for Spinoza.
In the former the I is sundered from nature, and nature is excluded from
God's essence; in Spinoza the I as thinking is, in God, in unity with the
essence of nature.}, and while the demand is made, to develop knowledge
out of itself, no way is left open by which it might, from itself alone,
arrive there again. The concept of \emph{God} must here pave the way.
Descartes arrives at it in that doubt, which, more deeply understood,
was thought's holding fast to itself in the separating-out of everything
that was alien to it, is understood as the sign of finitude, and the
doubting, finite I now finds in itself the idea of the infinite being,
an idea that guarantees the existence of this being, since it can have
its origin only from such a being, not from the finite I, as otherwise
the effect
% --- p. 19 ---
\opage{19}\setcounter{footnote}{0}would contain a greater reality than the
cause. In the representation of God as perfect is now found the
representation of his veracity; it gives us confidence in the truth of
our clear and distinct ideas\footnote{Consistently understood, the clear
and distinct idea itself is the criterion of truth, since in it thought
is in unity with itself and confirms itself. Cogito ergo sum is the
first form of the thinking I's self-confirmation, the idea of God as the
absolute essence of thought its highest form, and, in the concept, its
first. But when doubt is understood as finitude, the clear and distinct
idea becomes uncertain, and must be supported by the representation of
God.}; for the infinitely perfect being cannot will to deceive; that
would presuppose an imperfection. But among the clear and distinct ideas
Descartes counts, by giving the concept of thought an extent that stands
in opposition to his deeper conception of it (in that not only
intelligere and velle, but also \emph{imaginari} and \emph{sentire} are
counted under it), the idea of a body that is united with the mind as
well, and through this body the connection is opened with the world of
bodies, with nature, which, despite its opposition to the essence of
thought, constantly exercises a magical power over Descartes's thought
and becomes its essential and dearest object, while he as it were shuns
the more concrete domains of spirit, especially the ethical.

In the opposition to mind, the independence of nature from mind is
given; as the mind is conceived as substance, the body must also be
conceived as such; the thinking substance gets over against itself an
\emph{extended} substance; for in the case of nature, conceived \emph{through}
thought, abstractly as its opposite, quantity, extension, becomes the
essential determination. --- Separated from these finite substances,
\emph{the infinite substance, God,} comes to take its place. Since
\emph{thought} is the starting point, God is determined in such a way
that thought finds itself again in him; extension is excluded from his
essence. But since the \emph{I} is the starting point, and fixes itself
from pure activity into a thinking \emph{being}, it does not vanish in
God, although his essence
% --- p. 20 ---
\opage{20}\setcounter{footnote}{0}becomes the absolute essence of thought;
the consequences of this concept are not drawn, they are merely hinted
at\footnote{In that the thought of the infinite is given a priority
over the thought of the finite (cf. Ep. I, 119), and it is posited as
originally being in us (pr. phil. I, 20-22 and elsewhere).}, and the
concept of substance is held in an ambiguity that makes it possible to
preserve a doubtful independence for the finite I's. Descartes does
indeed determine substance as: une chose, qui existe en telle façon,
qu'elle n'a besoin que de soi-même pour exister (principes de la philos.
I, 51)\footnote{I cite from: \emph{Oeuvres philosophiques de Descartes}
publiées d'après les textes originaux par L. \emph{Aimé-Martin}. Paris
1839.} and by this definition it is restricted to God; but it is at once
extended again to those of the created things that do not need
\emph{other created} things in order to exist, but need for that only
God's concours ordinaire, while the rest of the finite is called the
\emph{qualities} or \emph{attributes} of these substances. From the
``\emph{uncreated substance}, which thinks and is independent'' (I, 54),
are then distinguished the two classes of \emph{created substances}, the
\emph{corporeal} and the \emph{mental}, each of which has its own
peculiar attribute, which constitutes their nature and essence, and on
which their other attributes depend. Thus \emph{extension} in length,
breadth and height constitutes the essence of the extended substance,
\emph{thought} that of the thinking substance (I, 53). The two classes
of finite substances have nothing in common, but the consciousness of
their union in man, who is conceived as a \emph{composition} of both,
forces Descartes to the hypothesis of a point of contact in the glandula
pinealis (conarion), the immediate seat of the soul in the brain,
through which the motions of the body communicate themselves as
sensations to the soul, and the soul's thoughts as motions to the
body\footnote{This hypothesis Descartes does not, however, hold to
consistently, since already in him there are hints of the theory of
Causæ occasionales developed by his school, by which divine causality
becomes what is at work in both mind and body, mediating between the
changes of both.}; but the possibility of this
% --- p. 21 ---
\opage{21}transformation in the communication he does not demonstrate, and
it cannot be demonstrated from his principles.

Instead of ending with a unity of knowledge that is coherent and grounded
in itself, Descartes therefore ends with a double, unreconciled
opposition: the opposition between the thinking and the extended
substance, and between the finite substances and the infinite. For the
thinkers who attached themselves to him, the two problems decisive for
his whole thought were thus set: to find the unity between mind and
body, and the unity between God and the finite, and thereby to reconcile
the double discord.

The consequences of Descartes's conception of the essence of the two
finite substances were drawn by \emph{de la Forge}, \emph{Geulincx} and
\emph{Malebranche}, in that the difference of essence was held fast,
and the union was sought, not in man, through a fantastic hypothesis,
but through a postulate in God, who in both series of substances
produces corresponding changes, while the finite substances are only
occasions (causæ occasionales) for each other's changes. And
\emph{Geulincx}, \emph{Malebranche} and \emph{Clauberg} alike carried
the consequences of the relation set up by Descartes between the
infinite and the finite thinking substances so far that only the
presuppositions of religious representation prevented the fusion of
the finite thinking spirit with the infinite. It was also only these
presuppositions that kept \emph{Malebranche} from positing, with
decisiveness, \emph{real} extension as an attribute of God.

Only by abandoning those representations could the strict consequence
of the Cartesian concept of substance be drawn, and the second problem
be carried further; but as the consequence was drawn with regard to
this problem, the changed conception acts directly on the conception of
the first problem, and a new formula is set up for its solution. This
development takes place in \emph{Spinoza}, and since he has thought to
the end what his predecessor had left standing ambiguous, he becomes the
real representative of the direction indicated by the latter.

% --- p. 22 ---
\opage{22}\setcounter{footnote}{0}For Descartes \emph{thought in the form of
the I} had been the center; from the I he had taken his starting point,
and even over against the infinite, which was conceived under the same
fundamental concept, he had held fast to it, and had rejected the
problems that threatened its existence\footnote{Thus especially the
problem, sharpened by Descartes's conception of conservation, of the
relation of divine omnipotence to the self-activity of the finite
substances, in particular to human freedom.}. But since the I was
determined as \emph{thinking}, to the exclusion of the conditions of
individuality, its essence was determined only by the
\emph{universality} of thought, and the consequence had to be drawn that
lets the independent subsistence of the finite form vanish, and posits
the essence, thought in its infinity, as the one being. The thought of
the infinite does not tolerate limitation by a finite subsisting by
itself, and all the less since it is not a secondary but the original
thought; for the limited --- Descartes had already stated this --- is
grasped only through the unlimited; the latter is the positive, that
which lies at the ground; limitation is a determination within it and
presupposes it\footnote{How Spinozism's peculiar determination of the
essence of the infinite substance leads to determining the general
proposition that the finite cannot exist \emph{alongside} the infinite
in such a way that what makes it finite is only something negative
(that limitation is negation), so that the finite is not merely in the
infinite, but is in it \emph{only as modification}, will be shown later
(cf. chapter 3.).}. The \emph{first thought} of the thinking substance is
the thought of infinity, the thought of God; but in this thought it is
not separated from the thinking infinite, it is in unity with it: an
independent position for the finite thinking substance becomes an
impossibility. What Descartes called finite thinking substances must
therefore change their name and become \emph{modifications of the one
infinite substance}; they must have their reality only in it. And the
same holds of the \emph{extended} substances. They are posited by
Descartes
% --- p. 23 ---
\opage{23}as created by God, existing, like the thinking ones, solely by
his continued conservation, but the essence of extension is not posited
in God; he is determined only as thinking, although in Descartes this
determination has significance only as opposition to extension. But the
cause can be the cause only of what is grounded in its nature; it can
impart only what it itself has in itself; if, therefore, God is posited
as cause of the extended substances, extension as infinite must be
posited in God, and just as the thinking substances had their reality
only as modifications in God, the same must hold of the extended ones.
Only thereby is the contradiction in Descartes also removed, that the
thinking and the extended substance are grasped solely through the
attribute that constitutes their essence, without there lying in their
concept a relation to God, of whom they thus become independent with
regard to essence, while with regard to existence they are dependent on
him. This contradiction between dependence and independence is removed
only when that which Descartes conceived as independent substances is
conceived as attributes of the one infinite substance, which is
undivided in both, so that both express its essence wholly, and while
they are independent of each other, in that each by itself expresses the
\emph{whole} essence of substance, they are at the same time in unity by
expressing the \emph{one} substance (cf. chapter 2.)

If now, as the presentation of Spinoza's philosophy will show, that
philosophy is wholly concentrated in his concept of substance, and this
concept develops immediately as a consequence of what was given in
Descartes, then it follows that it is not necessary to pursue the origin
of Spinozism further back.

Some have wished, deceived by a superficial similarity and misunderstanding
the peculiarity of Spinozism, to trace it back to the doctrine of
emanation of the \emph{Kabbalah}. The development in the following
chapters will show how incompatible Spinoza's conception of eternal
being, of the finite as being only in the infinite and, according to its
truth, with its limitation cancelled, is with the emanation doctrine's
conception of finite existence as having arisen through an unfree
unfolding in
% --- p. 24 ---
\opage{24}\setcounter{footnote}{0}\emph{time}, as independent and as fallen
away. The scarcely appreciative manner in which Spinoza expresses
himself about the Jewish tradition (cf. tract. theol.-pol. Cap. 7 and
9), and the few and indefinite references to the views of Jewish
thinkers (Ep. 21, 29. Eth. II, 7, schol.) are not suited either to prove
a closer relation of dependence.

But while such a relation of dependence with regard to what is
essential, to the fundamental thought of the system, is rejected, there
is not thereby rejected either an influence of other philosophers on
particulars of the execution, or a general influence of certain
tendencies of the age and of Spinoza's Jewish nationality. On Spinoza's
doctrine of the state \emph{Hobbes} had an effect, perhaps also
\emph{Hugo Grotius}, although Spinoza does not name him; for the
determination of the relation of religion to philosophy \emph{Herbert of
Cherbury} may have had significance, for the sharper conception of the
problems prepared by Descartes \emph{Arnold Geulincx}\footnote{Especially
with regard to the conception of individual minds as modi in God, and
of extension as res singularis and infinite.}. On Spinoza's pantheistic
fundamental view the pantheistic elements that had remained in the
spirit of the age from the period of transition may have had influence;
on the strict conception of substance, God, as the one being, the Jewish
representation of God. But every external influence, whether through
definite theories or through the whole spiritual atmosphere of the age,
is in Spinoza worked up into inner consistency; his system forms a whole
that is understood from itself, from one principle.

If, ahead of the actual presentation of Spinoza's philosophy, we wish
for preliminary orientation to indicate this principle and its immediate
consequences for the form of his thought and for his method, then the
elements for this lie in what was developed above concerning his
relation to Descartes. As the concept of substance is conceived in such
a way that in its unity it fills the whole of existence, it governs
thought as its not merely essential but, strictly speaking, its only
% --- p. 25 ---
\opage{25}\setcounter{footnote}{0}object. In substance thought finds all
reality; in it, it finds itself. But the fact that all reality is in
substance grounds the \emph{pantheistic} conception. The fact that
thought is immediately in unity with substance, which is its own
infinite, objectified essence, and at the same time the essence of
extension, gives thought its immediate certainty and the system the
character of \emph{dogmatism}. The fact that all \emph{reality} is
enclosed in the concept of substance entails that this concept must
unfold itself for thought, and the character of this unfolding, which is
in the element of eternity, not of time, must be determined by the form
of natural law, necessity. For since thought has first separated nature
out from itself as its opposite, and determined it by the abstract
concept of extension, this conception acts back upon thought itself,
which must seek its content in nature; and as thought and extension are
united in the unity of substance, mechanical necessity, the law of
abstractly conceived nature, becomes the law of thought and of substance.
And even though in substance, as the one existent, necessity is not an
external necessity, not compulsion, but an inner necessity, determined
by its own essence, a libera necessitas, the concept of \emph{causality}
nevertheless displaces that of purpose, and this both on account of the
influence of the concept of nature on the concept of substance, and on
account of the strict holding fast to the concept of infinity in
opposition to anthropomorphic representations.

The \emph{method} of the system is determined immediately by the
fundamental concept. Thought is not separated from its object; the object
is the reality in thought, thought's own reality; the unfolding of the
object must therefore give the form of thinking its stamp; the object's
necessity, conditioned by the concept of causality, determines the
\emph{mathematical} form of the \emph{development} of
thought\footnote{Not exactly the synthetic method, which Spinoza has
used only in the Ethics and in the Principia phil. Cartes., and whose use
instead of the analytic he ascribes especially to a didactic purpose,
but in general the method that proves with mathematical stringency,
which, with mathematical necessity, links to what is clear and simple in
itself the firmly chained series of consequences.},
% --- p. 26 ---
\opage{26}while the starting point and end point of thought becomes
\emph{immediate intuition}, corresponding to substance's being in itself.

A \emph{division} of philosophy from its principle Spinoza has not
given, but only hinted at. It will be brought out in the following
chapter what peculiar difficulties had to arise for him on this point,
and these difficulties will become still clearer in the closer
development of his doctrine.

\bigskip
\begin{center}\rule{2.5cm}{0.4pt}\end{center}

\newpage

\chapter*{Chapter One.}
\addcontentsline{toc}{chapter}{Chapter One. Spinoza's Life and Writings}
\markboth{Spinoza's Life and Writings}{Spinoza's Life and Writings}

\begin{center}
\textbf{Spinoza's Life and Writings.}
\end{center}

% --- p. 27 ---
\opage{27}\setcounter{footnote}{0}The life history of a thinker is the history
of his thought, for the content of his life is his thought. His outward
life does not acquire an immediate interest, as does that of an
outwardly acting personality, whose life stamps itself on his own time
as it itself bears that time's stamp. It becomes either wholly
indifferent, or of significance only insofar as its agreement with his
thinking offers a guarantee of that thinking's power to permeate a human
life. The interest becomes a derived one; it cannot become essential,
for in the appropriation of the content of thought the personal relation
to the thinker vanishes, and must vanish, if the appropriation is to
become complete and free. One can therefore, in the case of a thinker,
even if one demands more than the genesis of his thought, be content
with a light outline of his life, provided only that this outline is
definite enough for the stamp of his thought to become visible through
it; one can renounce all those details, significant even in their
apparent contingency, that make the outline into a finished picture. As
regards Spinoza, such contentment is necessary; he has found only one
somewhat more detailed biographer, the Lutheran pastor
\emph{Colerus}\footnote{La vie de Spinoza, tirée des écrits de ce fameux
philosophe et du témoignage de plusieurs personnes dignes de foi, qui
l'ont connu particulièrement. The Hague 1706. 8. An earlier Dutch
edition appeared in Utrecht.},
% --- p. 28 ---
\opage{28}and this simple-minded, honest narrator, who despite his
theological antipathy is appreciative of the philosopher he did not
understand, has found in Spinoza's simple, outwardly uniform life only
scant material to impart, and his reports are completed only a little by
the details that are found in the preface to Spinoza's posthumous
writings, in \emph{Bayle} (in his Dictionnaire historique et critique),
in ``La vie de Spinoza par un de ses disciples'' (probably the physician
\emph{Lucas}), Amsterdam 1719, and in \emph{Nicéron}, Kortholt and
Leibnitz.

From the character of Spinoza's thinking, as it already shows itself in
the way in which he conceived his predecessor's principle and carried it
further, one might expect that his life would not be like that of many
of the thinkers of the period of transition, especially the Italians,
whom thought seized as a restless, never satisfied, consuming passion,
and, filling them as with prophetic frenzy, tossed them unsteadily about
the world as its heralds, and let them, proud of their calling, collide
defiantly with the established order, to end as martyrs of the idea. In
Spinoza we must expect that the calm of the fundamental view will cast
its light over his life, that the vanishing of the finite in the
infinite will stamp itself in resignation, in a giving up of all the
more concrete relations of life in order to live only for knowledge, for
the unity of thought with the divine in intellectual love.

This we also find confirmed when we read the few and meager reports
that have been handed down to us about Spinoza. They are insufficient to
form for us a whole and full picture of a human life, but sufficient to
let us follow a break with everything that is not the infinite, with
nationality, with the faith of the fathers, with love for a woman, with
outward esteem, in order to live solely in and for the highest. ---

\emph{Baruch Spinoza} was born in Amsterdam on 24 November 1632. His
father was a Jewish merchant who on account of religious persecution had
left his fatherland, Portugal. From the few traits of him to be found in
the
% --- p. 29 ---
\opage{29}son's biographers, one gets the impression of a clear intellect
and of freedom from prejudice in religious views. As the family's
circumstances were only modest and the son early revealed excellent
gifts, he was destined for the learned estate and entrusted to the
instruction of rabbis, among whom \emph{Moses Morteira} has a name in
modern Hebrew literature. This instruction continued after the early
death of Spinoza's parents, and to it was later joined, in accordance
with the precept of the Talmud, training in a craft, the making of
optical glasses, in which he gradually attained a great perfection. It
was not long before his clear intellect, sharpened by Talmudic subtlety
but driven by a thirst for truth that did not let it find rest in the
artificial edifice of the rabbis, embarrassed his teachers and at once
roused both the pride of his co-religionists and their secret unease.
Already before he was 20 years old, faith in the authority of the rabbis
had been destroyed in him, and his thought, which found no satisfaction
in their doctrine, sought the truth in itself. Unfortunately it is
precisely for the years that lie between the first trace of an awakening
independence in Spinoza and his decisive break with his people and his
religion that the biographers are poor and in disagreement. For internal
reasons the author of ``la vie de Spinoza'' seems on this point to
deserve credence in preference to Colerus, who imagines the course of
Spinoza's education after the pattern of the Christian learned schools,
and therefore brings in at an earlier point humanistic elements of
development with which Spinoza only later came into contact. According
to that biographer's account it seems as if Spinoza let his thoughts
ripen in silence; but as they took on a firmer form, and his doubt
extended from Jewish learning to the Jewish religion, the mistrust of his
teachers and fellow believers was aroused, and it was strengthened by
his withdrawing more and more from the synagogue and seeking the company
of Christians. His teachers, who had seen in him a future pillar of the
synagogue, now became hostile to him, and sought by peremptory commands
to bring him
% --- p. 30 ---
\opage{30}back; the congregation sought to buy him back with an annual
stipend of 1000 guilders, and, when that did not help, to get rid of him
by assassination. This too failed, and now they resorted to the last
means, the excommunication of the apostate; and soon afterwards the
rabbis, supported by the Reformed clergy, succeeded in prevailing upon
the magistrate of Amsterdam to banish him from the city for a time. From
this time on Spinoza stands as a stranger to his people and its
religion. His Jewish name Baruch he exchanges for its Latin form
Benedictus, and philosophy becomes the occupation of his life.

Already at the time when he began to feel a stranger among his people,
he had sought, outside the narrow circle of Jewish learning, to procure
the means for a richer scientific development and for humanistic
culture. From the freethinker, the physician \emph{Frants van den Ende},
he learned the classical languages, especially Latin, and with its help
he gained access to the study of the natural sciences and of philosophy,
in which he soon chose Descartes as his guide. When he had been
excommunicated, van den Ende took him into his house for a time, and a
romantic light falls over this period of his life through his love for
his teacher's gifted daughter. She preferred one of his fellow pupils,
and his love now sought its object in the divine.

Spinoza's excommunication from the synagogue falls about the year 1660,
and in the same year, probably, his removal from Amsterdam. He spent a
short time with a friend in the country near the city; the following
year he settled in the village of \emph{Rynsburg} near Leyden, where in
1663 he published his first work: Principia philosophiæ Cartesianæ, more
geometrico demonstrata, to which Cogitata metaphysica form an appendix.
In 1664 he moved to \emph{Voorburg} near The Hague, and from there,
after a five-year stay, at his friends' urging, to \emph{The Hague}
itself, where he lived the rest of his life.

% --- p. 31 ---
\opage{31}On this part of Spinoza's life Colerus is more detailed, at least
detailed enough for his narrative to give us an impression of its
strict, earnest, solitary, energetic character. We see Spinoza working
strenuously and without interruption toward his goal: the knowledge of
God, of nature, and of the unity in which we are with nature, and toward
the calm elevation above the human passions that is acquired through
the clarity of knowledge. Once, Colerus relates, he had not left for 3
months the simple room he occupied at the house of the painter van der
Spyck in The Hague, so deeply did he immerse himself in his studies. By
his extraordinary frugality he made it possible for himself, although he
had to work for his living, to devote most of his time to them. His
biographer tells us circumstantially and with naive details how meagerly
he lived, how little he spent, although in his outward appearance he
always maintained a certain modest elegance, how few and simple were the
effects found after his death, and how disinterested he was in his
dealings with his sisters and with his richer friends and patrons.

Besides the many whom curiosity drew to him, there were many who were
affected by Spinoza's personality. His deep knowledge impressed them,
while his lively and witty conversation, his humanity, his openness and
fidelity, and the nobility of his character (cf. Oldenburg's remarks in
Ep. 1) attracted them; but only few of them understood him, and his
relation to them became more a humane relation of goodwill, and theirs
to him one of admiration. From his letters, which are all replies, one
sees that only one single one of his correspondents, the anonymous writer
who wrote the letters 61, 63, 65, 67, 69 and 71, designated by various
marks, was capable of penetrating more deeply into his philosophy, and
for Spinoza, whose life was his thinking, only a unity in thinking could
ground a real relation of friendship. How little he was understood by
those who drew near to him and felt attracted by him one can best infer
from his correspondence with \emph{Oldenburg} (cf. p. 40), which, with a
single interruption, continued from 1661 until Spinoza's death. Even from
the
% --- p. 32 ---
\opage{32}\setcounter{footnote}{0}last of these letters it appears that
Oldenburg, who after all spent his life in the company of philosophers
and men of science, had still not managed to grasp a single one of
Spinoza's thoughts, and the same must naturally have been the case with
all those who wished to follow Spinoza without becoming aware that they
would have to give up the whole foundation on which their circle of
representations rested.

Even before Spinoza's writings had drawn more general attention to him,
we see from his letters that he had been noticed by scholars and drawn
into the correspondence that was so important and so eagerly used an aid
to the scientific life of that time\footnote{Spinoza's correspondence
seems, according to remarks in L. Meyer's preface to his Opera posthuma,
to have been fairly extensive, and the collection of letters found in his
writings (74) has probably made up only a small part of it. In the Great
Royal Library here in this city there is an unprinted note from Spinoza
to \emph{Grævius}, but its content is without significance.}. His first
work at once won him a name, but also provoked attacks from the strict
Cartesians\footnote{On Spinoza's relation to these cf. Ep. 19.}. His
fame rose still higher through the tractatus theologico-politicus,
published in 1670\footnote{Not only scholars became attentive to him,
but also political personages: thus, besides Jan de Witt, the Prince of
Condé, who invited Spinoza to a meeting with him in Utrecht in 1673, and
Fr. of Luxembourg.}, but in a still greater proportion than his fame
grew the number of his opponents, especially among the clergy, and
their fanatical attacks were the reason why Spinoza's most important
writings were published only after his death. Only by this reserve
could he procure for himself the peace on which he set so high a value,
and which was so necessary to him, and for its sake he also declined the
honorable offer to take up a professorship of philosophy in Heidelberg,
which was made to him in 1673 by the liberal-minded Elector \emph{Carl
Ludvig} through his councillor, who was evidently only
% --- p. 33 ---
\opage{33}little pleased with this commission, the Heidelberg professor
\emph{Ludvig Fabritius}.

\emph{In the political events} of his fatherland Spinoza took no active
part, but he followed them as an observer, with the interest that his
view of the political relation was bound to give them. His conception of
the indispensability of freedom for the attainment of the end of the
state (cf. chapter 7) made him an adherent of the republican form of
state and attached him to those who defended it against monarchical
tendencies, especially to \emph{Jan de Witt}, who became his protector
and benefactor, and whose lamentable death filled him with bitter grief.

His relation to \emph{religion} was given in his conception of its
essence. For himself his philosophy was his religion: the knowledge of
God is love for him and unity with him. Over against others, for whom
the common religious representations had kept their significance, he
sought to hold fast to the practical and humane element in religion, and
to conceive it in such a way that it could become an approximation to
true knowledge and to true agreement among men. Passing over the dogmas,
he brought forward what could serve as an incitement to love, concord and
honesty, and emphasized the significance of obedience and the value of
common worship. Precisely this formed the most frequent subject of those
conversations with his housemates in which he sought rest from more
strenuous occupation, and in which he knew how to mix earnestness with
kindness, and often with cheerful humor and good-natured irony.

Spinoza's health had been weak from his youth; strenuous mental work and
night watches undermined it still further, and a consumption developed
which on 21 February 1677 put an end to his life in his 45th year. His
Christian biographer relates circumstantially how patiently and
steadfastly he had borne his sufferings, how calm and gentle his death
was, and he rejects the untrue rumors that were in circulation
concerning the manner of his death.
% --- p. 34 ---
\opage{34}\setcounter{footnote}{0}Another theologian (\emph{Kortholt})
wonders that an atheist could have so beautiful an end, and those who
could not comprehend this, and did not have such good information as
Colerus, spread rumors that Spinoza had died of fright at coming to share
the fate of Jan de Witt, or by suicide.

A portrait of Spinoza that is still preserved shows us a face in which
the dark complexion, the curved nose, the strongly marked mouth and chin
point to his Oriental nationality, while the high, powerfully formed
forehead, the firm and fine line of the eyebrows, and the lively, clear
eye betray the keen thinker. The expression of the face bears the stamp
of the earnestness of resignation, but also of the calm in the infinite
that is acquired through it.

Of Spinoza's \emph{writings}, the two that appeared in his lifetime have
already been named above. The first of these, \emph{Renati Des Cartes
Principiorum philosophiæ pars I \& II}, more geometrico demonstratæ par
\emph{Benedictum de Spinoza}, Amstelodamensem, --- to which
\emph{Cogitata metaphysica}, in quibus difficiliores, quæ tam in parte
metaphysices generali quam speciali occurrunt quæstiones breviter
explicantur, are attached as an appendix, originally came into being
through an external occasion, part of it having been worked out for the
use of a pupil whom Spinoza was instructing in the Cartesian philosophy,
but did not consider ripe for initiation into his own views, already
independently developed at that time (cf. Ep. 9). As Spinoza, at his
friends' request, worked it up for publication, there are indicated, in
the course of the presentation of Descartes's doctrine, the points where
difficulties that are insoluble from Descartes's standpoint make
necessary consequences that lead toward Spinoza's own\footnote{It is
especially in the problem of \emph{God's relation to extension} (cf. pr.
ph. C. I, 9. schol. 16. 21. II, 2. schol., Cog. met. I, c. 2.), in the
determination of the concept \emph{creatio} (cf. pr. ph. C. I, 12. Cog.
met. II, c. 10), and of the concept \emph{finite things}, existing
\emph{extra Deum} and \emph{in se}, in the question of the compatibility
of this last concept with God's properties, such as omniscience,
omnipresence, immutability (cf. pr. ph. C. I, 11 dem. 12, coroll. 2.
Cog. met. I, c. 2. II, c. 1. 3. 7. 10 f.), and finally in the problem of
\emph{human freedom} (cf. Cog. met. I, c. 3. II, c. 11 f.), that
difficulties show themselves, for whose solution from Spinoza's own
standpoint hints are given (thus, e.g., with regard to the concept of
creation, which is cancelled in that it is restricted to apply to
existentia, not essentia, and in that its application to modi is
denied). On his relation to Descartes Spinoza has also expressed himself
in Ep. 2 (from 1661).}, and in the whole design of the presentation the
way is paved
% --- p. 35 ---
\opage{35}\setcounter{footnote}{0}for his own conception. This character of
the work indicates the purpose of its publication.

While this work was thus meant to prepare the reception of Spinoza's
philosophy, the next was to make room for it by marking out the boundary
between theology and philosophy and thereby securing a neutral territory
for philosophy, and by demonstrating the compatibility of free thought and
free speech with the end of the state and with its peace. This is the
immediate task of the \emph{Tractatus theologico-politicus}, continens
dissertationes aliquot, quibus ostenditur, libertatem philosophandi non
tantum salva pietate et reipublicæ pace posse concedi, sed eandem nisi
cum pace reipublicæ ipsaque pietate tolli non posse, which appeared in
1670. The essential content of this treatise will be developed in
chapter 7, where Spinoza's conception of religion and the state is
presented. It appeared without Spinoza's name and was quickly suppressed
at the instigation of the clergy, but a few years later (1673) was
published again under various fictitious titles\footnote{Such as: % print: Som; (semicolon for colon)
\emph{Danielis Heinsii} operum historicorum Collectio prima. Editio
secunda priori editione multo emendatior et auctior. Accesserunt quædam
hactenus inedita. Lugd. Batav. 1673. 8. --- and: \emph{Francisci
Henriquez de Villacorta}, doctoris medicinæ, a cubiculo regali Philippi
IV et Caroli II archiatri, opera chirurgica omnia, sub auspiciis
potentissimi Hispanorum regis Caroli II. Amstelod. 1673. 8.}. Later it
was Spinoza's intention
% --- p. 36 ---
\opage{36}\setcounter{footnote}{0}to publish it anew, accompanied by
notes\footnote{Cf. Ep. 75 (added in Bruder's edition). The treatise is
mentioned besides in Ep. 8. 14 (``tractatum de scriptura''), 17. 19 f. 47
and 49.}, but he was prevented from doing so by his early death. The
notes he had wished to add are included in the French translation by
\emph{St. Glain}, which appeared in 1678 in Leyden under the title: La
clef du sanctuaire par un savant homme de notre siècle. Besides this
French translation, for which two other titles are also found: Traité
des cérémonies superstitieuses des Juifs tant anciens que modernes, and:
Reflexions curieuses d'un esprit désinteressé sur les matières les plus
importantes au salut tant public que particulier; --- there is a Dutch
one by \emph{Joh. Henrik Glasemaker} from 1693, and several more recent
French and German ones.

Only after Spinoza's death\footnote{Shortly before his death Spinoza
destroyed a treatise, de Iride, and a Dutch translation of the Old
Testament with notes that he had begun, of which the Pentateuch was
finished. Lost are the greater part of his letters, a treatise on the
devil, and a defense, written in Spanish, of his leaving the synagogue.}
were those writings published that are the real source of our knowledge
of his philosophy. His friend \emph{Ludvig Meyer} published this
collection in the year of his death (1677) under the title: B. d. S.
Opera posthuma (it was by Spinoza's wish that his name was not
mentioned). It contains, besides Spinoza's chief work, \emph{Ethica},
ordine geometrico demonstrata\footnote{It is divided into 5 parts: de
Deo; de natura et origine mentis; de origine et natura affectuum; de
servitute humana seu de affectuum viribus; de potentia intellectus seu de
libertate humana.}, the two unfinished treatises, \emph{tractatus
politicus} and \emph{tractatus de emendatione intellectus}, as well as
74\footnote{The 75th letter, which is added in \emph{Bruder's} edition,
was found a few years ago in Holland and published by \emph{H. W.
Tydemann}. It is to Lambert van Velthuysen.} letters from and to him,
and, as an appendix, an unfinished \emph{Compendium grammatices linguæ
hebrææ}. The collection was translated into Dutch in the same year by
\emph{Jarrig Jellis} under the title: \emph{De nagelate Schriften van B.
d. S.}

% --- p. 37 ---
\opage{37}\setcounter{footnote}{0}Of the \emph{Ethics} Spinoza had already
made a draft early on. In the oldest of his letters to \emph{Oldenburg},
from 1661, the first definitions are given, though in a form somewhat
different from the present redaction of the Ethics (cf. chapter 2),
together with the main propositions of the first part of the Ethics. This
is of importance for Spinoza's development, since it furnishes a proof
that, already when he published Principia phil. Cartes., he was clear
about his fundamental concept, which diverged from Descartes's\footnote{For
Spinoza's relation to Descartes it is also of significance that he began
the study of Descartes's writings only after long independent thought.}.
The essential defects he finds in his predecessor's philosophy he states
in the same letter (Ep. 2). In a letter from \emph{Simon de Vries},
written in February 1663 (Ep. 26), definite propositions of the first
book of the Ethics are cited. In a letter to \emph{Blyenbergh} from 1665
(Ep. 36, cf. Ep. 38) Spinoza refers, with regard to the concept of
appetite, to the development in his ``Ethica necdum edita'', where it is
found in the third book. In the correspondence with L. \emph{Meyer} (Ep.
63---71) from the last years of Spinoza's life (1675---76), propositions
from the various parts of the Ethics are frequently cited, which had then
probably been brought into its final form. Finally, \emph{Oldenburg} in a
letter from 1675 (Ep. 18) urges Spinoza to publish ``tractatum illum ---
--- quinque-partitum''. From these citations, and from many developments
in the letters that agree almost word for word with the Ethics, it is
thus seen that this chief work of Spinoza had long been finished in
draft, in which form it had been communicated to several of those who
stood closest to him\footnote{This is also said expressly in the
editor's preface: ``ea (i.e. the Ethics) namque ante multos annos a
diversis fuit descripta et iis communicata.''}; but from the variants in
the citations, which, where they are of importance for the
determinations of concepts, will be touched on later, one sees at the
same time that he constantly occupied himself with it anew in order to
find the most exact expression for his thought.

% --- p. 38 ---
\opage{38}The treatise \emph{de intellectus emendatione} is, according to the
editor's statement, one of Spinoza's first works, with which he
constantly occupied himself, but without ever bringing it to completion.
It is mentioned already in 1663 in a letter from Oldenburg (Ep. 8) as:
opusculum, in quo de rerum primordio earumque dependentia a prima causa,
ut et de intellectus nostri emendatione tractas. Later it is mentioned in
several letters (Ep. 42. 63 f.). The essential part of its content (cf.
chapter 5) is taken up into the 2nd part of the Ethics, especially in the
scholia to propositions 17 and 40. Having determined the knowledge of the
eternal and of our unity with it as the highest and true good, in
opposition to the untrue, perishable goods, and having set it as the
task to attain this knowledge, the treatise develops the various forms
of knowledge, of which the knowledge of reason proves to be the only one
that leads to the goal. Its instruments are the true ideas, which
contain certainty in themselves, and the true method is the one that
shows how the mind links the ideas together according to the norm of
its highest idea. In the \emph{first} part of the treatise the true ideas
are now distinguished from ideæ fictæ, falsæ and dubiæ, which belong to
the imagination, and the difference between the imagination and
intellectus, the source of the true ideas, is determined. In the
\emph{second} part the rules are given for how the unknown is derived
with certainty from the known, in such a linkage that our mind becomes
the subjective rendering of the objective essence of nature as a whole
and in its parts. Here the treatise is now led, through the distinction
between what is conceived through its proximate cause and what is
conceived through its essence alone (sola essentia), to treat the laws
for the definitions of ``created and uncreated'' things, since knowledge
advances only through the definitions, which express the peculiar
essence of things. As the transition was now to be made to the statement
of the means by which res particulares æternæ, the essential object of
reason, can be derived from the highest concept, the treatise breaks off
at the decisive point, after the powers of the intellect and its
properties have still been treated only briefly.
% --- p. 39 ---
\opage{39}According to the plan of the treatise (p. 371) it still remained,
first, to develop the order in the linkage of knowledge, and then the
idea of the most perfect being. It will be shown in the development of
the fundamental concepts of Spinozism (chapter 2) that precisely at the
point where the treatise was broken off an insuperable difficulty arises
for Spinoza, so that it is not accidental that it had to remain
unfinished.

The \emph{Tractatus politicus}, which is likewise only a fragment, was
written shortly before Spinoza's death. In the first 5 chapters he
develops with extraordinary sharpness and clarity his conception of
natural right, of the concept of the state, of the right of the sovereign
power, of the extent of its authority, of the relation of states to one
another, and of the end of the civil community. In the 6th and 7th
chapters he treats monarchy, in the 8th---10th aristocracy, both that in
which one city is the head of the whole state and that in which several
cities form a federal state. In the 11th chapter the development of
democracy is begun, but here the treatise is broken off after a few
general determinations have been given. Besides the complete presentation
of this form of state, Spinoza would also (cf. the preface) have treated
``the laws and other special political questions.''

The collection of \emph{letters} found in the Opera posthuma probably
makes up only a small part of the correspondence Spinoza carried on.
Besides one letter (Ep. 48), written by a stranger to one of Spinoza's
acquaintances concerning the tractatus theologico-politicus, it contains
32 letters to and 41 from Spinoza. Most of the letters were originally
written in Latin, 24 (namely 30---41, 43---47, 50 and 55---60) in Dutch,
and were translated into Latin by L. Meyer. Besides a good deal of
information about Spinoza's personal circumstances and about his
writings, the letters contain numerous developments, partly of points
that are touched on in his works but here as a rule come out in a
modified presentation, called forth by questions and misunderstandings;
partly of such points as are either not at all or only in passing
% --- p. 40 ---
\opage{40}mentioned elsewhere (thus of various Christian dogmas, of
physical, mathematical and optical problems, of logical determinations
and the like). A great part of the letters, and precisely of those most
important in philosophical respect, are to and from unnamed authors.
Under several of the different marks one and the same author, the editor
of his posthumous writings, seems to be hidden. For not only Ep. 29, but
doubtless also Ep. 62, 64, 66, 68, 70 and 72 are addressed to him,
although the letters to which these are replies are dated from various
foreign cities, and in part purport to have been written by someone
other than the one who sent them to Spinoza. Among the named
correspondents, next after \emph{Leibnitz}, with whom Spinoza exchanged a
few words on the occasion of a writing sent to him, \emph{Heinrich
Oldenburg} is the best known, and the one to whom the greatest number of
letters is addressed. Under Cromwell he was envoy of the Lower Saxon
Circle in London and later held the same post under Charles II, who was
well disposed toward him and, when the English Society of Sciences
became royal in 1662, made him its secretary. He was an eager
intermediary between several of the philosophers and scholars of that
time, and, despite his slight capacity for grasping a deeper thinking,
not without an eye for the peculiarity of talents and recognition of it,
and, though not free from a motive of vanity in his approaches to
superior personalities, yet animated by a real interest in the
development of the sciences. The other named correspondents are:
Spinoza's friend \emph{Simon de Vries}, the Dordrecht merchant
\emph{Willem van Blyenbergh}, the Heidelberg professor \emph{Joh. Lud.
Fabritius}, \emph{Albert Burgh}, who made an unsuccessful attempt to
convert Spinoza to Catholicism, and \emph{Peter Balling}. By their
initials alone are designated, among others, \emph{Ludvig Meyer} (Ep.
29) and the Spaniard who had gone over to Judaism, \emph{Isaac Orobio de
Castro}, as well as perhaps \emph{Joh. Bredenburg} (who in 1675 wrote
against the tractatus theol.-pol.) and the Mennonite \emph{Jarrig
Jellis}, the translator of Spinoza's Opera posthuma.

% --- p. 41 ---
\opage{41}Spinoza's \emph{Compendium grammatices linguæ hebrææ}, which forms
an appendix to his posthumous writings, is of no significance for his
philosophy; at most, particulars of the terminology (thus the use of the
designation \emph{attribute}) and the conception of the relation of the
parts of speech to one another, in that they are all, with the exception
of interjections and conjunctions and certain particles, brought under
the concept Nomen, can have some interest.

Besides the writings left unfinished at Spinoza's death, he had also,
according to Meyer's statement, laid the plan for an ``\emph{integra
philosophia}'', to which there are frequent references in the tractatus
de emendatione intellectus (in the preface and in many of the notes) and
perhaps in Ep. 74; and for an ``Algebra breviori et magis intelligibili
methodo''.

In what relation this \emph{integra philosophia} would have stood to the
Ethics, and into which particular sciences it would have divided itself,
Spinoza has not indicated. In the introduction to the treatise de
intellectus emendatione he enumerates, but, as he expressly adds, without
systematic order, the various sciences that all aim at the same goal,
the highest human perfection, i.e. the knowledge, common to all men, of
the unio, quam mens cum tota natura habet. To this, he says, there belongs
first the knowledge of nature, and next a society that can lead to that
knowledge as easily and securely as possible. Besides, attention must be
given to moral philosophy and to the doctrine of the education of
children, and, since health is no slight means to attaining the intended
goal, a complete medicine must be set up; and since much that is
difficult is made easy by art, and art can save us much time and afford
us much advantage in life, mechanics must by no means be neglected. But
above all a way must be found to heal the intellect, and, so far as this
can be done before science (initio), to purify it, so that it grasps
things with ease, perfectly and without error. --- This enumeration
indicates in essentials the same sciences as Descartes's
% --- p. 42 ---
\opage{42}\setcounter{footnote}{0}division (cf. Descartes's letter to the
French translator of the principia philosophiæ); but whereas with
Descartes they form a single series, in which morals becomes the
concluding but unattainable science, with Spinoza the series is divided,
in accordance with the parallelism between the two attributes that man
knows. If we wish to bring the enumerated sciences into a systematic
order, the doctrine of knowledge becomes an introduction, but as such
incomplete, since it can become clear only after the fundamental
thoughts of the system have been developed. The fundamental science
becomes the knowledge of nature, i.e. of substance, God. This becomes
the metaphysical part of the system, in which the general concepts both
of ethics and of physics are developed. To it there then attaches, on
the one side, ethics\footnote{Ep. 38: magnam Ethices partem, quæ, ut
cuivis notum, \emph{metaphysica et physica} fundari debet.} with
politics and pedagogy, on the other side mechanics and medicine. But on
Spinoza's fundamental conception there can be no question, in the strict
sense, of a real separation of the sciences, and the parallelism he set
up between thought and extension, the difficulty of carrying out which
would have shown itself precisely in the more special physical domain,
was bound to prevent both series from being carried out, since this
would have led either to too great a fusion or to too great a
separation. It is therefore surely not accidental that Spinoza treated
only the one series, and he can hardly have intended to give more of the
other than ``generalia in physicis'' (cf. Ep. 63). At a later point we
shall see that the opinion expressed by the editor of his posthumous
writings, that Spinoza in his integra philosophia would have developed
``veram motus natu\emph{ram atque qua ratione a priori tot varietates
in materia deducendæ essent}'', could not, by the nature of the case,
have been fulfilled.

In Spinoza's chief work, which has come down to us under the title
\emph{the Ethics}, the essential content of the enumerated sciences is
found. The doctrine of \emph{knowledge}, which can only
% --- p. 43 ---
\opage{43}be developed completely after the metaphysical foundation, which
is the object of an immediate certainty, has gone before, is given in
the second part of the Ethics, summarily in prop. 40, schol. 2. The
general part of \emph{metaphysics} is found in the first book of the
Ethics; the general \emph{physical} propositions, which are
presuppositions for the more special ethical part, are given in the
second part of the Ethics, prop. 13. \emph{Politics} is given in its
foundation in the doctrine of the affects, in which the principles of a
\emph{pedagogy} also lie, and its fundamental concepts are determined in
Eth. IV, 37. schol. 2. For \emph{mechanics} and \emph{medicine} there
are, by the nature of the case, only few elements to be found, but in
the development of the parallelism between the mind's power to think and
the body's power to act (cf. especially Eth. III, 28. dem. V, præfat.,
Ep. 45) lies the fundamental principle for the conception of them. ---
Knowledge of Spinoza's system thus rests essentially on the Ethics; the
presentation must proceed from it, and the other writings serve only to
complete it on particular points.

Of \emph{editions} of Spinoza's writings there are, besides the
original, separate ones, 3 collected editions: by \emph{Paulus} (Jena
1802---3. 2 vols. 8.), \emph{Gfroerer} (Stuttgart 1830. 8. In this,
however, the Hebrew grammar is missing) and \emph{Bruder} (Leipzig
1843---46. 3 vols. 12. Stereotype edition.) The Ethics, tractatus
politicus and de emend. intellectus are included in: Renati des Cartes
et Spinozæ dissertationes philosophicæ ed. C. \emph{Riedel} (Leipzig
1843. 2 Vol. 12.) Of more recent translations there are two complete
ones: a German one by B. \emph{Auerbach} (Stuttgart 1841. 5 vols.) and a
French one by E. \emph{Saisset} (Paris 1843. 2 Vol.).

Where page numbers are given with the citations in this treatise, they
refer to the original editions when another edition is not expressly
named.

\bigskip
\begin{center}\rule{2.5cm}{0.4pt}\end{center}

\newpage

\chapter*{Chapter Two.}
\addcontentsline{toc}{chapter}{Chapter Two. The Fundamental Concepts of Spinozism}
\markboth{The Fundamental Concepts of Spinozism}{The Fundamental Concepts of Spinozism}

\begin{center}
\textbf{The Fundamental Concepts of Spinozism.}
\end{center}

% --- p. 44 ---
\opage{44} What was developed in the Introduction concerning Spinoza's relation to
Descartes leads immediately to the point from which his system can be
surveyed. When the concept of substance is taken in its strictness, it becomes
the all-embracing concept, hence also the starting point, and a
starting point that can no longer be abandoned. Whereas Descartes had to proceed
from the finite I, Spinoza proceeds from the absolute
essence of this I, from the infinite, and can begin only with this because \emph{only} the
infinite is. It is therefore not through a deduction that this is reached;
it is the object of an immediate certainty; in it thought thinks only itself
in its infinity. It is not in conflict with this that proofs of the existence,
infinity and unity of substance are found in Spinoza, and
that not merely in principia phil. Cartes., where in essentials the
Cartesian demonstration is reproduced and Spinoza's distinctive conception
is only hinted at, but also in the Ethics (I, 7. 8. 14). For this demonstration
has only the significance of something formal, of a development of concepts for
the subject, not of something that grounds. This will become clearer in its place
when the essence of human thought in its relation to
substance has been developed.

I. The stated conception of \emph{substance} now underlies
its definition: Per \emph{substantiam} intelligo id, quod in se est,
et per se concipitur; hoc est id, cujus conceptus non indiget conceptu
alterius rei, a quo formari debeat. (Eth. I, def. 3, cf. Ep. 27.).
Substance ``is conceived through itself'' (\emph{concipitur per se}) in that
its concept is not formed, i.e.\ derived, from a higher concept (as, e.g.,
that of motion from that of extension), but is itself the highest concept. It
is ``in itself'' (\emph{in se}), not in another (as modus is in
substance).
But in the determination that it is in itself and is conceived through itself,
% --- p. 45 ---
there \opage{45}\setcounter{footnote}{0}lies immediately that it also is \emph{per se}, that it is the
\emph{cause of itself}\footnote{De emend. intell. pag. 386
puts side by side: esse in se, and: causa sui.}; for the concept of an effect
includes the concept of the cause (``effectus cognitio a cognitione causæ
dependet, et eandem involvit.'' Eth. I, ax. 4); what therefore includes only its own concept cannot be an effect except of itself. % print: closing « missing
But since it is its own cause, its existence is not contingent, but
\emph{necessary}, posited with its essence: Per \emph{causam sui} intelligo id,
cuius essentia involvit existentiam, sive id, cuius natura non potest
concipi nisi existens. (Eth. I, def. 1.). This unity of existence and
essence is \emph{eternity}: per \emph{æternitatem} intelligo ipsam existentiam,
quatenus ex sola rei æternæ definitione necessario sequi concipitur
(Eth. I, def. 8.)\footnote{Æternitas = infinitas et necessitas
existentiæ; Eth. I, 23. dem.}. And since eternity is the absolute
affirmation of the existence of a nature, there follows from it immediately
\emph{infinity}, while everything finite is ex parte negatio (Eth. I,
8. schol. 1.). But in the concept of infinity lies \emph{unity}, both
as exclusion of \emph{plurality} (for this presupposes not merely a
mutual limitation, but also a separation between essence and
existence, and moreover, since it is not given by the definition of the thing, an
\emph{external} cause, cf. Eth. I, 8 schol. 2), --- and as exclusion of
\emph{divisibility}\footnote{In these determinations of the concept of substance
we find again the traces of the genesis through which it has arisen.
When substance is posited as unity of essence and existence, it is posited as
the objective essence of thought distinguished from the sensible; for
sensibility separates essence and existence, limits the intellect, and
the intellect finds only itself again where that separation is sublated, where it
affirms itself in the unity of thought and existence. In the determination of
substance as causality we find again the consequence of the fact that thought,
distinguished from the sensible, in order to arrive at a content, must seek
this in nature; for it is from nature that the determination of causality is taken, and
transferred to thought. But it will appear that it is again deprived of its
distinctive conditions and distinctive significance, so that nature
asserts itself only as thought, i.e.\ as sublated nature.}.

% --- p. 46 ---
\opage{46}\setcounter{footnote}{0}Substance is thus determined as the eternal, infinite, one, which is through itself and is
conceived through itself. All reality is concentrated in
its concept; in it thought has the infinity of its essence and the unity
of knowledge.

Predicated of substance, the concepts of unity and of infinity denote the
same; for as a merely \emph{numerical} determination unity has no application
to substance\footnote{Cf. Ep. 50. 39 ff., princ. ph. Cart. I, 11.
schol., Cog. met. I, 6. Eth. I, 8.}. But infinity is not conceived
as the negation of a positive finite, but as the positive essence of the
negative finite\footnote{On the relation between the infinite and the
indeterminate cf. Eth. I, 8. schol. 1. II, def. 5. Ep. 29. 41. 64. Cog.
met. I, 6. De int. em. p. 360.}. As negations both
determinations would lead to an abstract, undifferentiated emptiness and lifelessness,
but thought seeks in them precisely the expression of fullness. And in the determination:
cause of itself, it intuits activity, and with activity the absence of
difference is removed. When unity is determined by infinity and
activity, we are led to the two concepts which in the Spinozist
system, together with substance, make up the fundamental concepts: the concept
\emph{attribute}, and the concept \emph{mode}.

Infinity excludes limitation, for limitation negates. But
limitation takes place only through what is of the same kind, not through the
heterogeneous: Ea res dicitur \emph{in suo genere finita}, quæ alia
eiusdem naturæ terminari potest. Ex. gr. corpus dicitur finitum, quia
aliud semper maius concipimus. At corpus non terminatur cogitatione,
nec cogitatio corpore (Eth. I, def. 2.). Infinity can therefore include within
itself, without being cancelled as infinity, as it were a new
infinity, a multiplicity which, on account of the difference of kind,
entails no limitation, and such a multiplicity becomes, as
compatible with infinity, compatible with unity. It is even demanded
by infinity, since everything that
% --- p. 47 ---
\opage{47}\setcounter{footnote}{0}expresses essence and includes no negation belongs to it (Eth. I, def. 6.
explic.).

II. With this multiplicity, which without cancelling infinity is posited
within the unity of substance, the concept of the \emph{attribute} is given: per
\emph{attributum} intelligo id, quod intellectus de substantia
percipit tanquam eiusdem essentiam constituens (Eth. I, def. 4). This
definition must not be misunderstood as if the words: id, quod intellectus
--- --- percipit, were to make the attribute a \emph{merely subjective}
conception\footnote{To elucidate the relation between substance and
attribute, the use of the determination \emph{essentia} in Princip. phil.
Cart. and Cog. metaph. may serve. It is set on the one hand as synonymous
with \emph{definitio} (pr. ph. C. II, 15. schol.) and with
\emph{conceptus}, on the other hand as synonymous with \emph{natura}
(cf. pr. ph. C. I, 4. ax. 11 not. and def. 9, with II, 15. schol.), with
\emph{vis}, qua in meo esse me conservo (pr. ph. C. I, 7. schol., Cog.
met. II, 1), with \emph{potentia} (pr. ph. C. II, 2. schol.) and with
\emph{perfectio} (Cog. met. I, c. 6), and there is talk of \emph{essentia
formalis}, i.e.\ objective (Cog. met. I, 2). The subjective and the objective
conception are thus, in accordance with the whole character of Spinozism, in unity.
In accordance with this character substance too had to be determined both as
that which \emph{is in itself} and as that which \emph{is conceived through itself}.}. The misunderstanding
has, however, a great authority to lean on; I must therefore
refute it more fully. Already the closing words of the definition: tanquam
eiusdem essentiam constituens, contradict the subjective conception,
still more the whole character of Spinozism (cf. chapter 5, the doctrine of
knowledge) and definite statements of Spinoza's. If one first compares the
two places in the letters where definitions of the attribute are given,
one will find that in the one place (Ep. 2, written 1661, cf. Ep. 4)
it is determined in essentials in the same way as the Ethics determines
substance: notandum, me per \emph{attributum} intelligere omne id, quod
concipimus per se et in se, adeo ut ipsius conceptus non involvat
conceptum alterius rei; --- likewise, that in this letter the definition of
the attribute is attached directly to the definition of God without the
% --- p. 48 ---
\opage{48}\setcounter{footnote}{0}concept of substance coming in between\footnote{One also sees from Ep. 2 ff. that the first 4
propositions of the Ethics were originally set out as axioms, and that
the series of propositions then began with the present prop. 5. Just as in
the use of the concept of substance a formal deviation has taken place from
the Ethics in its present form, so too the determination accidens is
used instead of the determination modus, whereas Spinoza later distinguished
them.}. In the other place (Ep. 27, from 1663) \emph{substance} is first defined
as that, quod in se est et per se concipitur, hoc est,
cujus conceptus non involvit conceptum alterius rei. Then there is
added: \emph{Idem} per \emph{attributum} intelligo, nisi quod attributum
dicatur respectu intellectus, substantiæ certam talem naturam
tribuentis. Many places in the Ethics point to the same thing. Thus
already the definition of \emph{God} as the substance \emph{consisting} of infinitely many
attributes (\emph{constantem} infinitis
attributis) states the real unity. Likewise the expressions in Ethics I, 15.
schol.: hinc conclusimus, \emph{substantiam} extensam unum ex infinitis
Dei \emph{attributis} esse. The proofs of propositions 10, 19 and 20 in the first book
of the Ethics rest on the same presupposition. In the proof of prop. 10
the peculiarity of substance, that it can be conceived through itself, is transferred
to the attribute as \emph{constituting the essence of the former} (tanquam eius
essentiam constituens). The \emph{eternity} of substance is transferred in prop.
19. demonstr. to the attribute as that, quod divinæ substantiæ \emph{essentiam}
exprimit, hoc est, id quod ad substantiam
pertinet\footnote{cf. the reference to prop. 11 in the proof of prop. 21
in the first book of the Ethics.}. The same expressions are used in the proof of prop. 20,
and therefore in Coroll. 2 to this proposition it can be concluded: Deum \emph{sive}
omnia Dei attributa esse immutabilia\footnote{With this one may
compare princ. ph. C. I. ax. 4, where it is said that substance has
more reality than an accident or a modus, but not that it
has more reality than the attribute.}. It leads to the same result when
Eth. II, 40. schol. speaks of essentia
% --- p. 49 ---
\opage{49}\setcounter{footnote}{0}for\emph{malis}\footnote{Formalis means in Spinoza and the philosophers of that time
the same as in our philosophical terminology: objective. What we call
subjective was then called objective, e.g.\ esse objectivum, i.e.\ the
essence of things as objectified in representation. In some of the scholastics,
e.g.\ Occam, subjective (i.e.\ existing as subject,
substance) conversely denotes what we call objective.} attributorum. Likewise the
way in which Spinoza conceives our knowledge of the attributes (cf. Ep.
28. Eth. II, 40 schol. 2) points to their objective reality. And finally
the proofs of Eth. I, 4 and 30\footnote{cf. Cog. met. I, c. 2.
pr. ph. C. I, 7 schol.} state expressly that it is not merely in the
intellect (intellectus) that the attribute is equal to substance.

The attribute\footnote{The designation attribute is used more loosely not only
in many places in Spinoza's earlier writings and in the Hebrew
Grammar (e.g.\ C. 33), but also in the Ethics; thus when res,
aliquid and the like are called attributes (Eth. II, 40. Schol. 1), and when it is
said (Eth. IV. 37. schol. 2): apparet, iustum et iniustum, ---
notiones esse extrinsecas, non autem \emph{attributa}, quæ mentis naturam
explicent.} is thus of one essence with substance; it is conceived,
like substance, through itself, and is eternal and infinite. It is infinite
despite the plurality of the attributes, since limitation has significance only
within the same genus. It is \emph{in suo genere} infinitum, since we
can deny infinitely many attributes of it; it is not
\emph{absolute infinitum}, since to the essence of the absolutely infinite
belongs everything that expresses essence and includes no negation (Eth.
I, def. 6 explic.). This \emph{ens absolute infinitum}\footnote{ens
absolute \emph{indeterminatum}. Ep. 41.} is \emph{God}, i.e.\ substantia constans
infinitis attributis, quorum unumquodque æternam et infinitam
essentiam exprimit (def. 6). In the concept: \emph{God}, the concept of substance has
first reached its complete expression. Substance is the infinite;
only in the unconditionally infinite, therefore, does it reach its concept. The
infinite substance, God, must, as containing infinite reality,
consist of \emph{infinitely many} attributes (Eth. I, 9), each of which
is conceived by itself, but which do not therefore
% --- p. 50 ---
\opage{50}\setcounter{footnote}{0}constitute several beings or several different substances: Id enim
est de natura substantiæ, ut unumquodque eius attributorum per se
concipiatur, quandoquidem omnia, quæ habet, attributa, simul in ipsa
semper fuerunt, nec unum ab alio produci potuit; sed unumquodque
realitatem sive esse substantiæ exprimit. Longe ergo abest, ut
absurdum sit, uni substantiæ plura attributa tribuere; quin nihil in
natura clarius, quam quod unumquodque \emph{ens} sub aliquo attributo debeat
concipi, et quo plus realitatis aut esse habeat, eo plura attributa,
quæ et necessitatem sive æternitatem et infinitatem exprimunt, habeat,
et consequenter nihil etiam clarius, quam quod ens absolute infinitum
necessario definiendum (ut def. 6. tradidimus) ens, quod constat
infinitis attributis, quorum unumquodque æternam et infinitam certam
essentiam exprimit (Eth. I, 10, schol.).

At this point, however, an essential difficulty arises. The
different attributes are conceived through themselves, are in themselves, without
dependence on one another, without mutual limitation. In the single
attribute there thus lies nothing that points to the others; it is
closed in itself, contains a reality infinite in its kind, and this: in
its kind, denotes no comparison. If one proceeds from the
attribute\footnote{From Spinoza's standpoint one must no doubt say that it is
unwarranted for thought to separate the attribute from substance, since they
must immediately be seen in unity, and that thought, when it conceives
substance as separated from the attributes, conceives it under a new,
self-contradictory attribute, that of indeterminacy. But however
much one sees the matter from his standpoint, the result remains the same, and
the difficulty only takes another form. In the doctrine of knowledge (cf.
chapter 5) we shall see that idea corporis and idea ideæ express one thing and
are both in the mind. With one idea the two attributes are thus thought in
their difference. But the objective essence of extension as
different from thought is not thought; it expresses its actuosa
essentia through an operation parallel with thought, which remains
independent and is only ideally taken up into the mind. Here, therefore,
dualism returns.}, % print: Aanden, Her (comma for full stop)
one does not therefore arrive at a unity of the
% --- p. 51 ---
\opage{51} several attributes in substance; one comes in any case to
conceive the attribute immediately as substance, but then of course as
substance with \emph{one} attribute.
Nor, by taking one's
starting point from the thought of the absolute substance, will one of necessity
arrive at the \emph{plurality} of the attributes. The absolutely infinite, absolutely
perfect being must, as containing infinite reality or being,
have infinite attributes, Spinoza concludes. But in the premises there lies,
if one disregards Spinoza's presupposition, the absolute
difference of the attributes, only the exclusion of negation, and since negation, strictly
taken, becomes equal to limitation, and limitation has significance only
within the single attribute, the inference would lead only to the attribute
in its infinity without any limitation, not to infinitely many
attributes. If one imagines another historical presupposition for
Spinozism, the inference could have led to positing thought as the
sole attribute, so that the others became forms of thought,
or extension as the sole one, so that thought became a form of it.
The independence of the attributes in their multiplicity is an assumption that
rests on the historical presupposition of Spinozism, on the dualism between
thought and extension, whose unity was demanded. The two attributes, thought
and extension (cf. chapter 5), confirm themselves immediately to the
human mind, since this is corporis idea, and although already in
the fact that the body can become an ideatum there lies a sublation of the
abstract opposition, they are posited, by virtue of that dualism, as absolutely
different, so that each is conceived by itself, and they are in
unity only in so far as they express substance, hence in so far as their
difference is abstracted from, not through themselves as these determinate ones, such that their
difference is preserved. But since a \emph{plurality} of
attributes has been posited, the infinity of substance must express
% --- p. 52 ---
\opage{52}\setcounter{footnote}{0}itself in the infinity of plurality, the \emph{infinity of number}, and the attributes are posited
as \emph{infinitely many}\footnote{The relation between the infinite number of the
attributes and the duality of attributes that is expressed in the
human being, together with the difficulty that arises at this
point, will be touched on in chapter 5.}.

The determination of the multiplicity of the attributes is thus conditioned by
Spinoza's historical presuppositions, but the task had nonetheless to become
that of bringing this determination into harmony with the unity of substance. In the
rebuttal of objections given in the passage cited above,
the presupposition of the several attributes is taken as starting point, and from it
their independence is inferred; but the possibility of their unity despite
the difference is not demonstrated. Nor does this happen through the theory of the
parallelism between the series of modifications that belong to the
different attributes (Eth. II, 7); for this difference remains
standing as a presupposition, which is as little led back to a
unity as it has issued from one. It is not sufficient
(with Feuerbach) to say that precisely in that spirit and matter are conceived
each by itself and through itself, in that both are independent in their concept,
they express not themselves, but substance, God, and, since both are
substance, they express only this, that they are substance, so that what is
affirmative in them is only substance \emph{purely as such}, which is
indifferent to whether it is spirit or matter\footnote{In Spinoza's
conception the \emph{peculiarity} of the attribute affirms itself precisely immediately
to thought.}, and for which they are therefore both, \emph{in so far as} they are
something \emph{determinate}, something \emph{different} from each other, and
in so far as one looks only at this determinateness \emph{as} determinateness,
not at the fact that it is and expresses substance, only attributes,
properties, moments, that constitute the idea of God. --- For in this way
the peculiarity of the attribute is conceived as separable from
substance, hence as something vanishing; substance, which becomes
indifferent to the difference, either becomes an abstraction, a notio
universalis,
% --- p. 53 ---
\opage{53} really a subjective imagination (cf. chapter 5), or the difference
becomes something subjective and substance indeterminate. If the
difference is conceived in this inessentiality, one cannot even say that
the determinacies as moments constitute the idea of God; for they become merely
something subjective, the idea something real.

A closer grounding and development is, however, not to be found in Spinoza.
Being led, in the endeavour to give the concept of substance
fullness, by the presuppositions of his time to seek this in the
infinite multiplicity of the attributes, he has only sought a determination that made
it possible to conceive the single ones as infinite, and this he has
found in the peculiarly conceived concept of limitation, and in the
distinction between what is infinite in its kind and the unconditionally
infinite. He has next given the proof that, \emph{on the
presupposition of} the infinitely many attributes in the one substance,
these had to be conceived in and through themselves, mutually independent. But
the taking up of plurality into unity, the compatibility of the independent differences
in substance, whose essence is expressed \emph{only} in the attributes,
he has not demonstrated, and, as will appear later, could not
demonstrate. This is an incompleteness in his system, which the historical
exposition has to point out and to illuminate, not to remove.

III. Substance, which in its unity includes the multiplicity of the
attributes, now receives its further determination through the
\emph{concept of causality}. Substance is causa sui; this is
not merely an expression of the negative fact that it is not effected by
another, but has \emph{positive} significance, since it determines God's essence
as activity. This shows itself in that essentia is conceived as one with
potentia (Eth. I, 34: potentia Dei, qua ipse et omnia sunt et agunt,
est ipsa ipsius essentia) and as containing vis agendi (Eth. IV, 26
dem.); it appears likewise from the expression: actuosa essentia (Eth. II,
3 schol. cf. Cog. metaph. II, c. 11), from the conception of the attributes,
thought and extension, as active, from the important Proposition 16 in
the first book of the Ethics
% --- p. 54 ---
\opage{54}\setcounter{footnote}{0}(``ex necessitate divinæ naturæ infinita infinitis modis sequi
debent'')\footnote{Actiones, quæ ex ipsius [Dei] naturæ necessitate
consequuntur. Eth. IV. App. Cap. 4.}, from the proof of prop. 36 in the same
book, from the concept idea ideæ (cf. chapter 5), and from the doctrine of the affects,
which rests throughout on the view of activity.

When the determination \emph{causa sui} receives positive significance, there is posited
in the identity of what effects and what is effected a constantly
vanishing, or rather \emph{eternally vanished}, difference, and
through this difference as difference we arrive at the third
Spinozist fundamental concept, the concept \emph{modus}\footnote{As
synonyms are used: \emph{modi, modificationes, affectiones}. From \emph{modus}
Spinoza distinguishes \emph{accidens} (Cog. met. I, c. 1), which in Descartes is
synonymous with it and is still used thus by Spinoza himself in Ep. 4.
The designation \emph{affectio} is used in Cog. met. I, c. 3 as a synonym of
\emph{attribute}.}. ``Per \emph{modum} intelligo substantiæ affectiones, sive id
quod in alio est, per quod etiam concipitur''\footnote{Modificationes
are: Id, quod in alio est, et quarum conceptus a conceptu rei, in qua
sunt, formatur. Eth. I, 8. schol. 2.}. (Eth. I, def. 5). Since modus is
determined as being \emph{in another} and \emph{conceived through another}, it is
determined as \emph{posited by another}, by that in which it is and through
which it is conceived (Eth. I, 16. Coroll. 1, 25). Determined as effect,
modus is different from substance, its cause, and when it is now
determined by a plural (substantiæ \emph{affectiones}), and is brought
together with the determination: \emph{single things} (\emph{res particulares}
nihil sunt nisi Dei attributorum affectiones sive modi, quibus Dei
attributa certo et determinato modo exprimuntur. Eth. I, 25. Coroll.),
then this difference, since what is posited by substance and is in
substance is determined by plurality and limitedness, seems to become an
absolute one, and to acquire retroactive force to overturn the conception of the
concept of substance.

The determination: singularity and limitation, leads us to the problem of
\emph{the relation of the infinite to the finite}, and the question might
at first seem to have to
% --- p. 55 ---
\opage{55} be put thus: how does the infinite in Spinoza produce a
finitude, how does the \emph{transition} take place from the former to
the latter? But in this form the problem has already been given a skewed
position, since, forgetting the system's fundamental thought,
reality is granted to the finite. If the fundamental concept, the system's
starting point, is held fast, then the task must on the contrary become precisely this:
to conceive the
relation of the infinite to the finite in such a way that it is seen that there
can be no question of a transition from the former to the latter, \emph{that the
finite is not as finite}, that nothing is without substance,
and that the very knowledge of the nothingness of the finite does not come to be ---
hence is not subject to the form of finitude; --- but \emph{is
eternal}\footnote{With regard to the problem of \emph{finite things} cf. Eth.
II, 8. Cor. 9. Cor. Ep. 58. Pr. phil. Cartes. I, 5. 7. Cog. metaph. I,
c. 2. II, c. 4. 7).}.

It might seem that the determination of substance as
\emph{causality}, of God as causa omnium rerum (Eth. I, 16. Cor. 1),
would have to settle the problem of the being of the finite in the opposite way.
\emph{Real} causality includes a temporal relation, and with time
finitude is given, if not independently subsisting finite things.
But Spinoza evades this consequence partly through the concept causa
\emph{sui}, in which causality is as it were bent back upon itself; partly
by laying down that God is the cause of all things \emph{eo sensu} (in the same
sense) in which he is called causa sui (Eth. I, 25. schol.), for which reason
he must be determined not as \emph{causa transiens}, i.e.\ as acting over
upon another, but as \emph{causa immanens}\footnote{``\emph{agentem} sive
causam immanentem.'' Gr. hebr. c. XII. --- In the Index to the Opp. posth.
the determinations causa immanens and \emph{natura
naturans} are put side by side (see below).}, as indwelling cause (Eth. I, 18, cf. Ep.
21); partly, finally, indirectly, by developing the difference that,
on the ordinary conception of the causal relation, arises
between cause and effect, with respect both to essence and to existence
(Eth. I, 17. schol.)\footnote{That the argumentation in this place is
polemical against the ordinary conception of the concept of causality is seen
from other places, e.g.\ ``God has also formaliter (objectively)
something in common with finite things.'' Ep. 4. cf. Eth. I, 25. Coroll.
(Prop. 10 in the 2nd Book of the Ethics, which concerns the \emph{single} modus, has
no application here.)} ---
% --- p. 56 ---
\opage{56}\setcounter{footnote}{0}which, if finite things were conceived as independently subsisting,
would lead to the sublation of the concept of substance, and, applied to the
concept: causa sui, to absurdities.

If we look more closely at the significance of the concept of causality in Spinoza, we
also see that the peculiarity of its conception cuts off that
consequence. \emph{Causality coincides for Spinoza with
logical consequence}\footnote{Spinoza can therefore also
designate logical consequences as actiones. In this conception too it is confirmed that
the Spinozist substance is the objectified essence of the intellect.
Therefore in the doctrine of knowledge, too, agere is posited as equal to
Intelligere (Eth. IV, 24. dem.).}: ``Everything proceeds with necessity
or \emph{follows} always with the same necessity from God's infinite
nature, \emph{in the same way} in which it follows from eternity to eternity
from the nature of the triangle that its three angles are equal to two
right angles'' (Eth. I, 17 schol.); --- and in the 16th
Proposition of the 1st Book cited above the proof is drawn from the fact that from the
definition of every thing (i.e.\ from the essence of the thing) several
properties follow with necessity; since, therefore, from the absolutely
infinite attributes of the divine nature, each of which expresses an essentiality
infinite in its kind, infinitely much must follow in infinite modifications.

Through this logical conception of causality the temporal determination is
removed\footnote{God's priority is a priority of causality; Pr. phil.
Cart. I, 12. Coroll 4.}; the relation comes to lie within the element of eternity,
which is not determined by time, measure or number, the determinations of finitude and
of imagination (Ep. 29).

While the concept of causality, on Spinoza's conception, thus did not
lead to finitude, finitude is likewise excluded when in
causality we press the concept of substance as the infinite causa sui.
Causality is in unity with the essence of substance; through it
this essence realizes itself. But
% --- p. 57 ---
\opage{57}\setcounter{footnote}{0}since the effect of the causal substance is substance itself, the
infinite causality is exhausted in the infinite effect; there remains no
room for an effecting of something finite, and in the essence of substance there is
found no point of attachment for it. What follows from the essence of substance
follows with necessity; but necessity includes
infinity\footnote{cf. Eth. I, 7, demonstr. with I, 8. schol. 1.}.
And since substance is determined as the infinite, absolutely real
being, which is unity of essence and existence, there remains as little
room for an actual existence separated from it as for an essence
separated from it. Even if the existence of the modification, in Spinoza's
usage, is determined not by eternity but by
duratio\footnote{\emph{Duratio} est \emph{indefinita} existendi continuatio
(Eth. II. def. 5). Eth. II, 45. schol.: durationem, hoc est
existentiam, quatenus abstracte concipitur, et tanquam quædam
quantitatis species. With regard to the relation between æternitas and
duratio cf. Eth. I, def. 8. explic., prop. 23. dem., Cog. met.
I, c. 4. II, c. 1. 2; with regard to the relation between duratio and
essentia Eth. I, 24. dem., IV. præf.}, which is the expression of the
existence that follows not from the essence of the thing but from the nature of the cause,
yet duratio, on account of the infinity of the cause, is to be thought
as infinite when the effect is seen in its truth, i.e.\ in its unity with
substance, and only when it is separated from this by an abstraction
do the determinations of finitude come forth (Ep. 29).

Spinoza's conception of these determinations of finitude leads to the same result:
limitation, plurality, divisibility, time, number
and measure. \emph{Limitation} is conceived only as negation (Ep. 50); it
therefore becomes unthinkable when one proceeds from substance
(the attribute); for since limitation can take place only through what is of the same essence,
it would have to take place through the attribute itself, and this, as what limits,
would have to be conceived as the negation of what is limited,
hence as its own negation
(Ep. 50). Limitation is therefore excluded by the concept of substance; and
when it is seen from the standpoint of a hypothetically
% --- p. 58 ---
\opage{58} assumed finite, it becomes, as something merely negative, something
vanishing.

The determination \emph{plurality} does not lie in the definition of the thing, but
presupposes an external cause (cf. Eth. I, 14 schol. Ep. 41). If one follows
this proposition back, then, as the series of causes is led
back to the first cause, the unity of substance will exclude plurality,
hence finitude. --- The determinations \emph{divisibility, time, number, measure}, which
arise when duration and quantity are separated by an abstraction from
substance and from ``the eternal order'' (cf. Ep. 29), lead to
contradictions that cannot be resolved: an infinite series that can form no unity,
an infinity of parts that cannot manage to make up
a whole, a \emph{greater and lesser} infinity. But these
contradictions are removed when the representation of the finite, single things
is removed; only the reality that is ascribed to these as finite calls
those forth.

But while the consideration of all these concepts thus leads away
from an existing finite, it might seem that Spinoza himself has
made an attempt to derive the finite from the infinite in
Propositions 21--23 of the first book of the Ethics. They read: ``Everything that
follows from the absolute nature of one of God's attributes must always have
existed and existed as infinite, or is eternal and infinite through
this attribute.'' --- ``What follows from one of God's attributes,
in so far as it is modified by such a modification as
exists with necessity and as infinite through the attribute, must also
exist with necessity and as infinite.'' --- ``Every modus that
exists with necessity and as infinite must with necessity have
followed either from the absolute nature of one of God's attributes, or from
an attribute modified by a modification that exists both with
necessity and as infinite.'' --- A kind of gradation is given here:
the \emph{attribute}; that which is \emph{immediately produced by God} (immediate
a Deo productum), and \emph{that which is produced mediante infinita
quadam modificatione} (cf. Ep. 65 and Spinoza's answer in
% --- p. 59 ---
\opage{59} Ep. 66, where as examples of the second kind are adduced, under
the attribute of thought: \emph{intellectus absolute infinitus}, under
that of extension: \emph{motion and rest}; as an example of the third
kind: \emph{facies totius universi}, quæ, quamvis infinitis modis
variet, manet semper eadem). But the gradation remains within
infinity, and if one compares with these propositions the 28th
proposition of the same book: ``Every single thing, or every thing that is
finite and has a limited existence, cannot exist or
be determined to activity except by another cause which is likewise
finite and has a limited existence; and this cause in turn cannot
exist or be determined to activity unless it is determined
to existence and activity by another, which is likewise finite and
has a limited existence, and so on to infinity'' --- then
not only is the derivation not carried from the infinite to the
finite; \emph{but such a transition is expressly posited as
impossible}, since with the finite the series that runs on to infinity and is never
concluded is given. If one proceeds, therefore, from the infinite,
one arrives only at the infinite, never at the finite; if one proceeds ---
by a hypothesis --- from the finite, one arrives, in the infinite
series, only at the finite, never at the
infinite\footnote{Spinoza's conception of the relation of the different kinds of
knowledge to one another leads to the same consequence (cf. especially
Eth. V, 28). The clear and distinct
knowledge that comprehends under the form of eternity and necessity
(sub specie æterni) cannot arise from inadequate knowledge,
whose object is the imagined finite, but only from another clear
and distinct knowledge. Here is the same express denial of the
transition, which will become evident when the nature and
kinds of knowledge have been developed (cf. chapter 5 and chapter 8.).}.

With the determination modus we have thus not yet stepped out of
infinity. Already in the designation: modus or modificatio, there is
given a being, not in itself, but in another, whose affection it is.
This other is the infinite immanent cause, substance. And through
closer
% --- p. 60 ---
\opage{60}\setcounter{footnote}{0}determinations it is now secured that modus, although being in the
infinite, shall not by its heterogeneity with it come to be \emph{outside} the
infinite and thus cancel it as such. The intermediate determinations
by which Spinoza seeks to remove the incompatibility within the
infinite are chiefly three: the concept of \emph{infinite modi},
the conception of modus \emph{under the form of eternity and necessity} (sub
specie æternitatis et necessitatis), and finally the conception of modi
as a whole, as a \emph{universal individual}, as \emph{natura
naturata}.

If one puts the question thus: how can modifications, i.e.\
determinations, limitations, be thought in the infinite? then
the concept: \emph{infinite} modi, pushes the answer a step back. When % print: Begrebet; (semicolon for colon)
infinite modi are laid down, the question takes the form: how
can modi be thought as infinite? The answer to this is prepared by
the conception of modus under the form of eternity, and is attempted decisively
in the conception of \emph{modi as a whole}. When reason
\emph{thinks} modus, it conceives it not under the determinations of time and space,
--- which are the determinations of imagination, but under the form of
eternity and necessity, in God, i.e.\ according to its
truth.\footnote{For Spinoza the expressions: res in se spectata, ---
vere spectata, --- ad Deum relata, are synonymous.} But when it is
conceived thus, it is conceived not as single and finite, but in its
immediate connection with its eternal infinite cause; its definition
becomes the definition of substance (Ep. 29), its existence essential,
necessary, infinite existence (Eth. II, 45. schol.). Necessary
existence conflicts with the essence of finitude; through the conception of
modi as necessary and eternal one is therefore led to the conception of
the totality of the modifications as \emph{infinite unity}. Limitation
is negation, not something positive; what, when a finite is
presupposed, separates modi from one another in representation is, seen
in thought, something vanishing; but when this negative is removed, the
closed unity is posited. This
% --- p. 61 ---
\opage{61}\setcounter{footnote}{0}gives the determination of the \emph{universal individual} (Eth. II, 13.
Lemma 7, schol.; cf. chapter 4). This universal individual is
\emph{natura naturata}\footnote{\emph{Natura naturata} (the expression was
already used by Giordano Bruno and Hobbes) Spinoza determines (Eth. I,
29 schol.) by its opposition to \emph{natura naturans}. The latter is: id,
quod in se est et per se concipitur, sive talia substantiæ attributa,
quæ æternam et infinitam essentiam exprimunt; hoc est --- Deus,
quatenus ut causa libera consideratur. --- \emph{Natura naturata} is:
id omne, quod ex necessitate Dei naturæ sive uniuscuiusque Dei
attributorum sequitur, hoc est, omnes Dei attributorum modi, quatenus
considerantur ut res, quæ in Deo sunt et quæ sine Deo nec esse nec
concipi possunt.} as a whole (Cog. metaph. II. c. 7. 9.). One cannot say
of it that it \emph{is composite}, one cannot ascribe \emph{parts} to it;
for parts would presuppose a subsisting limitation, composition
the reality of limitation; but the limit is here posited precisely as the
non-real.

Through these intermediate determinations infinity is upheld, finitude and
its forms excluded from substance. Substance is, as causa sui,
\emph{the unity of natura naturans and natura naturata}. As cause it is
naturans, as effect naturata; but in substance, whose essence is
eternity, this difference has eternally vanished: limitation, finitude,
time, number, measure, are no reality, only a semblance, which cannot even
be said to vanish, but, seen in its truth, only \emph{to have
vanished}.

The interest of infinity is thus upheld; but the other
interest, which brought forth the determination modus: that of conceiving
substance as what is infinitely full in itself, not in a resting fullness,
but in one eternally producing
itself, not indeed in such a way that there could be question in its element
of a development, but yet of an
activity,\footnote{Therefore Spinoza can speak of an \emph{infinitus
amor}, quo Deus seipsum amat (Eth. V, 36). For this concept the
view of activity is the presupposition. (Cf. chapter 6 and chapter 8.)}
--- is not satisfied, and this interest, which is
% --- p. 62 ---
\opage{62}\setcounter{footnote}{0}thought's, philosophy's own interest, brings forth inconsistencies
in Spinoza which cannot be removed. In the indeterminate
infinity spirit affirms its abstract-universal essence, denies
itself in its singularity; but its aim is self-affirmation. The act of \emph{philosophizing}
has reality only through singularity, only through the fact that
human thought does not immediately melt together with the
infinite. If infinity is held fast in that abstract form, then
the being of the finite in representation too, however
decisively it is posited as semblance, as essenceless, becomes
incomprehensible.\footnote{It is this contradiction on which already
the Eleatics stumbled. When all non-being is excluded, so that only being is,
and being now as the One embraces thought, the existence of the non-being
in thought as semblance becomes inexplicable, precisely because thought
is posited as one with being.}

By positing activity, and by letting infinitely much follow from it in
infinite modifications, Spinoza sought to escape the emptiness of
abstraction, and to give thought a content and thinking a possibility. But
the fullness vanished in the same moment it was posited; it became a
fleeting semblance, since that which was to give the fullness became a semblance. The
finite things could be posited only hypothetically, as a semblance that casts
its reflection into itself, for they exist only in human
representation, which itself, as finite, is something vanishing. Nothingness
is mirrored in nothingness. Only when the limit receives a reality and
thereby the modifications subsistence can substance receive a content, and
therefore such a significance of the limit must come forth in Spinoza,
precisely through that inner conflict in thought, and even if it is again
dismissed, it nonetheless asserts a kind of validity. Determinations must be sought
for what does not immediately merge into substance. This
cannot, even if it is determined negatively, since substance is determined
as infinite reality, be wholly removed by thought, for thought
would thereby negate itself. When substance is thought, thought has
as it were
% --- p. 63 ---
\opage{63}\setcounter{footnote}{0}the \emph{memory} of finitude; this must force itself forward again and
thereby the inconsistencies arise.

These, therefore, Spinoza has not been able to avoid either. It is already
inconsistent to occupy thought with finitude in its
apparent subsistence; for through this very occupation with it
the finite receives a significance it could not receive according to the fundamental
concepts.\footnote{Yet here the dialectical element comes forth again, that this
significance, as resting on time, ceases again to be significance with the vanishing of time
as something non-real.}
But the
inconsistency comes forth more definitely when there is talk of a \emph{difference in the
finite}; for the old Aristotelian proposition: \textit{τοῦ μὴ ὄντος
οὐκ ἐστιν εἴδη}, always retains its validity.
The difference is posited as a difference of \emph{degree} in the being of the finite; but
such a difference of degree is, even if by holding fast the proposition that
all limitation is negation one forces it back into being merely
quantitative, nonetheless sufficient not to let the finite
vanish entirely. It is likewise an inconsistency when there is talk
of: the \emph{order} in nature, the series of causes, ordo et constitutio totius % print: Aarsagerne Række (genitive s missing)
naturæ, series et ordo causarum, which are distinguished from leges naturæ
æternæ; for in this \emph{series} the single things assert a significance \emph{as}
single. It is an inconsistency when from the active eternal laws
an infinitely \emph{manifold} is derived as consequence (Eth. I, 16). It is an
inconsistency when it is said of natura naturata, which is conceived as unity,
that it is one ``although its parts vary in infinite
ways'' (quamvis partes eius infinitis modis varient); for even if
limitation and change are posited as constantly sublated, they are nonetheless
posited again in order to be sublated. It is an inconsistency when there is talk of
res particulares æternæ, of æterni cogitandi modi, which are determined as
parts of infinitus intellectus; for through eternity the
limitedness that lies in the concepts: singularity and part, is sublated. It
is an inconsistency when
% --- p. 64 ---
\opage{64}\setcounter{footnote}{0}Spinoza ascribes \emph{positive} significance to modi for the knowledge of
God (Eth. V, 24); for although this is not to be understood as if
experience were to supplement an originally abstract idea of God, but that the
idea of God eternally revealed to us comes to consciousness in its fullness and makes
the vanishing imperfect knowledge a semblance, yet, since
positive significance for this full concept is ascribed to modi \emph{as} modi,
that which makes them modi, the limit, receives positive reality. It is
finally an inconsistency to speak of a knowledge of God and an
amor intellectualis issuing from it in the human mind determined as a modus,
albeit as an eternal modus; for in this determination the
heterogeneous, limitedness and infinity, are tied together (cf. chapter 9).

These contradictions thus remain in Spinoza's doctrine of the
modifications. They are not accidental, but necessary consequences of the fact
that \emph{human} thought thinks the concept of substance. The
inconsistency of ascribing significance to the finite, \emph{if only
in order to sublate it again}, is the condition for there being any philosophizing,
for thought not coming to a halt at once in the immobility of
abstraction.\footnote{How much this problem occupied
Spinoza, and how clearly he was conscious of the degree to which his
conception of the finite must be a stumbling-block for his
contemporaries, is seen everywhere in his writings, and especially in the
evasive manner in which he expresses himself on it still in the last years of his life, and
toward the one of his friends who seems to have understood him best
(cf. Ep. 70--72). Just as for Spinoza substance
is the real starting point that is no longer abandoned, so the
finite is a hypothetical starting point that is constantly \emph{only} abandoned; but
precisely herein lies a contradiction.}

\bigskip
\begin{center}\rule{2.5cm}{0.4pt}\end{center}

\newpage

\chapter*{Chapter Three.}
\addcontentsline{toc}{chapter}{Chapter Three. The Immediate Consequences of Spinoza's Conception of the Concept of Substance}
\markboth{The Immediate Consequences}{The Immediate Consequences}

\begin{center}
\textbf{The Immediate Consequences of Spinoza's Conception of the\\
Concept of Substance.}
\end{center}

% --- p. 65 ---
\opage{65}\setcounter{footnote}{0}The constituting determinations of the Spinozist concept of substance:
\emph{infinity} and \emph{causality}, must become the decisive ones for the
further unfolding of the system. The concept of substance was, after all, for Spinoza
the all-embracing concept; its basic stamp must therefore be found again in
every one of his thoughts. In the following chapters this will appear
in the more special developments; at this point it will be
shown with respect to the general metaphysical relations.

We saw that the concept of \emph{substance} in its completeness coincided
with the concept of \emph{God}, and that \emph{God} was determined as the absolutely infinite
being. When infinity is now strictly held fast, there is excluded first
and foremost from the concept of God every representation that could
make it finite in an anthropomorphic way. Therefore there are excluded
from it the determinations: \emph{intellect, will, desire, love}\footnote{Love
as passio is naturally excluded from God. On the other hand, in
the 5th Book of the Ethics there can be talk of an active amor, quo Deus se ipsum amat,
and of which man's amor intellectualis Dei is ``a part.''
This love is to be paralleled with infinitus intellectus
(cf. chapter 5), since it is idea Dei seen from another side.} etc.,
because they do not express the absolutely infinite attribute of
thought, but only a determinate modification of it (cf. chapter 5).
Therefore the fundamental Christian representation of God's becoming
man is excluded, since on Spinoza's conception it is as
self-contradictory that God, the infinite substance, should have
taken on the finite nature of man as that the circle should have
taken on the nature of the square (Ep. 21).

% --- p. 66 ---
\opage{66}\setcounter{footnote}{0}The necessarily existing active infinite (causa sui) must
act according to the necessity of its essence, since outside it there is nothing
that could determine it. This necessity, as inner, is
\emph{free} necessity (libera necessitas, Ep. 62),\footnote{Since this
necessity becomes an inner one, God's very essence as expressing itself
in activity, it does not become fatalism, since the latter separates fatum
from God (Ep. 23. 49. 62.).} freedom; God is causa libera (Eth. I,
def. 7). As one with necessity this freedom is not a
liberum decretum, not indifference (cf. Ep. 60); God's power is
not, like the power of kings, arbitrariness; it is bound by the law of his
own essence (tr. polit. II, 7). The relation to time, the
separation between essence, intellect and will, which a liberum
arbitrium presupposes, conflicts with eternity, by which God's
essence is determined. The essence is the power; the necessary existence
gives the necessary activity. But the concept of necessity
is sharpened by what was the deeper ground of Spinoza's thought.
Substance is the objectified essence of the intellect; but this essence
must, according to Spinoza's historical presupposition, be determined by the
abstract expression for that in which the intellect has its content, i.e.\ by
the abstract concept of nature, and therefore the necessary
causality is determined in such a way that the \emph{concept of purpose} is excluded.
The denial of this concept, which is decisive for the system's
character, is no doubt motivated from the concept of God's perfection
(the concept of infinity), but what is the deeper ground of the polemic,
and what conditions the one-sided conception of the concept of
purpose,\footnote{In the representation that God should act sub
ratione boni Spinoza sees only a subjection of God to a
fatum, since it ``posits something outside God which God in his
acting has before him as a model and strives after as a
goal'' (Eth. I, 33. schol. 2).} is the unconscious
influence of the concept of nature on Spinoza's thought.

In the Appendix to the 1st Book of the Ethics Spinoza develops the
psychological origin of this concept. The representation of
% --- p. 67 ---
\opage{67}\setcounter{footnote}{0}God as acting according to a purpose, he says, is grounded only in
human prejudices, which arise from the fact that men are born
without knowledge of the causes of things, but with appetite for what is
useful to them and with consciousness of this appetite. The consciousness of
appetite, combined with unconsciousness of its grounds, calls forth
the fancy of human \emph{freedom}, and through this
fancy the consciousness of appetite is transformed into a
representation of acting according to a purpose, and this representation
now governs, as conclusive, the conception of things, and since
these are viewed from man's standpoint, and at the same time as
independent of man, it is transferred to the divine,
despite all the absurd consequences and despite all the contradictions
that arise from this representation. But the transference of the
concept of purpose to God turns nature upside down. For that
which in reality is cause is regarded as effect and
conversely; what by its essence is the earlier is made
the later, what is the highest and most perfect the
most imperfect,\footnote{The more immediately something proceeded from God,
the more perfect it had to be posited as being, whereas on that
representation, conversely, the last became the most perfect.} and
finally, --- what is decisive, --- an imperfection is thereby ascribed
to God; for when he acts for the sake of a purpose,
he necessarily desires something which he is not in possession
of. (Eth. I, 36. schol. IV, præf.).

The exclusion of the concept of purpose meets with the exclusion of the
anthropomorphic determinations in removing the representations
that underlie the view of a \emph{personal} God. In its
place enters the \emph{necessity of the power of nature}, which also becomes the
form of thought, and the significance of this concept of necessity for
Spinoza's whole system he himself states when in Ep. 21 he calls
it præcipuum fundamentum eorum omnium, quæ in tractatu, quem
edere destinaveram (the Ethics), habentur, and in his summarizing
% --- p. 68 ---
\opage{68}\setcounter{footnote}{0}survey of the first book of the Ethics (in the Scholium to prop. 36)
puts the determination: quod [Deus] necessario existat, which
is supplemented by: quod ex sola suæ naturæ necessitate sit et agat,
before all the others.\footnote{Eth. IV, præf: æternum illud et
infinitum ens, quod Deum seu naturam appellamus, eadem qua
existit necessitate agit.}

If we pass from natura naturans over to natura naturata, we find
again the stamp of the same fundamental thought. The necessity of causality, which
expresses the essence of the cause, must express itself in the effect, which only
has its being in the cause. When the \emph{infinity} of substance entailed that
everything that is had to be in God, and could not be conceived without
him (Eth. I, 15); that the single things had to be conceived as
affections or modi of God's attributes, by which these are
expressed in a determinate, limited way (Eth. I, 25 Coroll.);
that God is not only the efficient cause of the existence of things,
but also of their essence (Eth. I, 25); then it lies therein that God
was called the cause of all things in \emph{the same sense} as he was called causa
sui, that every determination to activity in the thing is a
determinedness by God (Eth. I, 26); that, since God's acting is in
unity with the necessity of his essence, there can be in nature
nothing contingent (Eth. I, 29); that things could not have been
produced by God in any other way or in any other order than they have been
produced (Eth. I, 33), and that all that we conceive as being in
God's power must necessarily be (Eth. I, 35). And in this existence
held together by the bond of necessity no
independent starting point for an acting can form itself; there can be no
question of an individual \emph{freedom}, of freedom of choice (cf. Ep. 62),
any more than of contingency and possibility. These representations
have their origin only in the imperfection of the human faculty of
knowledge. Validity for things belongs to the concepts of
\emph{necessity} and \emph{impossibility} (although the latter, as mere negation,
naturally cannot be an affection of a being --- entis
affectio. Cog. met. I. c. 3). For a thing is called
% --- p. 69 ---
\opage{69}\setcounter{footnote}{0}\emph{necessary},\footnote{Necessity thus includes, in
its double significance, freedom and compulsion, and one can speak of a
libera and a coacta necessitas (Ep. 62, cf. 58. 60. Cog. met. I,
c. 3).} either with respect to its essence or with respect to its
cause. For a thing's existence follows with necessity either from
its own essence or definition, or from a given efficient
cause (Eth. I, 33. schol. 1). The opposite of the necessary
is the \emph{impossible}, a determination we use when either
the essence (the definition) contains a contradiction or when there is
no external cause that is determined to produce the thing.
On the other hand the designations: \emph{the contingent} and \emph{the possible}, express only a
defect in our knowledge. For when, considering merely
a thing's essence, we find nothing that with necessity affirms
or denies its existence, we call it \emph{contingent}; \emph{possible}
we call it when, looking to the causes by which
it would have to be produced, we do not know whether they are determined to
produce it (Eth. IV, def. 3 and 4).\footnote{Cog. metaph. I,
c. 3. Res \textit{possibilis} --- dicitur, quum eius causam efficientem
quidem intelligimus, attamen an causa determinata sit ignoramus.
--- Si autem ad rei essentiam simpliciter, non vero ad eius
causam attendimus, illam \emph{contingentem} dicemus --- --- quia ex
parte essentiæ nullam in ipsa reperimus necessitatem existendi
--- --- neque etiam implicantiam sive impossibilitatem. In the Ethics
I, 33. schol. 1. the expressions are used as synonyms.}

Just as little as the determinations: contingency, possibility and
freedom, can those determinations find application to natura
naturata which do not express it in its necessary connection with
its cause, but only in its contingent relation to human
imagination. Therefore such determinations are excluded
as order and confusion, good and evil,
perfection and imperfection, beauty and ugliness; and
such determinations as are not ``notiones castæ, quæ
naturam, ut in se est, explicant'', but ``notiones ex vulgi usu
factæ'', which merely express nature as it is referred to
the human senses,
% --- p. 70 ---
\opage{70}\setcounter{footnote}{0}e.g.\ visible and invisible, cold and warm, fluid and solid and the
like (Ep. 6). In nature everything proceeds with eternal necessity and
with the same perfection from God's necessary essence; nothing belongs
to any modus but what lies in its concept, and this
concept is determined not by the contingent, untrue relation to us,
but by the necessary, true relation to God (cf. Eth. II, 32.
Ep. 58). The closer development of this will be given in chapter 5.

If, finally, we pass from natura naturata conceived in its true relation to the
divine, in which, on account of the inconsistencies in Spinoza developed in the
previous chapter, it has nonetheless not been possible to avoid bringing in
determinations that strictly speaking are excluded by its concept, over to
natura naturata separated in abstraction from its cause, to the
infinite series of the hypothetically held-fast finite, then we find
here again the determination of necessity. The detached finite in
its limited existence has not in itself a determining ground
to be and to act, but must be determined to it by another. But
this other becomes here, in the domain of the finite, another
finite, which again demands another finite cause, and so on to
infinity (Eth. I, 28). Through this abstract conception of the
finite, therefore, the determining of things by God (cf. the proof
of the cited proposition) is conceived as their determining of one another, running on
in an infinite series.\footnote{Determinedness
becomes in finite things a \textit{finite} determinedness, hence not
by God, but by something finite.} At this point the system takes on
the character of a strict determinism; but it is easily seen that
determinism in Spinoza must acquire an essentially different significance
than in the systems in which a real existence is ascribed to the single.
Since the single, conceived in its truth, is conceived in
the infinite, the determinedness of the one finite by the other, coming from without,
is transformed into its determining by the
infinite, which is its true essence. Determinism thus becomes
a merely apparent one, and
% --- p. 71 ---
\opage{71} it will appear at a later place (in chapter 6) that the same
holds of the egoism that becomes a consequence of the conception
of the unity of essence and power in the single as well; for since
the limit of singularity is something merely negative, that which
constitutes egoism vanishes.

The representation of necessity as \emph{external} compulsion is thus connected
with the untrue representation of a finite subsisting for itself;
when this representation gives way to the true idea
of the infinite, external necessity becomes an \emph{internal} one, becomes
freedom, i.e.\ determination by the law of its own essence; the causal relation
in substance, its relation to itself as causa sui, embraces
every relation of necessity, just as we saw above that the
concept of substance included the modifications within itself. In God, the
only causa libera (Eth. I, 17. Coroll. 2), the single modi,
conceived under the form of eternity, in their unity with him, receive their
freedom (cf. the 5th Book of the Ethics: de libertate humana). The problem of
the finite is here the same as in the previous chapter; the
difficulties that remain here follow immediately from what was
developed there.

\bigskip
\begin{center}\rule{2.5cm}{0.4pt}\end{center}
\bigskip

\begin{center}
{\large\textbf{Chapter Four.}}\\[4pt]
\textbf{Spinoza's Concept of Nature.}\\[6pt]
\rule{2cm}{0.4pt}
\end{center}
\addcontentsline{toc}{chapter}{Chapter Four. Spinoza's Concept of Nature}
\markboth{Spinoza's Concept of Nature}{Spinoza's Concept of Nature}
\bigskip

When one speaks in Spinoza of a concept of nature, as distinct from
the concept of substance, it must be remembered, as was developed in the second chapter,
that this distinction is a vanishing one, an eternally
sublated one. Substance as causal posits as effect natura
naturata, but in this it posits not another, but itself; as
cause of it, it is causa sui, unity of what effects and what is
effected, and this unity is expressed in that the same designation
determines cause and effect, \emph{natura} naturans and \emph{natura} naturata.

% --- p. 72 ---
\opage{72}\setcounter{footnote}{0}\emph{Nature} becomes the common name; God and nature become synonyms.
``We have shown that nature does not act according to a purpose,''
thus Spinoza refers back to what was developed in Eth. I, 36 schol.,
where this is proved with respect to God. ``Æternum
illud et infinitum ens, quod \emph{Deum seu naturam} appellamus, eadem
qua existit necessitate agit'' (Eth. IV. præf.) --- ``Ipsa \emph{Dei sive
naturæ} potentia'' (Eth. IV, 4. dem.). --- ``Cognitionem unionis,
quam mens cum \emph{tota natura} habet'' (De intell. emend. pag. 360).
--- ``Ex \emph{natura}, sub quovis attributo considerata, infinita
sequi'' (Eth. III, 2. schol. cf. Ep. 21. 58)\footnote{cf. Eth.
IV. 50. schol.: qui recte novit, omnia ex \emph{naturæ divinæ}
necessitate sequi, et secundum æternas naturæ leges et regulas
fieri.}. In these quotations, therefore, natura denotes that which embraces all
attributes, the whole of existence, and produces itself;
its concept is not restricted to what we call nature in
opposition to spirit, even if the name bears traces of the
significance that nature as nature had for Spinoza's conception.

In nature, conceived as synonymous with causa sui, as unity of
natura naturans and natura naturata, there are now included (cf. chapter 2)
the independent attributes and their infinite modifications. Of
the infinitely many attributes only two are perceived by the human
mind (cf. chapter 5), namely \emph{extension} and \emph{thought}, because only
these two are comprised in the idea of the mind. To the development of these two
attributes and their relation the development of the concept of
``nature'' therefore restricts itself. The determinate limited thought, the \emph{idea}, and
the determinate limited extension, the \emph{body}\footnote{Per corpus
intelligo modum, qui Dei essentiam, quatenus ut res extensa
consideratur, certo et determinato modo exprimit. Eth. II, def.
1.}, have their positive subsistence in the attributes of thought and extension,
which both belong to the infinite substance (Eth. II,
1. f.), but as independent. The same substance therefore acts in
both attributes, posits the modifications of both, and that in such a way
that the modifications of the
% --- p. 73 ---
\opage{73}\setcounter{footnote}{0}single attributes form coordinate and \emph{parallel}
series\footnote{cf. Ep. 4. De intell. emend. pag. 378 f., princ.
phil. Cart. I, 21.}. The concept of each single modification includes
only the concept of the attribute of which it is a modification, and has God
for its cause in so far as he is considered solely under this attribute,
and not under any other. Ideas are thus not imprints of
things, not effects of objects, and these are
just as little a consequence of the ideas. The objective being of things follows,
in so far as they are not modifications of thought, not from the
divine nature as thinking, but from the attribute through which
they are conceived. The parallelism in the independent series is grounded
in the fact that the concept of the effect includes that of the cause (Eth. I, axiom.
4); thus \emph{the order and connection of ideas} (i.e.\ of the
modifications of thought) becomes \emph{the same as the order and
connection of things} (Eth. II, 7)\footnote{The concept of a
connection and agreement in the whole of nature Spinoza develops in
Ep. 15. Against the vulgar representation of an ``order'' in nature
determined by its relation to the imagination, on the other hand, he
polemicizes. Eth. I, 36 schol.}. ``Just as the extended and the
thinking substance were seen to be one and the same substance, which
was conceived now under the one, now under the other attribute, so
a modification of extension and the idea of this modification are one
and the same thing, but expressed in two ways. --- --- Thus, e.g.,
the circle that exists in nature, and the idea of the existing circle,
which is also in God, are one and the same thing, which is determined
(explicatur) by different attributes. --- For the objective being
(esse formale) of the idea of the circle can be conceived only through
another modus of thought as proximate cause, this again through another, and
so on to infinity, so that as long as things are considered as
modi of thought, we must explain the whole order of nature
or the connection of causes by the attribute of thought alone;
but in so far as they are considered as modi of extension, the whole order of nature
must just as fully be explained by the
% --- p. 74 ---
\opage{74}\setcounter{footnote}{0}attribute of extension alone, and the same holds of the other attributes.'' (Eth. II,
7. schol.)\footnote{Through this parallelism of thought and
extension the Spinozist fundamental concept again expresses itself:
causality --- in opposition to purpose. In the concept of purpose
lies the priority of thought, the determining of extension by it; this
preponderance is excluded in Spinoza along with the concept of purpose, and in that
parallelism the mode of operation of natural necessity is stamped.}.

If, after having seen the general determination of the
relation of the attributes to one another within nature in the wider
sense, we now consider the attribute that constitutes what we in the narrower
sense call \emph{nature}: the attribute of \emph{extension}\footnote{In Eth. I, 15.
schol. the expression is used: substantia \emph{corporea}.}, whose
concept is presupposed where the essence of the human mind is to be
developed (cf. Eth. II, 13), then it is at once striking how
different Spinoza's conception of the concept of extension becomes
from that of his predecessor Descartes. In the latter, extension was
excluded from God's essence and posited as a determination of
resting, passive matter, or rather: the determination matter
merged into the determination extension, conceived as resting,
so that motion and difference are brought in by a new postulated
act of God; and since extension is conceived as passive, it is conceived
as divisible to infinity. In Spinoza extension is
conceived as \emph{attribute}\footnote{On the relation of extension to God
cf. princ. phil. Cartes. I, 9. schol. Cog. met. II.}, hence
not only as infinite and \emph{indivisible}, but also as \emph{active}. The
last determination is given in that the determination of causality,
under which substance was conceived, immediately acquires validity for
the attributes. Each attribute expresses an \emph{actuosa} essentia; from
each follows an infinity (Eth. III, 2. schol. II, 3. schol.).
Where thought and extension are set against each other by Spinoza,
cogitandi potentia and \emph{agendi} or
\emph{operandi} potentia are sometimes opposed (e.g.\ Eth. III, 28. dem. Ep. 45).

% --- p. 75 ---
\opage{75}\setcounter{footnote}{0}Just as the determination of extension as active is an immediate
consequence of its conception as one of the attributes of substance,
so also are its \emph{infinity} and \emph{indivisibility} (cf. Eth. I, 8.
12. 13. Coroll. 15. schol.). These two determinations, of which
Descartes could admit only infinity, are on Spinoza's
conception inseparable. What deceives us with regard to indivisibility
is ``that quantity is conceived by us in a twofold
way, namely abstractly\footnote{``A substantia abstractam.'' Ep.
39.} or superficially, namely as we imagine it, and
as substance, which happens only through the intellect (intellectus).
When, therefore, we have our attention directed to quantity
as it is in representation, which happens frequently and
more easily, we shall find it finite, divisible\footnote{With
regard to infinite division cf. Ep. 6. De int. emend. pag.
385; princ. phil. Cart. II, 5. schol. --- On the concept: \emph{empty
space} cf. Ep. 6.} and composed of parts; but when we have
our attention directed to it as it is in the intellect,
and conceive it as substance, \emph{which is very difficult}, then
it will be found to be infinite, unique and indivisible. This will
be sufficiently clear to anyone who knows how to make a
distinction between representation and intellect, especially when it is
borne in mind that \emph{matter} is everywhere the same, and that one does not
distinguish parts in it except in so far as we conceive matter
as determined (affectam) in different ways, for which reason its parts
are distinguished only modaliter, not realiter. Water, e.g., \emph{as} water
we can imagine divided, so that its parts are separated,
but not in so far as it is corporeal substance; for
in so far it can be neither divided nor separated. Likewise
water comes to be and passes away \emph{as} water; but in so far as it is
substance, it neither comes to be nor passes away'' (Eth. I, 15.
schol.)\footnote{cf. Ep. 29: As the expression of substance,
extension cannot be limited or divided without its concept
being cancelled. ``Quare ii prorsus garriunt, ne dicam insaniunt, qui
substantiam extensam ex partibus sive corporibus ab invicem
realiter distinctis conflatam esse putant. Perinde enim est, ac
si quis ex sola additione et coacervatione multorum circulorum,
quadratum aut triangulum, aut quid aliud tota essentia diversum
conflare studeat.'' --- At this point the element of spirit
in the concept of substance determines the element of nature, just as the latter at many
other points determines the former.
}.

% --- p. 76 ---
\opage{76}\setcounter{footnote}{0}In the conception of the essence of \emph{extended substance}, therefore,
Spinoza stands in sharp opposition to Descartes. On the other hand he comes
closer to him in the conception of the finite modi of extension (see
below). Not only does he lay down, like Descartes, that
extension receives no other real modifications than those
that are given with motion and the determinations connected with it:
figure and position, so that all other determinations, e.g.\
qualitates tactiles, express only an imaginary relation to our
senses (Ep. 6), but, according to his conception of the finite, he must
posit every determinedness to motion in the single
modus of extension as called forth \emph{from without}, even if the moved thing's
own nature becomes a contributing factor (cf. Eth. I, 28 and II,
13. schol.), and he must let the changes called forth from without in
the world of bodies take place solely according to the laws of mechanics (Ep. 9).
But this agreement again loses its foundation as the
finite, the condition of the mechanical, vanishes.

The expressions: \emph{extension}, \emph{extended substance} (corporeal
substance), \emph{quantity}, \emph{matter}, Spinoza uses as
synonyms, and in so far as the determination extension most nearly
leads to the representation of what is intuited as at rest,
the determination \emph{matter} seems to be a more suitable designation. In
one of his last letters (Ep. 72, of 15 July 1676)\footnote{cf.
the quotation of Eth. I, 15. schol. given above.} he also seems
to be inclined to substitute this latter designation for
the former, no doubt precisely in order to get the active moment
expressed\footnote{That Spinoza did not adopt materia as attribute
instead of extensio may well also have its ground in regard for
the common representation of God.}. Spinoza's utterances
% --- p. 77 ---
\opage{77}\setcounter{footnote}{0}in this letter were called forth by an inquiry from a friend: (Ep.
69) how the existence of bodies determined by motion and figure
could be derived a priori from the concept of extension.
Spinoza first restricts his answer (Ep. 70) to extension
as Descartes had conceived it, namely as a \emph{moles
quiescens}, and shows that on this conception it becomes absolutely
impossible to deduce the existence of bodies a priori, since
matter at rest will continue to be at rest when it is not
set in motion from without by a stronger cause. When
the question is repeated (Ep. 71) and he is urged to state
his own view, he answers evasively (Ep. 72):
``quod petis, % print: opening » missing
an ex solo extensionis conceptu rerum varietas a priori possit
demonstrari, credo, me jam satis clare ostendisse, id impossibile
esse, ideoque materiam a Cartesio male definiri per extensionem;
sed eam necessario debere explicari per attributum, quod æternam
et infinitam essentiam exprimat. Sed de his forsan aliquando, si
vita suppetit, clarius tecum agam. Nam hucusque nihil de his
ordine disponere mihi licuit.'' --- The evasiveness in Spinoza's answer
no doubt has its ground in the fact that the question that was put touched
the two points of his doctrine that were most foreign to the conception of
his contemporaries and the hardest stumbling-block for it, namely:
\emph{the unity of extension, i.e.\ of matter, with God}, and \emph{the non-reality of
finite things}. Seen from substance and in substance, and
that is: seen in its truth, the finite is sublated as finite;
there can therefore be no question of deducing it a priori from the
infinite, indivisible extension; for precisely for a priori
knowledge, which is the knowledge of substance, it does
not exist\footnote{In Spinoza's answer, therefore, emphasis must no doubt also be laid
on the words: rerum \emph{varietas}. For if the single
things are thought as really different, hence as really subsisting
singulars, then they naturally cannot be deduced, since as
such they are untrue and are such only through something negative.}. But
% --- p. 78 ---
\opage{78}\setcounter{footnote}{0}precisely because the way out sought by Descartes, the arbitrary
intervention of a personal God, was cut off, the problem had to
present itself more sharply; the difficulties that emerged in chapter 2 with the
general question about the finite had to repeat themselves here
within the single attribute, and the intermediate determinations
given there had to be sought applied here.

As the \emph{infinite modi} of extension\footnote{In scrutandis rebus
naturalibus ante omnia investigare conamur res maxime universales
et toti naturæ communes, videlicet \emph{motum et quietem}, eorumque
leges et regulas, quas natura semper observat et per quas
continuo agit, et ex his gradatim ad alia minus universalia
procedimus. tract. theol.-pol. pag. 88. --- Corpus humanum non
est absolute, sed tantum certo modo \emph{secundum leges naturæ
extensæ per motum et quietem} determinata extensio. Princ. phil.
Cart. præf. editoris.} there are now determined \emph{motion} and
\emph{rest}\footnote{This: motion \emph{and} rest, must be taken copulatively,
so that it becomes synonymous with: the relation between motion
and rest, which then again finds its expression in facies totius
universi. Rest can, on Spinoza's conception of substance,
only become a relative determination, a difference of quantity within
motion. Only in single bodies can there be talk of a
motion by itself and a rest by itself, but this separation
falls away with the real existence of finite bodies.} (Ep. 66,
cf. Eth. I, 32. Coroll. 2.); as their product, \emph{facies totius
universi}. --- The \emph{finite} modi are the \emph{bodies}. \emph{They are not deduced
from the former}; for such a deduction is, according to what was
developed in chapter 2, impossible, and it is not replaced by the postulate that,
since the nature of the universe is absolutely infinite, ``this
absolutely mighty nature must determine its parts in infinite ways and
compel them to undergo infinite
changes''\footnote{Ab hac infinitæ potentiæ natura eius
partes infinitis modis moderantur et infinitas variationes pati
coguntur. Ep. 15.}; for precisely the concept \emph{variatio} cannot be
derived from the infinite. The finite bodies are \emph{presupposed
hypothetically}, in such a way that the conception of them
% --- p. 79 ---
\opage{79}\setcounter{footnote}{0}as single is held in suspense (Eth. II. def. 7), and \emph{are then led
back} to the infinite modifications\footnote{In the preponderance of the
concept of causality too lies the impossibility of holding fast
a subsisting single; for the effect does not separate itself from the
cause; there is no principle of individuation.}.

In the 2nd Part of the Ethics, 13th Proposition, schol., \emph{the general
propositions of physics} are stated as the foundation for
the subsequent development of the ethical fundamental concepts.
``Simple bodies (corpora
simplicissima)\footnote{The concept: corpora simplicissima, is
relative and arbitrary; for when bodies are conceived as
limited, single, they are separated from substance, and must be
conceived as \emph{divisible to infinity} (Ep. 29). Precisely through this
contradiction of Spinoza's definitely developed conception that
determination shows itself to be hypothetical.} are either at rest or in
motion, and now in quicker, now in slower motion.
Only thereby, not by a difference in substance, are they distinguished from
one another; they all agree in this, that they include the
concept of the same attribute and are determined by motion and rest.''
Since they are conceived as single, the determination to motion and
rest must proceed from another body, which again must be determined by a
third, and so on to infinity (cf. Eth. I, 28. and above
p. 70). The determinedness received they retain until an
opposite one is posited in them by another body, and every such
affection is determined both by the nature of the acting and of the affected
body. With this the determinations are given for the bodies conceived as
simple (simplicissima), which are distinguished only
by motion and rest, quickness and slowness; for the
\emph{composite} bodies (composita) closer determinations are required.
``For when several bodies of the same or different size are
so held together by the others that they lie close against one another
(ut invicem incumbant), or when, moved with the same
or different degree of speed, they communicate their motions to one another in a determinate
ratio, then these bodies can be called
\emph{mutually united}, and are all together said to make up \emph{one
body or individual}, which by this
% --- p. 80 ---
\opage{80}\setcounter{footnote}{0}union is distinguished from the rest.'' The concept of individuality is
thus not an inner determination from an abiding centre; for
such a centre is made impossible by the conception of the finite; ---
but a determination given from without, contingent and relative. Its
relativity now appears also from the extensions that are given it,
and which precisely from the standpoint of mechanism dissolve it at the
subordinate levels\footnote{Spinoza's conception of the limit
as merely negative must also make the concept of individuality relative.}.
Partly it is extended, as regards the body composed of the most simple parts,
in that the individual is regarded as the
same\footnote{retinebit suam naturam absque ulla formæ (i.e.\
of the essence) mutatione.} even if some of its parts are separated off, provided
only that as many others of the same nature come in their place; --- likewise
even if the parts become larger or smaller, provided only that during
the increase or diminution they retain the same ratio of
motion and rest; --- likewise, even if some of the parts change
the direction of their motion, provided they continue to move and
to communicate their motion to one another in the same way
as before, and finally, even if the composite individual as a
whole is moved or rests, or takes this or that
direction, provided only that each of its parts retains its motion, and,
as before, communicates it to the others. Partly the
concept of the individual is extended to what is composed of several different individuals,
whose affections can therefore become so much the more
manifold without harm to the identity of the individual. And since
such new compositions can again be composed to
infinity, one comes to see ``\emph{that the whole of nature is one
individual}\footnote{cf. Cog. metaph. II, c. 7.}, whose parts, i.e.\
all bodies, vary (variant) in infinitely many ways without
any change in the whole individual.''

The direction in this development is thus the same as in the
general problem of the finite. The single bodies are not derived
from infinite extension, but
% --- p. 81 ---
\opage{81}\setcounter{footnote}{0}since one provisionally proceeds from the representation of \emph{single} bodies,
which conflicts with the fundamental principles, the concept of individuality
is given an elasticity that first makes the
determination relative, and finally brings it to vanish
as it coincides with infinity. But during this
retrograde movement, this venous current, which leads everything
back to the centre, without an arterial current leading it out again
from there, there is nonetheless an endeavour, through the very
vanishing, to give the infinite a fullness, and for intuition
there forms a fleeting image of an infinite life of nature, eternally
moved in itself, ordered by eternal laws and embracing everything finite,
in which the thought of the organism is as it were mirrored\footnote{The
view of eternal laws of nature, which govern a
change and bind it under unity, must, carried through
consistently, lead to the view of a mutual communication of an
original motion, a circulation of motion, hence of an
organism; but the concept of the organism is foreign to Spinoza, and
already a reciprocal action cannot be thought in the series of
causes and effects. But that concept is hinted at, and only through
it does thought receive a content.}. But this image is a
vanishing one; for the finite is not merely in the process of vanishing, but has
eternally vanished\footnote{On the relation of this conception of nature to
the Leibnitzian doctrine of monads cf. chapter 10.}.

Only in the form of lemmas to the Ethics has Spinoza
developed the most general propositions of his physics (``generalia in
physicis''). A more detailed elaboration was cut off by his
principle, and however much Spinoza occupied himself with the study
of nature, his deeper interest could nonetheless not rest on
nature in its concrete fullness\footnote{For the conception of the
concrete objects of nature in their particulares differentiæ
receives its interest essentially from the fact that these prove capable of being
reduced to the general mechanical principles (Ep. 6).}, which
vanished in the radiance of substance, just as, according to the belief of antiquity, the
finite had to perish when it had seen the
infinite; --- it could rest only on the abstract
metaphysical concept of nature. And yet it is
% --- p. 82 ---
\opage{82} psychologically and dialectically interesting, in him, as in general in the idealist thinkers of that period, to see how
nature, which is pushed back by thought and is conceived only in
thought's form, because thought everywhere seeks itself, continues
as it were to tempt, so that not only is the thinker drawn to
the study of it, but also the determinations of nature, however volatilized,
unconsciously insinuate themselves into the conception of the
essence of thought. In Spinoza we have already seen samples
of this in the determination of the concept of substance; we shall see it
still more clearly in the determination of the human mind.

The general propositions of physics followed as simple
consequences of the general metaphysical foundation and
the hypothetically taken starting point. The closer carrying-through
of them, had such become possible, would only
have formed a parallel to the Ethics, and would therefore have yielded no more
philosophical gain than we already have in the latter. The Ethics
lent itself, as we shall see, to such a closer elaboration;
in it we shall therefore find the real complement to the
abstract metaphysical foundation.

\bigskip
\begin{center}\rule{2.5cm}{0.4pt}\end{center}
\bigskip

\begin{center}
{\large\textbf{Chapter Five.}}\\[4pt]
\textbf{The Human Mind.}\\[6pt]
\rule{2cm}{0.4pt}
\end{center}
\addcontentsline{toc}{chapter}{Chapter Five. The Human Mind}
\markboth{The Human Mind}{The Human Mind}
\bigskip

Through the second of the two attributes known to man,
the attribute of \emph{thought}, we arrive at the determination
of \emph{the human mind}. Thought is immediately
certain to us as an attribute of God; for every single thought includes its
concept (cf. chapter 4). God is \emph{res cogitans, ens cogitans
infinitum}. The \emph{infinite modification} of the attribute of thought
is \emph{intellectus absolute
% --- p. 83 ---
\opage{83}\setcounter{footnote}{0}infinitus} (Ep. 66)\footnote{This is also designated as intellectus
\emph{actu} infinitus (Eth. I, 30), not in order to distinguish it from an
intellectus potentiâ, which Spinoza does not recognize; but
in order to express that this intellectus is equal to intellectio itself
(Eth. I, 31. schol.).}. As a modus, intellectus, even
as infinite, belongs to natura naturata, not to natura naturans;
it is not the absolute cogitatio itself, but its infinite
effect (Eth. I, 31 f.). Synonymous with intellectus
infinitus\footnote{The difference between intellectus and idea is only
a difference of name, just as, in the sphere of finitude, the difference
between mens humana and the idea which constitutes man's
essence. When intellectus infinitus \emph{Dei} is spoken of, this
expression becomes perfectly parallel to \emph{idea Dei}.} stands \emph{idea
Dei}. Like the former, it follows from ``the absolute nature of the attribute''
(Eth. I, 21 dem.), it expresses God's essence and everything that
follows with necessity from it (Eth. II, 3), and by the unity of God's
essence it is \emph{one} (Eth. II, 4)\footnote{Cog. metaph. II,
c. 7: idea dei, per quam omniscius vocatur, unica et
simplicissima est. Nam revera deus nulla alia ratione vocatur
omniscius, nisi quia habet ideam sui ipsius, quæ idea sive
cogitatio simul semper cum Deo exstitit; nihil enim est præter
eius essentiam, nec illa alio modo potuit esse.} and infinite
(Eth. II. 8). --- The \emph{finite} modi of thought are the \emph{ideas}: per
\emph{ideam} intelligo mentis conceptum, quem mens format propterea
quod res est cogitans (Eth. II. def. 3)\footnote{The idea is called
mentis \emph{conceptus}, not mentis \emph{perceptio}, in order to signify that
the mind does not relate passively to the object, but is active.}.
In the conception of the idea and its relation to intellectus infinitus
and idea dei, the determination of extension and its modi must
serve as a guide. Spinoza views the ideas as limiting
one another just as bodies do (Eth. I, def. 2), he speaks of
the \emph{composition} of ideas, just as of the composition of
corpora simplicissima, and since limitation in the sphere of ideas,
as in that of bodies, is conceived only as negation, the
individual unity of the idea, like the concept of the individual in the
corporeal, becomes a relative concept, from which, as it is
% --- p. 84 ---
\opage{84}\setcounter{footnote}{0}extended through the vanishing of the negative limitation, one arrives at
the universal individual in the sphere of ideas (corresponding to the
universal individual in that of extension), that is: at the
infinite \emph{idea Dei}\footnote{Since intellectus and idea are
synonymous, the same thing is expressed when it is said that all
eternal modi of thought, as a unity, constitute God's eternal and infinite
intellectus (Eth. V, 40. schol.).}, in which all limitation
is cancelled, and which, in the same way as its analogue in
the sphere of extension, must be conceived as unchanged under
the infinite alternation of the parts. When it is said of this idea Dei that
it is in \emph{God} (Eth. II, 3), this is in the same sense in
which everything that is, is in God (Eth. I, 15. 25. Coroll.);
it has its positive subsistence in him as natura naturata in
natura naturans\footnote{In the Ethics V, 30. dem., where
reference is made to Eth. II, 3, idea Dei seems to be attributed to God as
absolute infinitus, but in this place too it is God as
causatum that has this idea of himself as causa, and substance
as causatum becomes equal to natura naturata in its
infinity. Meanwhile the difference between idea Dei and
cogitatio in its immediate, infinite certainty of itself
is on the point of vanishing, and this is also the case where
absolutus intellectus Dei or intellectus divinus is spoken of without
closer determination (Eth. II, 11. Coroll. V, 40. schol.
I, 33. schol. 2. II, 40. dem.). --- As a motive for Spinoza
to guard against intellectus being attributed to God as natura
naturans, one may well assume the endeavor to cut off the
anthropomorphic representations that go with that concept
(cf. Ep. 58. pag. 570 and the line of proof in Eth. I, 33. schol.
2. and I, 17. schol., where one must not overlook that the
argument is ex concessis).}.

Represented as \emph{single}, the ideas stand parallel to the
single bodies, and form a series parallel to the series of
things (cf. chapter 4). God's cogitandi potentia is thus equal
to his actualis agendi potentia, i.e.: what follows objectively
(formaliter) from \emph{God's infinite nature} follows
subjectively (objective) in the same order and with the same connection
in God from \emph{God's idea} (Eth. II, 7. Coroll. cf. Ep. 15; de
emend. intell. pag. 368)\footnote{Therefore the ideas of
the non-existing finite things or modi must also be comprehended in the
same way in God's infinite idea as the objective essence of these things
is comprehended in God's attributes; from which it thus follows that,
so long as the single things do not exist except insofar as they
are comprehended in God's attributes, their subjective
essence (esse objectivum) or their idea does not exist except
insofar as God's infinite idea exists; and when the
finite things are \emph{said} to exist, not merely insofar as they
are comprehended in God's attributes, but also insofar as they are said
to endure (durare), then their ideas also include
existence, by which they are said to have duration (Eth. II.
8).}. To
% --- p. 85 ---
\opage{85}\setcounter{footnote}{0}every thing there thus corresponds an idea, which subjectively expresses
what is objectively in the thing. It is the thing's idea, its
mens, and every natural individual has such a mens; they are
all animate, though in different degrees (Eth. II, 13.
schol.), since the excellence of the idea, i.e. its greater reality, is
dependent on the object's degree of reality.\footnote{For the more
the object of the idea can at once act and be acted upon, the
more apt is the mind to perceive at once
(percipere) several things, and the more the object's actions are
dependent on itself alone, the more apt is the mind
to understand distinctly (Eth. II, 13. schol.).}

At this point, where the presuppositions are given for
determining the essence of the human mind, the problem of
the relation between the two attributes that constitute it and
the infinitely many attributes comes forward again. From the parallelism
between the attributes it would seem that one could draw the
consequence that every thing was determined as a modus of all
attributes, and that the mind apprehended them all, not merely
the modus of extension, the body (cf. Ep. 67). But Spinoza
avoids this last consequence by saying: quamvis
unaquæque res infinitis modis expressa sit in infinito Dei
intellectu, illæ tamen infinitæ ideæ, quibus exprimitur,
\emph{unam eandemque rei singularis mentem constituere
nequeunt, sed infinitas}; quandoquidem \emph{unaquæque} harum
infinitarum idearum nullam connexionem cum invicem habeat
(Ep. 68). But since the modi of all the other attributes are apprehended
through cogitatio, whose modus the idea is, and since the idea becomes
an idea
% --- p. 86 ---
\opage{86}\setcounter{footnote}{0}ideæ, i.e. idea cogitationis, one is led either to split
the attribute of thought into an infinite number of attributes, or
to set the infinite attributes in unity in the attribute
of thought, or to let the latter vanish in the others; in
any case one does not escape conflict with the fundamental view.

By the determinations of the relation of the ideas to things, \emph{the human
mind is} not posited as qualitatively different
from the mentes of other things. It is a determinate modification of
the attribute of thought; what constitutes its actual
being is the idea of an actually existing \emph{single} thing (Eth.
II, 11), and this thing is the \emph{body}, i.e. an \emph{actually}
existing determinate modus of extension (Eth. II, 13). This is
given by the fact that we have ideas of the body's affections,
which could not take place if the body were not the mind's
object. Man thus consists of mind and body, hence
of parallel modifications of the two attributes, and the
meaning of the union of mind and body (unio mentis et
corporis) is clear from the foregoing. The mind is \emph{idea corporis};
but since the idea, albeit in a limited way, expresses the
actuosa essentia of thought, it must be conceived, not as an
image of the body, but as active, and idea corporis thus
becomes an \emph{idea ideæ} (idea mentis), which is united with the mind
in the same way as the mind is united with the body.
The body and the idea of the body were the same thing, apprehended under two
different attributes; the idea and the idea of the idea the same thing,
apprehended under the same attribute, that of thought. The idea of the idea is
the essence of the idea (forma ideæ) insofar as it is considered as a
modus cogitandi without regard to the object (Eth. II, 20 f.).
If there were only an idea corporis, thought would be passive in
the relation; the activity of thought conditions the idea of the idea, the
determinate idea's knowledge of itself.\footnote{The comparison
between Spinoza's conception of idea mentis and the Fichtean
conception, already prepared earlier, of the synthesis of self-consciousness
is instructive (cf. chapter 10). Spinoza's idea mentis is
not abstract self-consciousness (I = I), but the mind's
\emph{concrete immediate self-certainty} in its knowing ---
\emph{self-feeling}, one might say, if this expression did not
contain a formal contradiction, of which, however, Spinoza himself
is guilty when he uses the expression sentire for percipere.}

% --- p. 87 ---
\opage{87}\setcounter{footnote}{0}The conception of the mind as idea corporis,\footnote{Eth. II, 19:
mens humana est ipsa idea sive cognitio corporis humani. ---
Eth. III, 57. schol.: idea seu anima eiusdem individui. --- On
the mind's relation to the body cf. Ep. 15. 68. De intell. emend.
pag. 368. Eth. V. præf.} and the transfer, conditioned thereby, % print: »Eth. V. præf.)« (stray closing parenthesis)
of the relations of extension, of the representation of limitation
and composition\footnote{The idea that makes up the objective
essence of the human mind is not simple, but
composed of many ideas, since the human body
is composed of many individuals, each of which is in turn
composite (Eth. II, 15. 13. post. 1.).} to the modi of thought,
the ideas, now governs Spinoza's doctrine of the human
mind as knowing and willing. It is only by constantly
remembering the parallelism between the two attributes that one understands
Spinoza's theory of knowledge, and it is the view of this parallelism that
underlies the distinction, essential to the theory of knowledge,
between the \emph{adequate} and the \emph{inadequate} idea.
``The human mind is a part of God's infinite
intellectus, and when one says that the human mind
perceives (percipit) this or that, one says nothing
other than that God, not insofar as he is infinite, but
insofar as he is determined (explicatur) by the nature of the human
mind, or insofar as he constitutes the essence of the
human mind, has this or that idea; when, on the other hand,
we say that God has this or that idea, not insofar
as he merely constitutes the nature of the human mind, but
insofar as he, together with that of the human mind, also has
the idea of another thing, then we say that the human mind
perceives a thing partially or \emph{inadequately}'' (Eth. II, 11
Coroll.).\footnote{Eth. III, 1. dem. [ideæ] quæ --- inadæquatæ
sunt in mente, sunt etiam in Deo adæquatæ, non quatenus
eiusdem solummodo mentis essentiam, sed etiam quatenus
aliarum rerum mentes in se continet. --- De intell. emend.
pag. 380: Certum est, ideas inadæquatas ex eo tantum in nobis
oriri, quod pars sumus alicuius entis cogitantis, cuius
quædam cogitationes ex toto, quædam ex parte tantum nostram
mentem constituunt. --- The human mind is composed of
adequate and inadequate ideas. Eth. III, 3. dem.}

% --- p. 88 ---
\opage{88}\setcounter{footnote}{0}In the development of the human mind as knowing,
the same conception of the finite is carried through that we
met in the development of the concept of extension and already
earlier in the development of the general metaphysical
foundation, and the same striving to lead the finite
back to the infinite repeats itself. Just as we saw there
that finitude could be conceived in a twofold way: according to its
truth, as vanished in the infinite, and confusedly, untruly,
as subsisting in and for itself, so it shows itself also
here, and this twofold mode of conception makes the essential
distinction between the \emph{kinds of knowledge}.

If one starts from the hypothetical presupposition of an
existing finite in its multiplicity, this leads to
the kind of knowledge that Spinoza designates as \emph{imagination}. The
human body is affected in manifold ways by the
other bodies, and the mind expresses these affections. Since
the body's affections are determined both by the nature of the affecting and of
the affected body (cf. Eth. II, 13. Lemma 3), their
ideas include the nature of both bodies. ``The affections
of the human body whose ideas represent to us the external
bodies as present we call the images of things (imagines),
although they do not reproduce the forms of things, and when the mind in
this way considers bodies, we call it
\emph{imagining}'' (Eth. II, 17. schol.).\footnote{Eth. III,
56. dem.: quatenus imaginamur, sive quatenus afficimur
affectu, qui naturam nostri corporis et naturam corporis
externi involvit; cf. de intell. emend. pag. 384. Eth. II,
49. schol.: Verborum et imaginum essentia a solis motibus
corporeis constituitur, qui cogitationis conceptum minime
involvunt.} But from the origin of imagination it follows that
the ideas of external things express the state of our
% --- p. 89 ---
\opage{89}\setcounter{footnote}{0}body more than their nature (Eth. II, 16. Coroll. 2 cf. 25
f.), and that they expose us to errors, since the body retains
the affection, and hence the mind the idea of the foreign body as
present, until a new affection excludes the former (Eth. II,
17). From the determination given above of the adequate and
inadequate ideas it follows that through imagination, which
apprehends things as single and distinct, no adequate
knowledge can be reached, since its ideas must all
contain what is not understood from the human mind
alone. For when the finite is conceived as finite, it becomes
a link in the infinite, continuing chain of finite
causes and effects, and as a link it points beyond itself
and requires, in order to be understood, a preceding link as cause, and
thereby it becomes impossible that the idea of the thing held fast as single and
isolated can be adequate; for that would
presuppose a detachment from the series as independent, an
understanding solely through the very nature of the finite link. Although,
therefore, the mind is idea corporis, cognitio corporis, this
idea, considered as single, is nevertheless not adequate, i.e. clear and
distinct, and no more is the idea of the human mind
(Eth. II, 28. schol.)\footnote{Through this contradiction: that
the mind, which according to its essence is cognitio corporis (et mentis),
becomes \emph{non}-knowing, imagination and the
corresponding conception of finitude show themselves as untrue, as
non-being.}. For the mind only knows the human
body, and only knows it as existing, through its
affections; for whereas the human body itself, as
effect, is not understood through itself but through its cause, the
affections that follow from its nature include the body's
concept, and through them the mind therefore knows the body, and
through the ideas of them, itself. But since the ideas of these affections
include the nature both of one's own and of the foreign body,\footnote{Therefore the concept imaginatio can be applied to
our knowledge of \emph{ourselves}.} the knowledge of the former,
just as little as that of the latter, becomes adequate through the human
% --- p. 90 ---
\opage{90}\setcounter{footnote}{0}mind alone, and hence neither does the mind's
knowledge of itself through itself as single. The ideas of the affections,
referred to the human mind alone, are like
consequences without premises, \emph{ideæ confusæ}. For they are
in God, not insofar as he is considered as determined by
the idea of the human mind alone, but insofar as he is
considered as determined by several other ideas.\footnote{In this
relation lies the possibility of the abstract conception of modi
as separated from their cause.} The same consequence follows
with regard to the human body from what is relative in
the concept of the individual, since the parts of the body retain an
independence as separate individuals, each of which must
be understood through its own particular cause, different from the human body
(cf. Eth. II, 19---29). The lack of
adequate knowledge also expresses itself as defective
knowledge of the \emph{duration} of our body and of the external bodies.\footnote{The concept
\emph{duratio} is, in its difference from eternity, a relative concept
which has no significance for true knowledge (cf. chapter
8).} For since their existence as single things does not depend
on their essence, but on ``communis naturæ ordo et rerum
constitutio'', which we cannot survey in its
infinity, the knowledge of existence becomes in high degree
incomplete, and we apprehend things as \emph{contingent} and
\emph{perishable}, determinations which thus express only a
defect of our thought, not a property of the thing.

Through imagination, then, which rests on the untrue conception of
the finite as subsisting, no true
knowledge is attained. It is called forth by the chance collision
of things (rerum occursu fortuito); the mind is determined
from without, and links its ideas together, not according to their eternal and
necessary connection, according to the order of reason, which is the same in all men,
by which it understands things through their first
causes; but by the \emph{memory} of the chance linking
of the
% --- p. 91 ---
\opage{91}\setcounter{footnote}{0}body's affections. Its connecting of the ideas is what
we, in modern terminology, would call that of the
association of ideas (Eth. II, 18. schol. 29, schol.).
Imagination is affected only by the single (De intell. emend.
pag. 383), it apprehends the infinite as finite, limited
and divisible (Ep. 29), is therefore unfit to apprehend the
divine,\footnote{Ep. 60: Deum enim non imaginari, sed
quidem intelligere possumus; cf. Eth. II, 47. schol.} and since
all ideas have their truth only through God, \emph{all}
ideas of imagination\footnote{Cf. Eth. III, 56. dem.}
become inadequate, mutilated and confused (mutilatæ et confusæ); it
becomes the source of ideæ fictæ, falsæ and dubiæ (cf. de intell.
emend.), from which errors arise, as well as of entia
imaginationis, which express not the essence of things but their
untrue relation to us as single beings (see below).

But the untruth of imagination does not exclude a true knowledge;
for imagination became untrue through the untruth in the
finite; but to the finite in its true conception,
conceived such that limitation is a mere negative, and
the finite as being in the infinite, there corresponds the true
kind of knowledge, whose foundation this conception is. For what was untrue
in the former kind of knowledge was not something
positive:\footnote{Cf. de intell. emend. pag. 392; princ.
phil. Cart. I, 15.} there would then be in God a positive false
modus of thought, which is impossible; rather it consisted in the
privation of knowledge that lay in the inadequate,
incomplete and confused ideas. Since the human mind,
whose individuality is relative, makes up only as it were a fragment
of a more comprehensive idea, a part in a composition, its
knowledge is incomplete, as it were mutilated; when
the limitation is cancelled and the whole is posited, i.e. \emph{when the idea is led
back to God,} it is true\footnote{Since all ideas, led
back to God, are true, i.e. adequate, and are seen as
inadequate only insofar as they are referred to the mind of a single thing,
the same law of necessity governs both the adequate and the
inadequate ideas in their consequences. Eth. II, 36.}, i.e.
it agrees perfectly with its
% --- p. 92 ---
\opage{92}\setcounter{footnote}{0}ideatum, and every idea that is complete in us is true. But
that which is common to all things, and which is equally in the parts
and in the whole,\footnote{In the sphere of bodies what was common
was that they all belonged to the same attribute, and that they had
the possibility of motion --- faster or slower --- and
rest. Eth. II, 13. Lemma 3. The meaning of ``the common'' for
the other attributes is evident from this.} does not
exclusively constitute the essence of any single thing, and can therefore only be apprehended
adequately. Hence the consequence that there are certain ideas
or concepts which are common to all men, and which are in all
as adequate, i.e. clear and distinct (Eth. II. 38.
Coroll.).\footnote{Per ideam \emph{adæquatam} intelligo ideam,
quæ \emph{quatenus in se, sine relatione ad obiectum consideratur,}
omnes veræ ideæ proprietates sive denominationes intrinsecas
habet. Eth. II. def. 4. The distinguishing mark of the true idea was
agreement with its ideatum (cf. Eth. 1. axiom. 6.
Ep. 64), and this is a denominatio \emph{extrinseca}. --- As
synonyms for ideæ adæquatæ are used: claræ et distinctæ, e.g.
Eth. II, 28; cf. Eth. II, 13. schol.: adæquate sive
distincte.} They are the ideas of that which is at once
proper to and common to the human body and the
external bodies, and which is equally in the single part and in the
whole. And since such a common element is the foundation of an adequate
knowledge, the mind must become all the more apt
to understand adequately, the more things the body has in common % print: »til et begribe« (misprint for »at begribe«)
with other bodies (Eth. II, 39. Coroll.). The adequate ideas
now become the source of true knowledge; for the ideas that
follow from them must also be adequate.\footnote{On
the significance of consequence for truth cf. de int. emend.
pag. 375 f.} The kind of knowledge proceeding from them is
that of the \emph{intellect} (intellectus). What is peculiar to it is not
merely to be \emph{true}, but to have certainty of its truth,
since the idea, as active, is not something mute like a picture on a
tablet, but the living act of thought, which affirms itself
and immediately gives certainty (Eth. II, 43.
% --- p. 93 ---
\opage{93}\setcounter{footnote}{0}schol.).\footnote{Intelligere is equal to certum esse, idea equal to
ipsum intelligere. Eth. II, 43, schol. Certainty can therefore
be posited as one with the essentia obiectiva itself (the idea), which also
becomes synonymous with truth. De intell. emend. pag. 366 ff.
Intellectus becomes the clearest of all (Eth. I, 31. schol.);
truth its own proof (Ep. 74. De intell. emend. 378.
382). Certainty as something positive is different from the
negative: not to doubt: per certitudinem quid positivum
intelligimus, non vero dubitationis privationem. Eth. II, 49.
schol. The idea \emph{as} idea contains affirmation and negation
(ibid). Since the knowledge of reason contains certainty of its
truth, \emph{probability} is excluded from its domain (cf.
Ep. 15. 60); likewise it is posited as independent of
\emph{experience}, which refers only to finite
\emph{existence} (nam experientia nullas rerum essentias
docet; sed summum, quod efficere potest, est, mentem nostram
determinare, ut circa certas tantum rerum essentias cogitet.
Ep. 28), and can never become a counter-proof against what we understand
clearly and distinctly (princ. phil. Cartes. II, 6. schol.
End.)} Since it apprehends things truly (vere), it apprehends
them not, in the manner of the imagination, as contingent,
but as \emph{necessary}, and since the necessity of things is the very
necessity of God's \emph{eternal nature}, it apprehends them under a
form of \emph{eternity} (sub quadam æternitatis specie. Eth. II, 44.
Coroll. 2).\footnote{Sub specie æternitatis \emph{et
necessitatis}. Eth. IV. 62. dem. --- Eth. V, 29. schol.:
res duobus modis a nobis ut actuales concipiuntur, vel
quatenus easdem cum relatione ad certum tempus et locum
existere, vel quatenus ipsas in Deo contineri et ex naturæ
divinæ necessitate consequi concipimus. Quæ autem hoc secundo
modo ut veræ seu reales concipiuntur, eas sub æternitatis
specie concipimus, et earum ideæ æternam et infinitam Dei
essentiam involvunt. --- Eth. V, 30. dem.: res --- sub specie
æternitatis concipere est res concipere quatenus per Dei
essentiam ut entia realia concipiuntur, sive quatenus per Dei
essentiam involvunt existentiam.} This lies also in the fact that
the foundations of reason (rationis fundamenta) are concepts which
express \emph{that which is common to all things}, and which therefore, since
they do not express the essence of any single thing, must be apprehended without
any relation to time, under the form of eternity. But this
true conception of the single existing things necessarily
includes \emph{the idea of God's eternal and infinite essence}. For the existence
of things is
% --- p. 94 ---
\opage{94}\setcounter{footnote}{0}not thought, separated from its cause, as duration; but
\emph{according to its nature} (``loquor de \emph{ipsa natura} existentiæ''. Eth. II,
45. schol.), as proceeding from the eternal
necessity of God's nature, from which infinite things follow in infinite
modifications. Things are thus thought in their existence in
God; for even though each single thing is determined by another single
thing to exist in a certain way, the force
with which they persist in their existence is nevertheless a consequence of the eternal
necessity of God's nature.\footnote{The leading back of the
single thing's idea to God therefore takes place not through a never
completed approximation through the series of finite things, but
by seeing things in God from the standpoint of infinity, as it were with one glance.
In this conception essence and existence
(\emph{ipsa} existentia, \emph{ipsa natura} existentiæ) are set in
unity, just as they are in substance.} Since the idea of every thing
includes the idea of God, the latter is necessarily adequate and
perfect, and \emph{the human mind thus has an
adequate knowledge of God's eternal and infinite
essence}\footnote{The fact that the human mind knows only two
of God's infinitely many attributes prevents it only from
knowing God completely (\emph{omnino} cognoscere), not from
knowing him clearly and distinctly, since each attribute is known
completely through itself. Cf. Ep. 60.} (Eth. II, 46 f.), and this
infinite being and its eternity must be known to all
(Eth. II, 47. schol.).\footnote{Cf. the enumeration of the
powers and properties of reason in the treatise de emend. intell. pag.
390 ff.}

From the untrue and inadequate knowledge of imagination, then,
the knowledge of reason distinguishes itself as true and
adequate,\footnote{Concerning the general relation between
intellectus and imaginatio cf. Eth. IV. 59. schol. De intell.
emend. pag. 384 f. Ep. 29. 60.} and this distinction is
grounded in the different conception of the finite, which in
the former was abstract, partial, under the form of contingency and
therefore different for different people; in the latter in unity with
the idea of God, under the form of necessity and eternity, and the same
for all. But within adequate knowledge
Spinoza again distinguishes a twofold form: \emph{the knowledge of reason}
(ratio), which
% --- p. 95 ---
\opage{95}\setcounter{footnote}{0}proceeds from notiones communes, and gives adequate ideas of
the properties of things, and \emph{intuitive knowledge}
(scientia intuitiva), which proceeds from the adequate idea of
the objective essence of God's attributes to the adequate knowledge of
the essence of things.

The different kinds of knowledge are thus the following three: 1)
\emph{The first kind of knowledge: opinion} (opinio) or
\emph{imagination}, which partly has its point of departure in the single
things that present themselves to us through the senses incompletely,
confusedly and without the order of reason --- sine ordine ad intellectum
(this Spinoza calls cognitio \emph{ab experientia vaga}); --- partly in
signs, e.g. in the fact that, through words which we hear or
read, an image of the things is called forth in us by the association of
ideas.\footnote{Cf. tract. theol.-pol. pag. 185: intellectio
pura mente \emph{extra verba et imagines}.} --- 2) The \emph{second}
kind of knowledge, or reason (ratio); --- 3) The \emph{third}
kind of knowledge, \emph{intuitive knowledge}. As an example (which, however,
must by its nature be inadequate) to illustrate the difference, he cites
the problem: from three given numbers to find the fourth proportional.
``Merchants have no doubt\footnote{Let the reader recall the difference
between: not to doubt, and: to have certainty (see above p. 93, Note
1).} that one must multiply the second by the third and divide
the sum by the first, because they have not forgotten what they have heard
from their teachers without any proof, or because they have frequently had
this experience with the simplest cases (1st kind). But it can
also be done by virtue of Euclid's proposition, according to the common
property of proportional numbers (2nd kind). But with the simplest numbers none of
this is necessary. For when the numbers 1, 2 and 3 are given, everyone
sees immediately that the fourth number is 6, and much more clearly,
because we infer the fourth from the very ratio which we see
at one glance that the first has to the second'' (Eth. II, 40. schol.
2).\footnote{In the treatise de intellectus emendatione (pag. 362 f.)
the division is somewhat different in form, but nevertheless forms a
parallel to the one given in the Ethics. Spinoza there enumerates 4 modi
percipiendi: 1) perceptio ex \emph{auditu} aut ex aliquo signo; 2)
\emph{ab experientia vaga}; 3) ubi essentia rei ex alia re
concluditur, sed non adæquate, e.g. when we infer from an effect
to the cause, or when an inference is drawn from something universal with
which a certain property always follows; 4) when the thing is understood through
its very essence or through knowledge of its proximate cause. The first two
of these kinds of knowledge correspond to what is later
comprehended under the designation: imaginatio. The third should
correspond most nearly to \emph{ratio}, but some uncertainty enters,
since Spinoza has not distinguished between notiones communes and the
abstract notiones universales (see below), and therefore does not
regard this, but the 4th kind, as adequate and secured against
error. As examples of the third kind of knowledge he cites:
when we have clearly apprehended that we sense this particular body and
no other, we infer with clarity that the soul is united with
the body, which union is the cause of such a sensation; but of
what nature the essence of that union is, we do not understand perfectly
from this, since by the union we as yet understand nothing other than
the sensation of the effect, from which the inference was drawn to the
unknown cause. Another example is: when I know the nature of sight,
and know that it has the property that, according to the different
distance, we see the same thing as larger or smaller, I infer
from this that the sun is larger than it appears; but this is only a
property (proprium), not rei essentia particularis. This
kind of knowledge is thus not sufficient in and for itself, since its
result is abstract. --- ``\emph{Through the very essence of the thing}, on the other hand,
the thing is understood when, from the fact that I know something, I know what it is
to know, or when, from the fact that I have known the essence of the soul,
I know that it is united with the body. By the same knowledge I know
that 2 and 3 are 5, and that two lines which are parallel to one and
the same third are parallel to each other.''}

% --- p. 96 ---
\opage{96}\setcounter{footnote}{0}In order to apprehend more definitely the relation between the three kinds of knowledge,
it is necessary to go somewhat more closely into the Spinozistic
terminology and into his relation to the concepts of earlier
metaphysics.

We saw that rational knowledge took its point of departure from
\emph{notiones communes}\footnote{Eth. II, 40. schol. 1: notionum, quæ
communes vocantur, quæque ratiocinii nostri fundamenta sunt. Cf.
tract. theol.-polit. pag. 430 (ed. Pauli): ad notiones quasdam
\emph{simplicissimas}, quas communes vocant; ibid. pag. 88: res maxime
universales et \emph{toti naturæ communes} (where motion and
rest are spoken of).} and that it was these
% --- p. 97 ---
\opage{97}\setcounter{footnote}{0}that gave the essential ground of distinction between the first and second
kinds of knowledge. It therefore becomes important to determine this
concept, in particular to distinguish it from the notiones universales,
notiones transscendentales, generalia, entia
rationis and the like that belong to imagination. When we remember that notiones communes are
adequate, that they express things according to their truth (vere), and
at the same time recall the proposition that limitation is only negation, it is easily
seen that the proper meaning of: \emph{that which is common}
cannot be: the abstract element common to several really separate
things; for according to their truth things do not exist as single;
but that the expression must designate the \emph{positive unity running through
the many apparently single
things}.\footnote{As a parallel one may here recall Geulincx's
designation of God and the corporeal world as res singulares, outside
which there were only modes of existence. One easily sees that for Spinoza
the determinations realism and nominalism in the scholastic sense
acquire no meaning; nor do they through propositions such as: universalium
--- --- quæ non sunt, nec ullam habent præter singularium essentiam.
Cog. metaph. II, c. 7.} On the whole we must remember that number,
plurality, has no real meaning for Spinoza, since
substance is not determined thereby; the consequences that follow from
the presupposition of number as a reality must therefore vanish.
Notiones communes thus become the ideas of the positive One, and they
are therefore also communes in the sense that they are common to all
the single men; for the existence of these as single, too, is after all
only a semblance; they are only parts of the self-knowing idea. The
subjective and the objective moment are in unity in this concept (cf.
Eth. II, 38. Coroll. with prop. 37, 38 and 47, dem.).

From notiones communes are distinguished first and foremost \emph{termini
transscendentales} (such as ens, res, aliquid), and
% --- p. 98 ---
\opage{98}\setcounter{footnote}{0}the commonly so-called\footnote{For at the beginning of Schol. 2
notio universalis is used in a wider sense. --- As
synonyms Spinoza uses: entia metaphysica (Eth. II, 48. schol.),
ideæ universales (Ep. 50) and notiones abstractæ; the knowledge obtained through those
concepts is a cognitio abstracta vel
universalis (Eth. IV. 62 schol.). Of God's essence we cannot, since he
is not determined by plurality, form a universalis idea (Ep.
50).} \emph{notiones universales} (such as man, horse, etc.). The former have
their origin in the fact that the human body, as limited,
is receptive only to a certain number of images (impressions) in distinct
form, so that the images, when this number is exceeded, begin
to become confused, and when it is exceeded to a high degree, are completely
mingled together. But the human mind can distinctly represent
to itself only as many bodies at once as there can be distinct images
simultaneously in the body. When the images in the body are completely
mingled together, the mind will represent them confusedly, without any
distinction, and comprehend them under one concept, such as res, aliquid.
--- The so-called notiones universales have a similar origin.
For so many images, e.g. of men, are formed at once in the human body
that they do not indeed completely surpass the power of
representation, but yet so far that the mind cannot
represent to itself the small individual differences and their determinate
number, but represents distinctly only that in which they all,
insofar as the body is affected by them, agree, and this
one designates by the universal name, and predicates it of the
infinitely many single ones. But since these concepts have such an
origin, they are not formed in the same way by all, but are
different in each individual, in proportion to that by which the body
has most frequently been affected, and which the mind most easily represents to itself
or remembers (thus at the name man the one thinks first
of the upright figure, the other of man's being an animal
bipes, risibile, sine plumis, rationale, etc.),
% --- p. 99 ---
\opage{99}\setcounter{footnote}{0}so that at bottom these universales notiones are \emph{universales
imagines}, which each forms according to the disposition of his body (Eth. II.
40, schol. 1.).

Next, there are distinguished from notiones communes, since these express something
really existing, the modi imaginandi,\footnote{Entia non rationis,
sed imaginationis. Eth. I, 36. schol. Yet the expression entia rationis
is not seldom used, just as Spinoza too can set, as the opposite of
what is in things, that which [is] merely in \emph{intellectu}
(de intell. emend. pag. 386).} which have no corresponding existence,
but have significance only as aids for imagination in
holding fast, determining and representing things; thus the determinations:
\emph{species, genus; time, number, measure, opposition, order, agreement,
difference} and the like, as well as those designations of something negative which express
a ``confused affirmation'', e.g. \emph{end, limit,
darkness};\footnote{Conversely, the infinite positive, which is only in
intellectu, not in imaginatione, is designated negatively, e.g.
incorporeum, infinitum.} likewise the determinations which designate only
a defect in our knowledge, such as: the \emph{possible}, the \emph{contingent} (see
above); fictions such as: accidentia realia, formæ substantiales,
qualitates occultæ, and finally such determinations as only
express a relation to us as single beings, not a determination of
the essence of things,\footnote{Ep. 6. Notiones ex vulgi usu factæ, vel quæ
naturam explicant, non ut in se est, sed prout ad sensum humanum
refertur.} such as \emph{order-confusion, good-evil, beautiful-ugly} and the like
(cf. chapter 6).

The \emph{general} determinations of imagination thus do not express something
existing and rich in content, but only an unreal abstraction, a
generalized single, whose content is blurred, or they express only
the state of imagination itself, not the nature of things. They are furthermore
differently modified in each individual, since what underlies them
is the single in its difference.

Through the different point of departure the character of rational knowledge
becomes a different one. Notiones communes express,
% --- p. 100 ---
\opage{100}\setcounter{footnote}{0}conceived from the fundamental concept, not a confused abstraction, but
the existing in its full unity; they express at the same time something real,
and, since in thought the unreal separation of the single is cancelled,
they express this real in the same original way in all, and
become universally valid axioms.\footnote{``Opinionibus et documentis
humano generi universalibus, hoc est notionibus \emph{communibus et
veris};'' (tract. theol.-pol. pag. 50). --- Ex axiomatis
intellectualibus per se notis (ibid. pag. 62). --- On the relation
between definition and axiom cf. Ep. 26 f.}

If one starts from this determination of notiones communes, which is the
consistent one according to Spinoza's fundamental view, one has in it an essential
distinction from the determinations of imagination, and intellectual
knowledge can simply be opposed to that of imagination as the one that
rests on the true ideas, whose series proceeds in the mind in
agreement with the series of real objects, in the concatenation of
cause and effect, and in such a way that the mind is an automaton
spirituale acting according to determinate laws.\footnote{De intell. emend. pag.
384. The expression automaton spir. is found already in Descartes
(passions de l'âme Art. 16).} But the difficulty comes forward when
the relation between the second and third kinds of knowledge is to be determined,
and the determination of this relation reacts upon the conception of
the relation between ratio and imaginatio and once more makes the distinction
unsettled. It will appear that the deeper reason for the
difficulties that arise here is the general stumbling block of
Spinoza's philosophy: the conception of the finite.

The knowledge of reason (ratio) is posited as adequate and
true;\footnote{Cognitio \emph{secundi} et tertii [generis] est
necessario vera. Eth. II, 41. cf. V. 28. dem. In the Ethics IV, 31
the designation ratio is used so that it comprises both the second
and the third kind of knowledge.} it apprehends things not as contingent
and untruly single, but as necessary and under the form of eternity. It
leads to the knowledge of God, and the idea of God that it gives is adequate and
perfect. But nevertheless the third
% --- p. 101 ---
\opage{101}\setcounter{footnote}{0}kind of knowledge, intuitive knowledge, is posited as something higher, which
gives a more excellent kind of knowledge (Eth. V, 36. schol.). If
the ground of this superiority is to be stated, it is stated thus: that in
that \emph{universal knowledge} (cognitio universalis) the proof of
the dependence of things on God, both with regard to essence and
existence, is conducted \emph{generaliter} and applied to the single
thing, whereas in intuitive knowledge it is conducted from the very
essence of the single thing. In this conception notiones communes then acquire
the meaning of concepts of what is common to homogeneous things, which
express not the individual essence of things but a
proprietas.\footnote{Cf. Eth. II, 40. schol. 2; 39, dem. V, 12,
dem. and the determination of the different kinds of knowledge in de
intell. emend.} They thus come to express something abstract, and
are therefore brought together with axiomata intellectualia, which are true
but abstract.\footnote{De intell. emend. pag. 380; cf. pag. 386:
ex axiomatis solis universalibus non potest intellectus ad
singularia descendere, quandoquidem axiomata ad infinita se
extendunt, nec intellectum magis ad unum quam ad aliud singulare
contemplandum determinant.} Since the second kind of knowledge is called
universal (``cognitio universalis, quam secundi generis esse dixi''.
Eth. V, 36. schol.), it is designated by the same expression that was used
for the determinations of imagination; its concepts are called, like
the general representations of imagination, notiones universales
(Eth. II, 40. schol. 2); and since universal knowledge is
conceived as general and abstract (de intell. emend. p. 381), but the
general and abstract leads into the domain of
imagination,\footnote{Cf. De intell. emend. pag. 363 Note h., 372, 376.
On pag. 383 f. abstract, universal knowledge is
expressly traced back to imagination.} the limits between
the latter and that of reason become uncertain.

In contrast to this abstract and universal knowledge, the
\emph{third} kind of knowledge proceeds from the concept of God (natura, fons
et origo naturæ, ens unicum, infinitum, quod est omne esse, et
præter quod non datur esse), which cannot be conceived abstractly or
universally, and from the adequate
% --- p. 102 ---
\opage{102}\setcounter{footnote}{0}knowledge of the attributes is derived the adequate knowledge of
the essence of things, which follows from the necessity of the divine nature.
It apprehends the ideas in their concatenation as one idea, so
that the mind subjectively expresses the objective essence of nature; its
knowledge of the effects becomes only a fuller knowledge of
the cause. It derives nothing from the abstract, but from ``essentia
particularis affirmativa'',\footnote{De intell. emend. pag. 388:
Optima conclusio erit depromenda ab essentia aliqua particulari
affirmativa. Quo enim specialior est idea, eo distinctior, ac
proinde clarior est. Unde cognitio particularium quam maxime nobis
quærenda est.} and ``ens physicum et reale'', or from ``a true and
lawful definition'', i.e. one which states not merely properties
of the thing, but the thing's intima essentia, from which all its
properties can be derived, and which does this in an affirmative form.
The series that knowledge forms in the ideas is ``the series of the fixed and
eternal things'', not the unsurveyable ``series of single, changeable
things'', since the existence of these things has no connection with their
essence, or, in other words, is not an ``eternal truth'' (æterna
veritas). These ``abiding and eternal things'', which are all simul natura,
are, although single, nevertheless, on account of their omnipresence
(ubique præsentia) and all-embracing power (latissima potentia), for
us as it were universal or general definitions (genera
definitionum) of the single and changeable
things,\footnote{Thus ``the laws of nature extend to infinitely
many things, and are apprehended by us under a form of eternity'', tract.
theol.-polit. pag. 72.} and proximate causes of everything; in them, as
in the true law-books, are inscribed the laws according to which all single
things come into being and are determined; indeed these changeable and single things
are so inwardly and essentially (intime et essentialiter) dependent
on those abiding ones that they can neither be nor be understood without them
(de intell. emend. 388 ff.).\footnote{Concerning the conception of
the existence of things through God's essence cf. Eth. V, 30, dem. 36, schol.}

% --- p. 103 ---
\opage{103}\setcounter{footnote}{0}If, however, one compares what has thus been posited as
peculiar to the third kind of knowledge with what is
peculiar to the second, the distinction becomes difficult. The latter too
is supposed to be able to give clear and adequate ideas, and to derive from these
other clear and adequate ones. But in deriving ideas from others,
it knows the effect from the cause, and this cannot merely
be conceived thus, that ratio, within the concatenation of the single ideas,
should be able to extract adequate ideas, but not set
them in connection with God, since this would conflict with the concept of
adequate apprehension; for it apprehends things as they are according to their
essence, and according to their essence they are in God\footnote{Cf. Eth. V, 29.
schol. 30, dem.}. Knowledge is adequate precisely only when the effect
is led back to God, who is the causa absolute proxima of all things
(Eth. I, 28), in whom they are contained, and from the necessity of whose nature
they follow (Eth. V, 29. schol.). But when this happens, the
difference between ratio and scientia intuitiva would vanish, since God,
who cannot be conceived abstractly or universally, immediately shows
himself as what really is in the thing, and the thing, apprehended adequately,
hence under the form of necessity and eternity, is immediately led
back to the first cause. --- If one would, as might seem
hinted in the treatise de intell. emend., by the second
kind of knowledge think chiefly of the inference from effect to
cause or from abstracta to abstracta, then the former would lead
to ``mutilated and confused ideas'', the latter to the universalia
of imagination, hence in neither case to an adequate
knowledge.

The difficulty arises from the ambiguous position of the finite.
If the strict concept of infinite substance is held fast, then, as
was developed above in chapter 2, one can neither immediately nor mediately
arrive at something finite, hence not at a subsisting plurality, but
the latter is the presupposition of the abstract, universal conception.
The object of knowledge would then be one, the apprehension of it in its
% --- p. 104 ---
\opage{104}\setcounter{footnote}{0}fullness necessarily intuitive, since its being is undivided, infinite and
eternal. The adequate idea would have only the infinite as its object,
for it expresses its object under the form of eternity and necessity\footnote{Thus no agreement could arise
between intellectus and imaginatio, which Spinoza nevertheless concedes
(cf. de intell. emend. pag. 384, Ep. 29. Eth. V, 12), although
the limits are not drawn and cannot be drawn. Since the mathematical
figures are called entia rationis, i.e. imaginationis (Ep. 72; de intell.
emend. pag. 387), they become a point of contact.}, and therein lies
infinity; it is the affirmation of thought, but the limitation of the finite
is a negation. But the borrowed, fleeting existence that
must be given to the finite (cf. chapter 2) gives a point of attachment for the
kind of knowledge that belongs to mathematics, and for which Spinoza's
historical situation must have made it natural to him to claim a place.
This, however, he cannot procure without confounding
the determinations of reason and of imagination. The foundation of imagination
is the single, abstract as single, which does not exist for
reason; abstract knowledge presupposes such a single,
even if only as leaving it behind; otherwise it would not become abstract,
universal, but individual and intuitive. The indeterminate thus comes
in this way to stand in a doubtful relation to the infinite,
``the true abstract'' in a doubtful relation to the untrue abstraction of
imagination; but this uncertainty and contradiction was
the consequence of the uncertainty in the metaphysical relation between the
infinite and the finite. According to the consistent conception of
the concept of substance, the second kind of knowledge would have to vanish in the
third; but in this infinity the individual thought itself
vanishes in the infinite, philosophy becomes inexplicable, the
philosophizing thought comes into contradiction with itself, and this
contradiction forces upon the finite a kind of subsistence, which nevertheless cannot
be determined without new contradictions coming forward.

The doctrine of the different forms of \emph{knowledge} is decisive for
Spinoza's conception of the human mind,
% --- p. 105 ---
\opage{105}\setcounter{footnote}{0}whose essence was determined by knowledge, and for which there was no
principle of individuation. The \emph{idea} is, as a modus of thought, by its
nature prior to the other modi of thought\footnote{Eth. II. axiom. 3.:
``Modi cogitandi, ut amor, cupiditas, vel quicunque nomine affectus
animi insigniuntur, non dantur nisi in eodem individuo detur idea
rei amatæ, desideratæ \&c. At idea dari potest, quamvis nullus alius
detur cogitandi modus.''}. Since it is conceived as active, it includes % print: closing „ with no opening mark in this note
the \emph{will} within itself: ``will and intellect are one and the same''
(Eth. II, 49. Coroll.). For the will (voluntas) is determined by
Spinoza as facultas affirmandi et negandi, i.e. as the faculty by
which the mind affirms and denies the true and the false (not as
the desire by which the mind desires or shuns a thing). But
the determination: \emph{faculty}, is only an ens metaphysicum sive universale,
which we form by abstraction from the single; and voluntas therefore
comes to relate to the single volitiones as lapideitas
to lapis, i.e. it is in reality not different from them. But since
the will thus becomes the single affirmation or negation,
it becomes equal to the idea, since in the mind there is no other affirmation and
negation than that which the idea, \emph{as} idea, includes. Thus, e.g.,
the volitio by which the mind affirms that the triangle's three
angles are equal to two right angles is not different from the affirmation that
lies in the idea of the triangle, which idea is the condition of that
affirmation, and in turn cannot be thought without it, so that their
concept coincides (Eth. II, 49. dem.). The single volitio thus becomes
equal to the single idea, voluntas to intellectus.
Conceived as single, volitiones, like every other single
thing, must be determined to existence and action by another single thing; they
are thus necessary (Eth. I, 32), and the human mind,
conceived as a determinate and limited modus of thought, can no more
have an absolute faculty of willing and not willing (an absolute and free
will) than it can, under the same conception, have an absolute faculty
of understanding (Eth. II, 48); it cannot become the free cause
% --- p. 106 ---
\opage{106}\setcounter{footnote}{0}(causa libera) of its actions, but is determined from without with
necessity.

In the conception of the will, therefore, Spinoza must distinctly separate himself from
Descartes, who, since he posited in God a freedom of arbitrariness,
at the same time, despite the consequences of his fundamental view, conceded
a similar one to man, and, by determining the human will
as infinite, hence wider in extent than limited
knowledge, sought a ground of explanation for error, whereby
the latter could be conceived, in relation to God, as a mere negation, and only in
relation to man became a privation. Spinoza indeed concedes that
the will has a wider extent than the clear and distinct ideas; but
not than perceptiones sive concipiendi facultas in general, since
this too, in the same way as the will, can \emph{successively}
take in infinitely many things. The conception of the will as infinite he derives from
the misunderstood use of the abstract concept: faculty, that is applied
to it; for since the abstract, the universal, is predicated equally of
one, of several and of infinitely many particulars, whose common element it
expresses, one can give it an infinite extension beyond the domain of
determinate knowledge; and the same abstract
conception\footnote{Eth. II, 49. schol.: hic apprime venit notandum,
quam facile decipimur, quando universalia cum singularibus et
entia rationis et abstracta cum realibus confundimus.} likewise brings it about
that the will, determined as affirmation, is predicated in
general of every idea without regard to its content, and of the
true and the false, so that with the power of the \emph{same} will we affirm
all ideas, and predicate the true and the false as true, whereas
the ideas according to their essence contain a different degree of reality,
so that the affirmations contained in them must also include
a similar one, and whereas the true and the false relate to
each other as ens to non-ens. The faculty of arbitrarily
suspending one's judgment, too, Spinoza denies, since suspending one's judgment
becomes nothing other than that we
% --- p. 107 ---
\opage{107}\setcounter{footnote}{0}see that we do not understand the thing adequately. Even in imagination
there lies an affirmation, as long as another representation does not
exclude the existence of the represented object. As long as this
does not take place, one cannot doubt the reality of one's representation,
although: \emph{not to doubt}, is not synonymous with:
\emph{to have certainty}, which only the true idea gives. If I represent
to myself, e.g., a winged horse, then I predicate (affirm) wings
of the horse, and I will, according to the nature of imagination, regard the
winged horse as existing as long as I have this representation
alone, and only doubt the representation or deny it when another
idea cancels this representation or shows it to be inadequate;
but then my doubt and my denial will be \emph{necessary}. --- The ground
for the freedom of the will that is drawn from its being indispensable in order to
cancel the perfect state of indifference, Spinoza rejects, mockingly, since it is
incompatible with his conception of the activity of thought:
dico, me omnino concedere, quod homo, in tali æquilibrio [like
Buridan's ass] positus (nempe qui nihil aliud percipit, quam sitim
et famem, talem cibum et talem potum, qui æque ab eo distant) fame
et siti peribit. Si me rogant, an talis homo non potius asinus quam
homo sit æstimandus, dico me nescire, ut etiam nescio, quanti
æstimandus sit ille, qui se pensilem facit, et quanti æstimandi sint
pueri stulti, vesani \&c.

After Spinoza, in the 2nd book of the Ethics, has developed his doctrine of the
nature of the human mind, he emphasizes, against the objections that
would naturally occur to his contemporaries, the advantages of his conception, and
he places these in: 1) that it, by teaching us that in our action
we are unconditionally dependent on God, and in our essence partakers of the
divine nature, and that in all the higher degree the
more perfectly we act and the more we understand God, --- calms our
mind and shows us wherein our highest happiness or blessedness consists,
namely \emph{solely in the knowledge of God}, by which we are led to
do only what love and piety bid, and to regard
virtue itself and
% --- p. 108 ---
\opage{108}\setcounter{footnote}{0}dependence on God as reward, freedom and beatitude; ---
2) that it gives us equanimity in the face of what is not
in our power (i.e. what does not follow from our nature), because we know
that everything follows with the same necessity from God's eternal decretum
as the properties of the triangle follow from its essence; --- 3) that it
contributes to furthering social life, since it teaches us not to
hate, despise, mock, be angry or envy, and since it makes
everyone content with his own and inclined to help his neighbor,
not from womanish pity, partiality or superstition, but
solely at the prompting of reason; --- 4) finally in this, that it
is of no small use to society, since it shows in what
way the citizens must be governed and led, namely so that they
do not become slaves, but freely do what is best (Eth.
II, 49. schol.). --- What Spinoza has here emphasized as the ``use''
of his doctrine will find its closer development and understanding in
the following chapters. For through the determination of the kinds of knowledge,
and of the relation between the idea and the other modi of thought, such as will,
desire and the like, the foundation is given for the conception of
what specifically makes up the content of the Spinozistic
ethics\footnote{``Ethica universalis'', tract. theol.-polit. pag. 46.}:
the doctrine of \emph{man in his bondage and in his freedom}, and of
this the next three sections will give an account.

\bigskip
\begin{center}\rule{2.5cm}{0.4pt}\end{center}
\bigskip

\begin{center}
{\large\textbf{Chapter Six.}}\\[4pt]
\textbf{Man as naturæ pars. The Affects. Man's
Bondage.}\\[6pt]
\rule{2cm}{0.4pt}
\end{center}
\addcontentsline{toc}{chapter}{Chapter Six. Man as naturæ pars. The Affects. Man's Bondage}
\markboth{Man as naturæ pars}{Man as naturæ pars}
\bigskip

The thought of necessity, which in Spinoza governs
the conception of nature in its unity with God and as the
infinite series of individual things, must also govern
% --- p. 109 ---
\opage{109}\setcounter{footnote}{0}the conception of the human being as acting and
suffering. Since Spinoza proceeds from the concept of substance,
he cannot conceive man as detached from nature's
eternal, all-embracing order, as a state within its state, governed
and determined only by his own groundless, unconditioned freedom.
He cannot regard that which, as occurring in the
human mind, betrays its impotence and inconstancy,
as something accidental, as an object of lament,
mockery and abhorrence, but only as something that must be comprehended in
its necessity through the power of nature, that must be derived with mathematical
strictness from the natural
causes\footnote{cf. tract. polit. Cap. I, 4: ut ea, quæ ad
hanc scientiam [i.e. politicam] spectant, eadem animi libertate,
qua res mathematicas solemus, inquirerem, sedulo curavi,
humanas actiones non ridere, non lugere, neque detestari, sed
intelligere: atque adeo humanos affectus, ut sunt amor, odium,
ira, invidia, gloria, misericordia et reliquæ animi
commotiones non ut humanæ naturæ vitia, sed ut proprietates
contemplatus sum, quæ ad ipsam ita pertinent, ut ad naturam
aëris æstus, frigus, tempestas, tonitru et alia huiusmodi, quæ
tametsi incommoda sunt, necessaria tamen sunt, certasque
habent causas, per quas eorum naturam intelligere conamur, et
mens eorum vera contemplatione æque gaudet, ac earum rerum
cognitione, quæ sensibus gratæ sunt.}. ``Nothing happens
in nature that can be ascribed to a defect in it;
for nature is always the same and everywhere one, and
its power (virtus) and power of action are always the same;
that is to say: the laws and rules of nature, according to which everything
happens and everything is changed from one form into another, are everywhere
and always the same, and the conception of the nature of any thing
whatsoever must therefore always be the same, namely
through the universal laws and rules of nature. The affects of hatred, of anger,
of envy and the like therefore follow, when they are considered
according to their essence (in se), from the same necessity
and power (virtus) of nature as the other individual things,
and therefore have definite causes, through which they are known,
and certain properties, which are just as worthy of being comprehended
by us as the properties of any other thing, in the mere
contemplation of which we delight. I shall therefore treat of the
% --- p. 110 ---
\opage{110}\setcounter{footnote}{0}nature and strength of the affects and the mind's power over them in the same
manner in which I have treated of the essence of God and of the mind,
and I shall consider human actions and
appetites just as if the question were of lines, planes or
bodies'' (Eth. III. præfat.)\footnote{Eth. IV. 57. schol.: Sane
humani affectus, si non humanam, naturæ saltem potentiam et
artificium non minus indicant, quam multa alia, quæ
admiramur, quorumque contemplatione delectamur.}.

The ethics that can find its place in the Spinozistic
system cannot, then, rest on the presupposition of an
individual freedom, cannot set \emph{tasks} which freedom
is to realize through a choice, cannot speak of a good or evil determined by this
choice. It becomes \emph{the natural science of the
human mind}; its task is not to make
a demand on man, but to see the human
being in the light of nature's eternal and necessary laws, and, by
these laws, which express an activity, to derive what
follows with eternal necessity from the concept of the human
being. This task at once takes on a different
form according to the twofold conception of the finite, inasmuch as
man's essence is considered either as abstractly finite
or in its unity with the infinite. In the former sense it is seen
under the law of external necessity, in the latter under that of
internal necessity, of freedom, since it is seen in unity with
the only causa libera; in the former, then, man is conceived in
his state of bondage, in the latter in his freedom. Spinoza keeps
these two conceptions apart from each other, except
insofar as the general determinations can be applied
to both points of view and the definitions comprise both.
He first develops the consequences of the first
way of considering, in order through them to lead it back to
the latter.

If man is first conceived in general as a single
modus in the infinite, then he is determined as \emph{acting}
and \emph{suffering}. ``I call it \emph{acting} when something happens
in us or outside us of which we are the \emph{adequate cause},
% --- p. 111 ---
\opage{111}\setcounter{footnote}{0}i.e. when something follows from our nature, in us or outside
us, that can be understood clearly and distinctly through it alone. On the other
hand I call it \emph{suffering} when something happens in us,
or when something follows from our nature of which we are only
the partial (inadequate) cause'' (Eth. III, def.
2)\footnote{``\emph{Causam adæquatam} appello eam, cuius effectus
potest clare et distincte per eandem percipi. \emph{Inadæquatam}
autem seu \emph{partialem} illam voco, cuius effectus per ipsam
solam intelligi nequit (Eth. III. def. 1). cf. Eth. IV. App. C.
1: Omnes nostræ cupiditates ex necessitate nostræ naturæ ita
sequuntur, ut vel per ipsam solam tanquam per proximam
causam possint intelligi, vel quatenus naturæ sumus pars,
quæ per se absque aliis individuis non potest adæquate
concipi.''}. % print: closing quotation mark missing

Since from the essence of every existing thing an effect always
follows, man must always be determined, and
this determinateness is the \emph{affect}\footnote{Eth. III, 56. dem.:
prout unusquisque --- --- afficitur, hoc est, prout eius natura
hoc aut alio modo constituitur. --- Eth. III. Affect. def. 1.
explic.: per affectionem humanæ essentiæ quamcumque
eiusdem essentiæ constitutionem intelligimus, sive ea sit
innata, sive quod ipsa per solum cogitationis sive per solum
extensionis attributum concipiatur, sive denique quod ad
utrumque simul referatur.}: per \emph{affectum} intelligo corporis
affectiones, quibus ipsius corporis agendi potentia augetur
vel minuitur, juvatur vel coërcetur, et simul harum
affectionum ideas (Eth. III. def.
3.)\footnote{This definition is general, since the affect
becomes an acting when we are the adequate cause of the
affection, in the opposite case a suffering (ibid. explic.).}.

This general definition of the affect now receives its
closer determination through what was developed in the preceding chapters
about the nature of the mind, about the relation of the attributes
to each other, and about substance as acting. The mind
is acting insofar as it has adequate ideas, since
their effects are comprehended through the essence of the human
mind alone; it is suffering insofar as it has inadequate ones,
i.e. insofar as it ``has something that includes a negation'', or
insofar as it is considered as a part of nature which
cannot be understood clearly and distinctly without other parts. The
% --- p. 112 ---
\opage{112}\setcounter{footnote}{0}more inadequate ideas the mind has, the more, then, is it
exposed to passions (passive states); the more adequate,
the more is it acting (Eth. III, 1. 3. schol. cf.
chapter 5 and Eth. II, 11. Coroll.). Reason and acting, imagination
and suffering become parallels. --- From the relation
of the attributes to each other it follows that neither can the body determine
the mind to think nor the mind the body to motion
or rest (Eth. III, 2), so that the acting and suffering states of the body
and of the mind are not consequences of each other,
but are simul natura, and the mind's appetite and the body's
determination are one and the same thing, which is only designated by
different names (decretum --- determinatio) according as it
is seen under the two different attributes. --- From the nature of
substance as active it follows that the individual modi, as things which
in a definite and limited way express God's
attributes, also express his potentia (which is equal to his
essence), and since a thing's essence can only affirm, not
deny, itself, every thing can be annihilated only by
an external cause, and it must, as far as it is in itself (in se),
strive\footnote{The conatus is conceived as a striving \emph{to
preserve} what is already present: Eth. III, 12. 13. cf.
25. 7. Coroll. 3. dem. 28 f.} without any temporal limitation, to % print: a damaged digit stands before `7.'
persist in its
being\footnote{Unaquæque res, quantum \emph{in se est}, in
suo esse perseverare conatur (Eth. III. 6). The expression \emph{in
se} is striking here, since strictly speaking it can be said only of
substance that it is in se. There perhaps lies in it an
intimation that the thing conceived as individual is active only in \emph{its
unity with substance}. Conceived as an individual link in
the infinite chain, the human mind could, consistently,
no more be conceived as active than it can have
adequate ideas. What Spinoza develops from a Cartesian standpoint
in Cog. metaph. II. c. 12 finds no application to
his own, which conceives the mind as modus, not as
\emph{finite substance}. --- Perhaps, however, the expression:
quantum \emph{in se} est, might be merely an incorrect usage for:
quantum in \emph{ea} est, which expression is always used in the tract.
theol.-polit. and the tract. polit. where this thought comes forward
(cf. also Eth. IV. App. c. 25). The old Dutch
translation gives no enlightenment, since it translates
quite according to the words: ``Yder ding poogt, voor zo veel als't in
zich is, in zyn wesen te volharden.''}, and this striving is % print: opening quotation mark missing
the very essence of the thing (Eth.
% --- p. 113 ---
\opage{113}\setcounter{footnote}{0}III, 6 f.) and is found in the human mind both insofar
as it is constituted by adequate and by inadequate
ideas\footnote{One must remember here that Spinoza, in
considering the mind as the idea of the existing body,
conceives this idea as composed of many other ideas, of
which some are adequate, others inadequate
(cf. Eth. III, 3. dem. II, 15. 19. 38), and that the inadequate ideas,
which, referred back to God, are adequate, develop with
the same consistency as the adequate ones (Eth. II, 36).}.
This striving, which is thus the simple expression of
the actual essence as active, is the first fundamental affect.
If it is referred to the mind alone, it is called
\emph{will} (voluntas); if it is referred both to mind and
body, it is called \emph{appetite} (appetitus), and this is
thus nothing other than the very essence of man, from whose nature
that which serves for its preservation follows with necessity,
in such a way that man is determined to bring it about. From
appetitus cupiditas is distinguished only in that it is as a rule
referred to man insofar as he is conscious of his appetite
(cupiditas est appetitus cum eiusdem conscientia,
Eth. III, 9. schol.)

To appetite as fundamental affect, when the
human being is conceived as an individual, limited link
in the infinite chain of causes and effects, and its
power of action is thus conceived as one which can be both increased
and diminished, two others must now attach themselves: the one which expresses
the transition to a greater, and the one which expresses the transition
to a lesser perfection: \emph{joy} and \emph{sorrow}\footnote{\emph{Lætitia}
is: passio, qua mens ad maiorem perfectionem transit;
--- \emph{tristitia} is: passio, qua mens ad minorem transit
perfectionem (Eth. II, 11. schol.) cf. Eth. III, 57. dem.:
lætitia et tristitia est ipsa cupiditas sive appetitus, quatenus
a causis externis augetur vel minuitur, juvatur vel coërcetur,
hoc est, est ipsa cuiusque natura. They are referred most nearly to
the mind alone, \emph{hilaritas} and \emph{melancholia}, on the other hand, to the whole
human being.}. Only
these and appetite can be fundamental affects; from them
all others must be derived.

% --- p. 114 ---
\opage{114}\setcounter{footnote}{0}Since the mind is determined as idea corporis, it follows,
with regard to the three fundamental affects, that there can be
no idea in the mind which excludes the existence of the body, since
such an idea would conflict with the essence of the mind, but that, on the
contrary, the nature of the mind must first and foremost express itself in
affirming the existence of the
body\footnote{When the mind is conceived as idea ideæ,
it expresses itself as a striving to represent to itself what
affirms its own power of acting, and when it considers itself
and its power of acting, it feels joy (Eth. III, 53 f. cf.
V, 15. dem.), whereas it feels sorrow when it
represents to itself its impotence, i.e. when its striving to
represent to itself what affirms its power of acting is hindered
(Eth. III, 55. dem.)}, and that it, since the idea of
all that increases or diminishes our body's power
to
act\footnote{Spinoza sets it up as a postulate that the
human body can be affected in many ways by
which its power of action is increased or diminished, and
likewise in other ways which neither increase nor
diminish it. He refers to the lemmas from
physics (Eth. II, 13), and from this reference it seems to
follow that the increase or decrease of the power of action
must be thought to take place through the parts of the body, while
the mutual communication of motion is retained, becoming
larger or smaller, and through the body as a whole, likewise
with that relation retained, being set in motion or
rest (Eth. III, postul. 1).} must increase or diminish the mind's own
power to
think\footnote{``Quum supra dixerim, mentis cogitandi
potentiam augeri vel minui, nihil aliud intelligere volui, quam
quod mens ideam sui corporis vel alicuius eius partis
formaverit, quæ plus minusve realitatis exprimit quam de suo
corpore affirmaverat. Nam idearum præstantia et actualis
cogitandi potentia ex objecti præstantia æstimatur'' (Affect.
gener. defin., explic.).}, must have a striving to hold
fast the idea which affirms the essence of the human body
and to remove the one which denies
it\footnote{The latter happens through an opposite idea being remembered
(held fast).} (Eth. III,
10---13). With the last determination we have already arrived
at the first of the derived affects: \emph{love} and
\emph{hatred}, which in reality are nothing other than \emph{joy
and sorrow accompanied by the idea of the external cause}
% --- p. 115 ---
\opage{115} (Eth. II, 13. schol.), and whose peculiar character becomes clear
from what has already been developed.

In the deduction of the derived affects the
Spinozistic fundamental view comes forward with a sharply marked peculiarity.
In it he gives a \emph{natural history of the
passions}, and regards, with the naturalist's dispassionate
calm, passion as a natural phenomenon proceeding with necessity
from the human being. In the fundamental affects
this being found, in its determinateness, its most general
expression, but this general expression is now individualized
by supervening factors, and from the different
combinations the derived passions proceed as necessary
products.

Their multiplicity is grounded partly in the nature of imagination
and memory, which makes it possible to preserve
the impression of the affections after the external objects
are no longer present, and thereby to combine the passions
and to set one against another; --- partly in
the nature of the human body as composite, since it
can thus be affected in several ways at once; --- partly
in the interaction with the external objects, in the affection
by their affections, in the collisions with their passions;
--- partly in the representation of the objects' difference of degree,
and finally --- what makes the nuancing infinite
--- in the infinite diversity of individuals and in the infinite
multiplicity of external objects. But through all
the derived passions the fundamental passions show themselves, and in
reality the former are only these very ones, named with different
names according to their different relations and denominationes
extrinsecæ (Eth. III. Affect. defin. 48. explic.).

An affection that has taken place simultaneously with
another must, according to the laws of the association of ideas, repeat itself
when this other repeats itself (Eth. III, 14), and thus
any object can per accidens become the cause of joy,
sorrow and appetite (thus arise sympathia, antipathia, propensio,
aversio, desiderium; prop. 15), and of the passions derived from them
(III, 50). In consequence of the property of
% --- p. 116 ---
\opage{116}\setcounter{footnote}{0}imagination that the image of the affecting thing is preserved until
the thing's existence is excluded by another idea, the
represented non-present must act as something present until
other ideas make it uncertain or deny it (spes,
metus, securitas, desperatio; prop.
18)\footnote{Fear and hope are always bound together.
Eth. III, 50. schol.}.

In the relation to the external object it follows from the nature of the appetite
for self-preservation that, just as with regard
to ourselves we strive to affirm what gladdens and to deny
what pains, we must have the same striving with regard
to what we love, the opposite with regard to what
we hate (pr. 25 f.), and that this striving must express itself
both in representation and in action (pr. 28) insofar as it is not
hindered by the representation of the consequences (metus, timor, verecundia,
consternatio; prop. 39). The representation of the
joy and pain of the beloved and of the hated must call forth in us
joy and pain in opposite
relation\footnote{Yet the joy over the annihilation of the hated
is mixed with sorrow, inasmuch as the hated is
like us. Eth. III. 23. schol. 47.}, and in
degree determined by the degree of love and hatred. The cause
of the beloved object's joy and of the hated one's
pain must become the object of our love, the cause
of the beloved's pain and of the hated one's joy, the object of our
hatred (commiseratio, favor, indignatio, invidia; prop. 19---24).
Love must include an appetite for
love in return\footnote{The proof of this proposition is
characteristic of Spinoza's conception. It is conducted thus:
We strive, as far as is in our power, to represent to ourselves the
object we love in preference to all other things. If it is like
us, then, since the affection of what is like us, through
representation (according to the laws of the association of ideas, cf. pr. 27),
calls forth a similar affection in us, we shall strive to
call forth joy in it rather than in others, or, in other
words, we shall strive that the beloved object feel
joy accompanied by the representation of us, i.e. that it love us
(Eth. III, 33). --- When the individual is conceived as individual, it must
be conceived as egoistic; but egoism in Spinoza is
vanishing, just like finitude (cf. p. 70 f.).}, and
the joy this calls forth is determined by the degree of
love. If the one we love prefers another,
% --- p. 117 ---
\opage{117}\setcounter{footnote}{0}love must turn into hatred (since the appetite to represent the love in return
is hindered), and the one preferred to us must become
the object of jealousy (zelotypia). --- The intensity of joy and sorrow,
of love and hatred, must determine the intensity of the appetite
arising from them. --- Hatred arising from love
must become more intense than if love had not
gone before (since an all the greater appetite is hindered); the
same must hold of love arising from hatred. The representation
of unmerited hatred must call forth hatred toward the
hater;\footnote{For likeness calls forth the same
affect, and the fact that the hater is the only cause
restricts the hatred to him.} the representation of undeserved love must call forth
love. --- Hatred must be strengthened by reciprocal hatred,
but brought to waver by the representation of love
in the one hated (ira, vindicta, gratia, gratitudo). Love
and hatred toward what we represent to ourselves as free must
become greater than toward what we represent to ourselves as necessary,
since the free is represented by itself alone, the necessary together
with something else, its
cause (prop. 34---44.
49)\footnote{For the same reason the object which has
nothing peculiar to itself, but only what is common to several, will affect us only
in a lesser degree (admiratio, consternatio, horror,
devotio, irrisio). Admiratio is not an affect, but a
reinforcing motive in the other affects.}.

Through \emph{likeness} to what we hate or love
an object can per accidens become an object of hatred or
love; but since likeness at the same time presupposes difference,
the same thing can become the object of a double
passion (animi fluctuatio, which can also arise from the
different influence of the same object on the composite human being;
prop. 16 f. 46.). --- An affect
in the object like us must, when we do not already
stand in a relation of affect to the object, through the representation
of the affect call forth a similar one (for imagination
includes the nature both of our body and of the external body,
and insofar as they agree, the affects of the one are transferred to
the other; --- commiseratio, æmulatio), and the appetite which
% --- p. 118 ---
\opage{118}\setcounter{footnote}{0}our own affections call forth is called forth by those of what is
like us. We must therefore strive to do what
others regard with joy, and to avoid what others
regard with displeasure (benevolentia, misericordia, ambitio,
humanitas, laus, vituperium, gloria, pudor, acquiescentia in
se ipso, poenitentia, superbia); we must be confirmed in our
love and our hatred by the representation of agreement
with those like us, while disagreement
with them must call forth wavering between the opposite
affects\footnote{If the object of love can belong only to
one, egoism must come forward (see Note 3, p. 116), and
we must strive to bring it about that the one whom we represent to ourselves
as rejoicing in it does not come into possession of it; and
thus the same peculiarity of
human nature which was the reason why men
were compassionate becomes the reason why they are envious and
ambitious (prop. 32. cf. 55. schol.). --- The same
appetite of self-assertion, which lies so deep in human
nature, also brings it about that humility and self-contempt (humilitas
and abjectio) are so rare as passions, and that they are as a
rule a cover for ambition and envy
(ambitio and invidia) (Aff. def. 29. expl.).} (prop. 27. 29. 30. 31. 45).

Finally, in the concept of the human body and of the affects
lies what gives the possibility of infinitely many
nuances and combinations: that the same object can
affect different men differently and the same
man differently at different times (pr. 51), that
of each affect there are as many kinds as there are objects
that affect us (since every affection includes the
nature of the external
object)\footnote{Eth. III, 56. dem.: natura uniuscuiusque
passionis ita necessario debet explicari, ut obiecti, a quo
afficimur, natura exprimatur.} and that every individual's affects
are different according to the difference of its essence, which expresses
itself in the affects. (The passions of animals are different
from the corresponding passions in men in the same
degree as the essence of animals and of men is different:
fertur quidem equus et homo libidine procreandi, at ille
libidine equina, hic autem humana).

% --- p. 119 ---
\opage{119}\setcounter{footnote}{0}In what has been developed, since man is conceived as a link in
the infinite series, the affects have been conceived as
passiones\footnote{In the conception of the passions as passiones
one proceeds from the body alone (Eth. III, App., IV, 7); in
speaking of them as actiones, from the mind; cf. tract. polit. II,
11: the strength of the mind, not of the body, determines
human power (potentia).}, joy too; for although it might seem that this, % print: passione (misprint for passiones)
as a transition to a greater perfection, would have to be
an actio, there is nevertheless, through the quantitative difference in
the relation between the adequate and the
inadequate\footnote{Remember that: ideas inadæquatas habere
is $=$ imaginari $=$ pati. cf. Eth. IV, 59 schol., V, 1. 11.} ideas in
man's mind, a possibility of thinking \emph{an approximation}
to the adequate which, so long as this has not yet
been reached, must be called a passio (cf. Eth. III, 57. dem.
58. IV, 39. dem.).

The general determination of the affect, as a passive
state in the \emph{mind} (animi pathema), can now be given thus:
it is
confusa idea, qua mens maiorem vel minorem
sui corporis vel alicuius eius partis existendi vim quam
antea\footnote{By these words there is not expressed a
comparison with an earlier state; but ``that the idea
which constitutes the essence (formam) of the affect affirms something of
the body which expresses a higher or lesser reality than
the earlier.'' Through this affirming the mind passes over to a
higher or lower perfection.} affirmat, et qua data ipsa mens ad hoc potius quam
ad illud cogitandum determinatur (Affect. gener. def.).

If one considers man, from the point of view that has been taken,
in his relation to nature, and compares
his power with the power of the external causes, then it will
appear that the state dominated by the affects is to be conceived
as a state of bondage, in which man's essence cannot
assert itself. It follows, as an axiom, from the thought of
the finite as individual, that there can be no individual
thing without there being another individual thing which is
stronger and by which the former can be destroyed (Eth. IV.
axiom.). The force with which man strives to
% --- p. 120 ---
\opage{120}\setcounter{footnote}{0}preserve his existence must therefore, as finite, be infinitely
surpassed by the power of the external causes, so that it becomes
impossible that he should not be subject to changes.
Since the power and duration of an affect must be determined by the
power of the external cause compared with ours, it can
have an overwhelming power over our actiones, and man
becomes constantly dependent on the passions and
unable to defend himself against their power; he must follow and
obey the general order of nature and accommodate himself to it
as far as the nature of things requires. With mechanical
necessity the affects struggle against each other in man;
the stronger necessarily overcomes the weaker, and only
an affect can remove an affect, since the bodily affection
which underlies it must persist until it is removed
by a stronger and opposite bodily affection.

The relative strength of the affects is calculated with mathematical
certainty according to the determinations developed earlier. The
present must act more strongly on us than the past
and the future, and of these the nearer more strongly than
the more remote. The necessary must have preponderance over the
contingent, the past over the non-present
contingent; the appetite that arises from joy must, other things
being equal, become more powerful than the appetite that arises from sorrow.
The knowledge of good and evil\footnote{For by good and evil is understood only what is useful or
harmful to us, i.e. what increases or diminishes our power of action,
thus that which calls forth the affect of joy and of sorrow, from which
the idea of the affect is distinguished only in thought (Eth. IV, 8 dem.)}, which is nothing other
than the very affects of joy and sorrow as conscious, cannot
overcome the affects by being true; for nothing
positive in them is removed by the presence of the true\footnote{This is a consequence both of the equal relation of all ideas to God as
adequate, and of the relation of imagination to the human body,
whose affection remains even after the error is removed.}; but
only insofar as it is considered as an affect; and as such
it can, since the appetite that arises from it must be comprehended
% --- p. 121 ---
\opage{121}\setcounter{footnote}{0}through our essence alone and be determined with regard to strength
by our power alone (Eth. IV, dem.), be hindered and
destroyed by many other appetites that arise from affects called forth
from without, and that especially when their object
is something present, while its object is something future or
contingent (Eth. IV, 14--17).

From the state of bondage there is therefore no way out so
long as man is conceived abstractly as individual. Just as
the consequences of this conception in the theory of knowledge
led to the impossibility of the mind's adequately conceiving
itself, so that, as a merely \emph{confused} idea, it came into
conflict with its own concept as the idea of the body, so
here too, consistently thought, the human
being, which always remains suffering, comes into conflict with its own concept
as active. Here too the task must be to remove this
contradiction by conceiving the human being in its
truth in the infinite, thus as free; but since in the sphere of the
infinite the relations of time have no significance, and
a \emph{transition} from the finite to the infinite becomes
just as impossible as a transition from the latter to the former, the
task cannot be to see man \emph{successively liberating
himself} from his state of bondage, but to comprehend him in his
\emph{eternally and necessarily existing freedom}, by which the
bondage proves to be a semblance. Thus Spinoza
conceives the matter in the last book of the Ethics, where he treats of
human freedom; but as a preparation for this
conception he gives developments which, while they still let
the common conception of man held by representation stand
as point of departure, and therefore seem to teach the \emph{successive}
liberation and to set this as a \emph{task} for an
action, on closer consideration show a necessary
development of concepts and a leading back of the finite to its
eternally existing essence. Yet it cannot be denied that
at this point some uncertainty enters in
Spinoza. While on the one hand (see below)
he conceives the adequate knowledge of God, in which freedom
is given, as \emph{belonging to man's}
% --- p. 122 ---
\opage{122}\setcounter{footnote}{0}essence, thus as eternally existing; conceives human
nature as a unity (res particularis æterna), in which
men represented as individuals are so united
that their drive of self-preservation, through this common element, is a
drive to the preservation of the nature of other men; and lets
all limitation, even that which seems to take on positive
forms, like crime, vanish as something merely negative,
yet on the other hand the separation of individuality
asserts itself: what belonged to the human essence
is posited as something rare (Eth. IV. App. c. 13. V, 42.
schol.), and what was conceived as vanishing gets
power to hold thought fast, so that thought, on the vanishing
foundation, can erect a kind of pedagogy and a
politics; and even where the human mind is conceived in
its freedom and in its being in God, it still preserves its
\emph{determinate} existence in the divine intellectus (Eth.
V. 40. schol.). The same urge which in the domain of the purely
metaphysical forced finitude forward makes itself felt
here; but here it has gained a greater intensity through
the unconsciously working sympathy for the humane.

As a preparation for the true conception of human
freedom Spinoza first develops \emph{what reason
prescribes with regard to the affects}, and,
since a standard for judging them is thereby given,
he determines them as \emph{good} or \emph{evil} according to their
agreement with human nature; and since
through these developments the relation between
the \emph{free man} and the \emph{slave} comes forward, the way is prepared for decisively
letting \emph{liberation} show itself as an \emph{existing} freedom.

The determination of what reason demands of us with
regard to the affects becomes, since ``\emph{the demand of reason}''
and ``nature's appetite'' are the same thing named with two
names, equivalent to what follows from the concept of the human
being, idea hominis, which is not an ideal of the
imagination\footnote{We form such a one, e.g., when, from confused images of the
outward shape of many individual men, we form a universal representation
of man, and now refer the individual men to
this ideal.}, but the expression for the human being
% --- p. 123 ---
\opage{123}\setcounter{footnote}{0}as res particularis æterna (cf. de intell. emend. p. 388 f.).
This is set up as \emph{exemplar humanæ naturæ, quod
intueamur} (Eth. IV. præf.), and since man represented as
individual is referred to it as to his truth,
the otherwise merely subjective and relative\footnote{As these concepts are usually conceived, they designate nothing
real in things considered in and for themselves; but are only modes
of representation which we form according to the accidental relation of things to us (in
such a way that the same thing can be called good, evil or indifferent according to
our individual difference), or when we compare things
with one another or with a universal representation, and then conceive the fact
that the individual thing expresses less reality as a defect in its
essence. But since the same law of necessity governs everything,
and since everything that is and happens depends on the necessity of the infinite
divine nature, everything becomes perfect according to
the degree of its essence (for perfection is equal to reality, being); there
remains no room for anything imperfect; nothing belongs to the nature
of any thing other than what, according to the concatenation of things,
follows from the determinateness of the efficient cause, and what follows from
it happens necessarily. Nothing can happen against God's will, i.e.
against the law of the divine nature, not even what we call
crimes and the like; for that expresses not a determination of essence,
something positive, but only a negation. In Nero's matricide,
thus, the positive was not the crime, but the outward action and
the intention, which were the same as Orestes' action and intention. What
constituted the crime was that by this action he
showed himself ungrateful, disobedient and without compassion. But none of
this expresses essence, and therefore God is not the cause of it, although
he is the cause of Nero's action and intention. The negative in the
action vanishes when it is led back to God, just as the negative
in the inadequate ideas vanishes by this leading back, and
just as the limitation and negation in the finite modi are cancelled
when they, under their apparent individuality and infinite
change, are seen to be, with limitation cancelled, in the unchangeable,
infinite One. In inadequate, imperfect knowledge
therefore lies the origin of all those modes of representation, and they lose
their significance with it (cf. Eth. I, 36. schol. IV. præf. 64, 73
schol. De intell. emend. pag. 360. Cog. metaph. I, c. 6. II, c. 8.
Princ. phil. Cart. I, 4 ax. 4. Ep. 32. 34. 36. tr. polit. II, 8. tract.
theol. polit. pag. 177.)} designations:
% --- p. 124 ---
\opage{124}\setcounter{footnote}{0}\emph{good} and \emph{evil}, \emph{perfect} and \emph{imperfect}, find
an application, justified from the standpoint of the individual, to
that which approaches or recedes from this model, in
such a way that \emph{good} becomes that which we know with certainty to be
useful to us, \emph{evil} the opposite (Eth. IV. def. 1 and 2)\footnote{\emph{Good}
thus becomes that which disposes the human body
to be affected and to act in many ways, since the mind's power to
perceive corresponds to the body's power to act. Likewise
that is good which preserves the relation of motion and rest between the body's
parts, since a change in this relation destroys the human
body and brings it to pass over into another form. (This
is \emph{death}, but the concept of death receives with Spinoza a greater
extent than is commonly given it, since it designates not merely
the transformation into a corpse, but every change of such an
extent that nature shows itself as another, Eth. IV, 39. schol.). --- Also everything must be useful that furthers the union of men,
while the opposite is harmful (Eth. IV, 38--40).}.

Reason cannot demand anything that conflicts with
nature; since, therefore, everyone according to the necessary laws of his nature
desires the useful and flees the harmful, reason
must demand that everyone love himself, that he
seek what is truly useful to him, that he desire
what leads to a greater perfection, and, in
general, that he strive, as far as his power reaches,
to preserve his being. This is the demand of reason because
it is the necessity of nature. \emph{This striving is virtue
itself.} Virtue is equal to man's very essence
conceived as active\footnote{Eth. IV. defin. 8: Per virtutem et potentiam idem intelligo; hoc
est: virtus, quatenus ad hominem refertur, est ipsa hominis essentia
seu natura, quatenus potestatem habet quædam efficiendi, quæ
per solas ipsius naturæ leges possunt intelligi.}, thus to humana potentia. Since it
is one with man's very essence, it is desired for its
own sake, and the fundamental appetite, the appetite for self-preservation,
becomes its first and only foundation, since being is
the condition for being happy, acting well and living
well\footnote{Eth. IV, 20. schol.: nemo igitur, nisi a causis externis victus, suum
utile appetere sive suum esse conservare negligit.}. --- Since virtue is posited as one with potentia,
% --- p. 125 ---
\opage{125}\setcounter{footnote}{0}only he can act virtuously who has adequate ideas; for only
insofar as we have these are we acting (``nos eatenus tantummodo
agimus quatenus intelligimus.'' Eth. IV, 24. dem.).
The complete expression for virtuous action thus
becomes: ex ductu rationis agere, vivere, suum esse conservare,
ex fundamento proprium utile quærendi (Eth. IV. 24). For ratio
expresses the mind insofar as it comprehends clearly and
distinctly\footnote{Since reason conceives things under the form of eternity and necessity,
it does not conceive them in relation to time, which is only
a determination of imagination, and must therefore prefer the
greater good to the lesser, the lesser evil to the greater, without
regard to whether they are present or future (Eth. IV.
62. 65 f.).}. And since it can only affirm its own essence,
that toward which it bids us strive can only be to
comprehend (intelligere), and only what leads to this, or
hinders it, can we know with certainty to be good or
evil. \emph{The good is therefore knowledge}\footnote{At this point the motives and peculiar character of Spinozism
come forward, to be decisively asserted at the conclusion of the Ethics. --- Spinoza does not develop the relation of the bodily to the proposition
that knowledge is the proper aim of self-preservation.}\emph{, the
highest good the highest knowledge: the
knowledge of God, and the mind's highest virtue is to
know God} (Eth. IV. 28)\footnote{Hominis, qui ratione ducitur, finis ultimus, \emph{hoc est, summa cupiditas},
qua reliquas omnes moderari studet, est illa, qua fertur
ad se, resque omnes, quæ sub ipsius intelligentiam cadere possunt,
adæquate concipiendum (Eth. IV. App. c. 4). --- Nulla igitur vita
rationalis est sine intelligentia, et res eatenus tantum bonæ sunt,
quatenus hominem juvant ut mentis vita fruatur, quæ intelligentia
definitur. Quæ autem contra impediunt, quominus homo rationem
perficere et rationali vita frui possit, eas solummodo malas esse
dicimus (ibid. c. 5). --- Beatitudo nihil aliud est quam ipsa animi
acquiescentia, quæ ex Dei intuitiva cognitione oritur. --- Intellectum
perficere nihil --- aliud est, quam Deum, Deique attributa, et actiones,
quæ ex ipsius naturæ necessitate consequuntur, intelligere. --- In hoc uno (i.e. intellectum perficere) summa hominis felicitas
seu beatitudo consistit (ibid. c. 4).}.

% --- p. 126 ---
\opage{126}\setcounter{footnote}{0}Just as necessarily as what reason prescribed \emph{with
regard to ourselves} followed from the concept of our essence, what
it prescribes with regard to \emph{the relation to others} follows
from the common concepts. Since things are individual
and different only through the negation of limitation, and this is
the source of impotentia, which is the opposite of power, i.e. of virtue,
then, the more limitation is removed, the more must impotence
be removed and virtue approached, since our power of action is increased by
connection with what is of like kind. That whose nature was absolutely
different from ours could be neither useful nor
harmful; that which agrees with our nature must
be useful, i.e. good, since there comes to be a unity in the common
nature, so that the appetite of self-preservation of the one thing immediately
benefits the other\footnote{This grounding is of significance for Spinoza's conception of \emph{the
common} (commune). The individual thing is in its self-preservation
immediately one with the other, since the common is their positive
unity (Eth. IV, 31 dem. cf. 35 dem.).}. Usefulness increases in the same
degree as agreement, in such a way that if two individuals
which were perfectly of the same nature were united, they would
form an individual that had twice as great a power
as the single one. For man, therefore, nothing is more useful
than man, and men can wish for nothing
better for maintaining their being than that they should all agree
so perfectly in everything that the souls and
bodies of all formed as it were one soul and one body, and that
they all in union, as far as they were able, strove to
preserve their being and to further the common useful. But
as dominated by affects, men cannot be said
to agree in nature; for this is a positive
determination, the affect (as impotentia)\footnote{Eth. IV. 37. schol. 1: impotentia in hoc solo consistit, quod homo a
rebus, quæ extra ipsum sunt, duci se patiatur, et ab iis ad ea agendum
determinetur, quæ rerum externarum communis constitutio,
non autem ea, quæ ipsa ipsius natura, in se sola considerata,
postulat.} a negative one. Even
the individual man in that state does not agree
% --- p. 127 ---
\opage{127}\setcounter{footnote}{0}with himself, still less the different men, who through the
passions can become opposed to each other. Only as led
by reason do they agree of necessity; for
only rational actions are comprehended through human
nature alone as proximate cause, and since \emph{reason's}
knowledge of good and evil is necessarily true,
the actions that proceed from reason must necessarily
contain what is a good for \emph{human
nature}, and thus for every man (cf. p. 122 f., Note 1
p. 126). When the individual is led by reason, his striving
for what is useful to him becomes immediately useful to the others,
and the more useful the stronger it is. For in reason, which
expresses the common human nature, separation is
removed, and the drive of self-preservation is no longer
egoistic. The highest good of the rational man, i.e. of the virtuous man:
the knowledge of God, the most perfect being, without which
nothing can be \emph{or be comprehended}, must be common to all
and capable of being possessed by all, \emph{``for it belongs to the
essence of the human mind} to have an adequate
knowledge of God's eternal and infinite essence'' (Eth. IV. 36.
schol.). In the appetite with which the rational man desires
this good for himself there lies also, immediately, the appetite
for it for the others, and that in so much higher a degree the greater
the knowledge of God is\footnote{A \emph{gradation} in the knowledge of God has significance only for the
common representation, not for Spinoza's peculiar conception;
cf. chapter 8.}. At this point there thus appears
the significance of the \emph{social relation} among men
(see the next chapter), and the true significance of \emph{religion}:
``whatever we desire and do, whatever we are the cause of
insofar as we have the idea of God, or insofar as we know
God, I reckon to religion'' (Eth. IV. 37. schol. 1)\footnote{cf. Eth. IV, 73. schol.: vera vita et religio. On religion in the
usual sense of the word cf. chapter 7.}.
It appears likewise here \emph{that there are also affects
which have their origin in reason,}
% --- p. 128 ---
\opage{128}\setcounter{footnote}{0}\emph{and are active.} These must be referred to two of the fundamental
affects: \emph{appetite} and \emph{joy}; for \emph{sorrow} presupposes a
suffering. The appetite arisen from reason Spinoza calls
\emph{fortitudo} (since it expresses man's power) and divides
it into: \emph{animositas}, i.e. the appetite by which one
strives to preserve one's being solely according to the prescription
of reason; --- and \emph{generositas}, i.e. the appetite by which
one strives, solely according to the prescription of reason, to
bind other men to oneself in a relation of friendship
(Eth. III. 59. schol.)\footnote{Modifications of \emph{animositas} are: temperantia, sobrietas, animi
præsentia etc., of \emph{generositas}: modestia, clementia etc.}. Generositas thus becomes not
different from \emph{pietas}, i.e. the appetite to do good which
is called forth by our living according to the bidding of reason\footnote{This will is dependent on idea Dei, which is posited as equal to spiritus
Christi. Eth. IV. 68. schol.},
or from \emph{honestas}, which in Eth. IV. 37. schol. 1. is defined
quite in the same way as it\footnote{\emph{honestum} is the name for what men led by reason
praise, \emph{turpe} for what hinders the relation of friendship (Eth. IV. 37.
schol. 1). \emph{Probi} is the name for those who have a clear idea of God, by
which all their actions and thoughts are determined; \emph{improbi} for those
who \emph{do not possess the idea of God} (qui Dei ideam haud possident)
but only the ideas of earthly things, by which their actions
and thoughts are determined (Ep. 36).}.

If the determinations \emph{good} and \emph{evil} are applied to the
individual affects, the consequence must be that those of them
are called good which designate a power, and those evil which
designate an impotence\footnote{``nostræ actiones, hoc est, cupiditates illæ, quæ hominis potentia
seu ratione definiuntur, semper bonæ sunt, reliquæ autem tam bonæ
quam malæ possunt esse'' (Eth. IV. App. Cap. 3).}. Virtue was indeed determined as power, as
self-assertion (conceived in the right sense), not as self-annihilation;
for nothing positive is removed by the unity with
substance. The affects that express self-assertion are
good, since they express a transition to a greater perfection;
those that express a hindering of our essence's
% --- p. 129 ---
\opage{129}\setcounter{footnote}{0}fundamental striving, which can take place only from without\footnote{Eth. IV. App. c. 6: quia omnia illa, quorum homo efficiens est
causa, necessario bona sunt, nihil ergo mali homini evenire potest
nisi a causis externis; nempe quatenus pars est totius naturæ,
cuius legibus humana natura obtemperare, et cui infinitis modis
pæne sese accommodare cogitur.}, are evil.
Good, therefore, are joy, hilaritas, and the appetite and the
self-feeling that arise from reason\footnote{Eth. IV, 45. schol. 2: quo maiori lætitia afficimur, eo ad maiorem
perfectionem transimus, hoc est, eo nos magis de natura divina
participare necesse est.}; evil are not only
affects such as sorrow, melancholy, hatred, despair, pride
and the contempt for oneself (abiectio) that is related to pride and to envy,
but also such as, even though in the man not led
by reason they can become useful as a counterweight against opposite and more harmful ones,
nevertheless include a sorrow and are therefore in and for themselves harmful,
thus: fear, hope, humility, repentance, pity,
shame. Such affects as refer themselves to a partial
affection, which can attain excess, like the feeling of pleasure and the feeling of
pain (titillatio and dolor), can per accidens become both
good and evil, yet the former is rather good, the latter evil. \emph{Appetites}
are on the whole good or evil according as they
arise from affects that are good or evil. But all those
that arise from affects that are passiones are blind, and
would be of no use to men if they
could easily be led to live according to the bidding of reason, for
\emph{by reason we can be determined to all the
actions to which we are determined by an affect
that is a suffering} (Eth. IV, 59)\footnote{For in every passive state our power of action is less than that
which follows from the necessity of our nature, considered in itself alone
(Eth. IV. 59). Joy too is a passio insofar as our power of action
is not increased to the point at which it comprehends itself and its actions
adequately (Eth. IV. 59. dem.). The transition to a \emph{greater
perfection} is thus, in the domain of the passive states, only
the transition to a \emph{lesser imperfection.} The \emph{positive}
expression has significance only so long as the representation of an \emph{existing}
finite is held fast; when this representation dissolves, first
the positive expression must be exchanged for the negative, and then the \emph{transition} as a whole must show itself as something negative, a non-being (cf.
chapter 8.).}.

% --- p. 130 ---
\opage{130}\setcounter{footnote}{0}Through these preparatory developments the difference
now shows itself between the \emph{free man} and the \emph{slave} (the wise man and
the ignorant), and of this Spinoza gives, before he passes
over to the decisive conclusion of his Ethics, a portrayal (Eth.
IV, 67 ff.), in which the sphere of finitude and relativity
is not yet left in the way it is left by that
conclusion. While the slave is led by affect, is determined
by external causes, and therefore, whether he will or not,
must do what he does not understand, the free man is led only
by his own essence, reason, clear and distinct knowledge,
and therefore does only what he knows to be the
essential, and for that reason desires (Eth. IV. 66. schol.).
Since he is determined only by reason, which is his essence,
and since reason is the mind \emph{as active}, no
other affects can be found in him than the active affects of joy and of
rational appetite\footnote{cf. Eth. IV. App. c. 25: in communibus --- colloquiis cavebit hominum
vitia referre, et de humana impotentia non nisi parce loqui
curabit: at largiter de humana virtute seu potentia, et qua via possit
perfici, ut sic homines, non ex metu aut aversione, sed solo lætitiæ
affectu moti, ex rationis præscripto, quantum in se est, conentur
vivere.}. He is thus not led by
fear to the good, for fear is a sorrow, but his appetite
seeks the good \emph{directly}, since it arises from the knowledge
of it, which is one with the affect of joy as
conscious (see above, p. 120). Likewise it flees evil only
indirectly, since the knowledge of evil, as the conscious
tristitia, is an inadequate knowledge\footnote{The knowledge of evil is an inadequate knowledge because it is
equal to the conscious tristitia. For sorrow is a transition to
a lesser perfection, which cannot be understood through man's
essence alone and therefore is a passio dependent on inadequate
ideas, so that the knowledge of it can only be inadequate. If
we had only adequate ideas, i.e. if we were born free, we would not form
the concept evil, and thus not its opposite either: the concept
good (Eth. IV. 64. 68).}.
He thinks least
% --- p. 131 ---
\opage{131}\setcounter{footnote}{0}of all of death; his wisdom is a meditation on life
(vitæ meditatio). He does not let himself be influenced by the affects of the
unfree, but sets love\footnote{In Eth. IV. App. c. XIX--XX animi libertas is posited as the ground of
love. --- Voluntas constans et in omnibus determinata is
conceived as a property of intellectus. Ep. 60.} and magnanimity against
hatred; he preserves independence of the ignorant; he
alone is capable of gratitude; he always acts with good
faith; he is freer in the state, where he lives according to the common
decision, than in solitude, where he obeys only
himself; he hates no one, envies no one, despises no one,
is angry with no one; what he desires for himself he desires
also for the others; he considers above all that everything
follows from the necessity of the divine nature\footnote{Eth. IV. App. c. 32: Quod (i.e. that we, as part of nature, must
follow its necessary order) si clare et distincte intelligamus, pars
illa nostri, quæ intelligentia definitur, hoc est, pars melior nostri,
in eo plane acquiescet, et in ea acquiescentia perseverare conabitur.
Nam quatenus intelligimus, nihil appetere, nisi quod necessarium
est, nec absolute, nisi in veris acquiescere possumus; adeoque,
quatenus hoc recte intelligimus, eatenus conatus melioris partis nostri
cum ordine totius naturæ convenit.}, and that
what he therefore represents to himself as troublesome and evil, and
what seems impious, unjust, dreadful and shameful,
only seems so because he regards things confusedly and
mutilatedly, and therefore he strives first and foremost to
comprehend things as they are in themselves, and to remove the hindrances
to true knowledge, such as hatred, anger, envy,
mockery, arrogance and the like, and, as far as it lies with him, \emph{to
act well and to be glad} (Eth. IV, 73. schol.); for
thereby he expresses his essence, reason.

Human freedom is thus seen in its opposition
to man's state of bondage within the domain of
finitude: it now only remains to see man
as free in his eternal and necessary relation to God, the
% --- p. 132 ---
\opage{132} only causa libera, and to let this relation cast a light back on what has been developed with regard to the common representation. But at this point it is the fitting place to insert Spinoza's doctrine of the \emph{state} and of \emph{religion}, which, taking this representation as its point of departure, supplements what has just been developed, and at the same time indirectly prepares the way for the more complete understanding of the doctrine of \emph{the eternally existing human freedom} which Spinoza develops in the last book of the Ethics.

\bigskip
\begin{center}\rule{2.5cm}{0.4pt}\end{center}
\bigskip

\begin{center}
{\large\textbf{Chapter Seven.}}\\[4pt]
\textbf{Natural Right. Religion.}\\[6pt]
\rule{2cm}{0.4pt}
\end{center}
\addcontentsline{toc}{chapter}{Chapter Seven. Natural Right. Religion}
\markboth{Natural Right. Religion}{Natural Right. Religion}
\bigskip
Spinoza's doctrine of the \emph{social relation} takes, like the Ethics proper, its starting point from man, represented as a single modus, as a dependent link in the order of nature, in the infinite chain of causes and effects. Conceived as such, man (cf. the previous chapter) is governed by the universal law of external necessity; he is constantly determined by passive affects, which act according to the mathematically calculable proportion of their strengths and, by the power they receive from external causes (cf. Eth. IV, 3. 6), are able to overwhelm the very knowledge of good and evil, which acts against them only in so far as it itself takes on the form of an affect. Man is here like one of the fleeting figures that are formed in infinite extension by the laws of motion, only to be reshaped again and at last to vanish into its unity, as the limitation, being merely negative, vanishes. From that standpoint, which is that of the imagination, we saw that liberation from the affects, from the state of bondage, stands as a goal that can be reached only approximately, since the finite as finite cannot become
% --- p. 133 ---
\opage{133}congruent with the infinite, and, within the \emph{relative} difference between the slave and the free man, the ignorant and the wise, the free and the wise make up only a small number. Since, therefore, human nature through the drive of self-assertion strives after the unity that increases its power, that which is to be able to act upon all men, and to be capable of forming the foundation of the social relation which, as experience shows, is found everywhere, not only among civilized men but among barbarians, must be, not the rational alone, but the affect. The concept of the affect must therefore become the starting point of the doctrine of the state; in its mechanical laws the laws of the state's coming into being and continuance must be sought, and since man as determined by it is single and isolated, the individual drive of self-preservation, the \emph{egoistic} appetite, must form the foundation, and \emph{laws of nature}, not \emph{ethical} laws, must be what is determinative. But just as in the doctrine of knowledge and of the affects we saw that the conception of man as single cancelled itself, so we shall see it here too. Through the affect as the phenomenal we shall see man's true essence, reason, in which all men are in unity, act as the deeper \emph{ground} of the association and determine its essence and \emph{end}. The current here, as everywhere in Spinoza, runs back to \emph{the unity in the infinite}, and even if this unity is not reached in this domain, by the nature of the matter, it is nevertheless intimated, and the lines drawn converge so as to meet in one point at the conclusion of the Ethics. The development will therefore point toward this point; but from the necessary inconsistencies in Spinoza's doctrine of the finite developed earlier, the consequence will be that the finite frequently holds thought fast as though it had the validity of reality, and that that point is seen only as in the distance.

Since the doctrine of society proceeds from the single man, governed by the universal laws of nature, his essence must first be determined in this isolation, \emph{man's state of nature}. The determination is drawn from the universal metaphysical relation between God and single
% --- p. 134 ---
\opage{134}\setcounter{footnote}{0}things. Since these have existence and continuance only through God's eternal power, \emph{the power of natural things, by which they are and act, is God's} (nature's) \emph{very power}. But since God has a right to everything, and his right is equal to his power, regarded as absolutely free, the right of the natural thing to exist and act is determined by its power, which expresses God's absolutely free power in a determinate way. \emph{Jus naturæ}\footnote{``Nature has a right to everything''; tract.-theol.-polit. pag. 175. As synonyms are set down in the same place: jus naturæ, institutum naturæ, regulæ naturæ uniuscuiusque individui.} is thus the power of nature (naturalis potentia), i.e.\ its laws, according to which everything happens and which are not disturbed by an imagined human freedom. What every single individual does according to the laws of its nature, it does with the most perfect right of nature (\emph{summo naturæ jure}). But as part of nature man is determined not only by reason but also by the passions; in both conditions the appetite for self-preservation is at work, and in both it is an expression of vis naturalis; for even if the passions, referred to men alone, are passiones, they are nevertheless determined by the laws of nature, effects of the power of nature, and in the same proportion as this power is expressed by the nature of this or that man, man shares in the right of nature, which is not separated from power. It therefore makes no difference whether man is wise or ignorant; whatever he desires and does in either capacity, he desires and does with the full right of nature (tract. polit. II, 8). In this state, then, egoism becomes the sole determinant; each is bound only by his own appetite\footnote{Eth. IV. 37. schol. 2.: Existit unusquisque summo naturæ jure, et consequenter summo naturæ jure unusquisque ea agit, quæ ex suæ naturæ necessitate sequuntur, atque adeo summo naturæ jure unusquisque judicat, quid bonum, quid malum sit, suæque utilitati ex suo ingenio consulit, seseque vindicat, et id, quod amat, conservare, et id, quod odio habet, destruere conatur. cf. Eth. IV. App. c. 8.}. A promise, e.g., to do something which one could by one's right refrain from doing, thus has
% --- p. 135 ---
\opage{135}\setcounter{footnote}{0}validity here only so long as the will of the one promising is unchanged. He who has the power to break the promise has in reality not given up his right, but has only ``given his word''. When, therefore, he sees or believes he sees---and by the right of nature he is judge in his own cause---that by fulfilling the promise he will have more harm than benefit, he will break the promise, and do so by the right of nature (II, 12).

This jus naturæ, ``under which all men are born, and most live'', therefore forbids nothing except what no one desires and no one is able to do, neither strife nor hatred, nor anger, nor deceit, nor anything whatsoever to which appetite prompts. ``And this is not surprising; for nature is not comprised by (non continetur) the laws of human reason, which have regard only to the true benefit and preservation of men, but by infinitely many others, which have regard to the eternal order of the whole of nature, of which man is a small part, and by which order, which we know only incompletely, every individual is determined with necessity in its existence and action''\footnote{cf. tract. polit. II. 1--8; tract. theol.-pol. pag. 176 ff. --- The propositions last cited get their validity in that men are conceived as single beings among innumerable other single things, and in that the difference between ratio as ``universal knowledge'' and intuitive knowledge is held fast.}.

Since men are thus determined only egoistically by the affects, they are different from and opposed to one another\footnote{Were they led by reason alone, their appetite and right would not come into collision (Eth. IV. 37. schol. 2).}, and thus more dangerous to one another than any other being. \emph{Men are by nature enemies} (tr. polit. II, 14); for he is most my enemy whom I must most fear and guard myself against. In this relation of enmity each individual's appetite for self-preservation stands in opposition to that of all the others, his right stands against theirs; \emph{it therefore becomes impossible for the individual to assert his right,} since this right is determined by power, and the individual's power is
% --- p. 136 ---
\opage{136}\setcounter{footnote}{0}insufficient against that of all the others. Natural right therefore exists more in representation than in reality (opinione quam re); the individual's power is checked by fear, and with power right is diminished. Thus it is shown here too how the conception of the individual as individual is untrue and cancels itself. Only by union with the others, by the cancelling of isolated existence, do the assertion of right and a human, and in particular a rational, life become possible, since connection with what is like us gives greater power, and thus also greater right. \emph{The appetite for self-preservation therefore drives with necessity to the appetite for union}\footnote{Tract. polit. II, 15: concludimus, jus naturæ, quod humani generis proprium est, vix posse concipi, nisi ubi homines jura habeant communia, qui simul terras, quas habitare et colere possunt, sibi vindicare, seseque munire, vimque omnem repellere et ex communi omnium sententia vivere possunt.}. And this appetite issues with such strength from the fear of solitude that all men must be said to desire by nature the social relation (status civilis), and they can never wholly abolish it (tr. polit. VI, 1). In man, regarded as determined by the affects, it must now be an affect (such as fear, hope, vindictiveness), not reason, that appears as what prompts to union and secures it\footnote{In order that individual men may be secure and have confidence in one another, their affects must be governed by an affect, and this happens through the power of the state, which instils fear (Eth. IV, 37. schol. 2). --- On the whole, in the arrangement of the state and in regard to its preservation, the essential thing becomes to calculate the balance of the affects, and to arrange the state in such a way that the realization of what is useful to it, as regards both the rulers and the ruled, is secured, whatever their disposition may be, whether, that is, they do it voluntarily or compelled by the affects. Nothing that concerns the common weal ought to be left to the trustworthiness (fides) of any one individual; for then the continued existence of the state would be abandoned to chance, while the task is precisely to preserve it, since the dissolution of the state does not take place without the most violent upheavals; and with respect to the security of the state it makes no difference by what disposition men are moved to govern it well, provided only that it is governed well. ``For freedom or strength (fortitudo) of mind is a private virtue, but the virtue of the state is security'' (tract. polit. I, 6. VI, 3. X, 4. 9).}.
% --- p. 137 ---
\opage{137}\setcounter{footnote}{0}But the deeper ground of union lies in reason; for since the rational man is most sui juris (II, 11), the appetite to assert his right must at the same time become an appetite for the rational, and this was indeed (see the previous chapter) what is alike in all men, which as the unity of human nature unites them in the unity of the highest good. This deeper ground of union will then also become what determines its end, and what essentially and \emph{positively} holds it together (cf. what follows).

Now, when the individual, in order to increase his power and assert his right, enters into union with others, he gives up his existence as isolated, and exists from now on only as a member of society. The extent of his right is determined by its relation to the power of society; he now has only that jus in naturam which the common right (jus commune) grants him. He is obliged to do what is laid upon him by common decision, and can rightly be compelled to it. The common right, which is determined by the power of the multitude, is the power of the state (imperium)\footnote{Tract. polit. III. 2: \emph{imperii} seu summarum potestatum jus. --- In society there arises the concept of a persona civilis (a collective personality), as distinct from persona naturalis, tract. polit. X, 2.}, which he possesses absolutely who by common agreement has the government of the state. The condition in society is, in contrast to the state of nature (status naturalis), a \emph{status civilis}; the whole imperii corpus is \emph{civitas}\footnote{Civitas is: societas legibus et potestate sese conservandi firmata. Eth. IV. 37. schol. 2.}; its common affairs, which depend on him who holds dominion, are \emph{respublica}. A citizen (\emph{civis}) is he who by virtue of jus civile enjoys the advantages of the political community. He is called a subject (\emph{subditus}) in so far as he is bound to obey the laws and institutions of the state. An enemy (\emph{hostis}) is anyone who lives outside the state and is not an ally.

% --- p. 138 ---
\opage{138}\setcounter{footnote}{0}Only in this new relation do the determinations peccatum, justitia and the like acquire meaning. In the state of nature there is no offence (peccatum i.e.\ actio, quæ jure fieri nequit), or, if anyone offends, he offends against himself, not against others; for the law of nature forbids only what no one is able to do. Men led by the affects, too, follow the necessary order of nature, i.e.\ God's institutum, which is eternal and cannot be violated, and from this it follows: ``that the ignorant and mentally powerless man is, by the right of nature, no more obliged to order his life wisely than the sick man is to live with a healthy body.'' The concept of offence can first become intelligible in society, where it is determined by the common right of the whole state what is good and what is evil, and no one does anything rightly except what he does with the common consent\footnote{Peccatum thus becomes equal to inobedientia, meritum to obedientia (Eth. IV. 37. schol. 2).}. The following of this common will\footnote{This is lex civilis, which depends on men's discretion, whereas lex naturalis depends on the necessity of nature; cf. tract theol.-polit. pag. 43.} is \emph{obedience} (obsequium), which can be determined as the will to follow the commands of reason (which is true freedom) only in so far as a real unity of the state can be brought about only by rational determinations of right, and the determinations of right of the best state must be laid down according to the bidding of reason.

What holds with respect to offence and obedience holds also with respect to \emph{justice} and \emph{injustice}. In the state of nature there is nothing that can rightly be said to be the one man's and not the other's; but everything belongs to all in so far as they have the power to assert it for themselves. But in the state, where it is determined by the common right what belongs to the individual, he is called \emph{just} who has a constant will to give each his own, and \emph{unjust} he who strives to appropriate what is another's.

% --- p. 139 ---
\opage{139}\setcounter{footnote}{0}Since the general proposition that \emph{right is determined by power} holds of society just as well as of the individual, its preponderant power gives it the right to determine, over against the individual, how he shall live, which cannot be left to the individual's discretion without the social condition reverting again to the state of nature. For the same reason the power of society must itself be the sole interpreter of its laws. The individual is thus not sui juris, but civitatis juris; what the state determines to be just and good must be regarded as determined by everyone, and even if the subject finds the state's decisions unreasonable, he must obey them. \emph{The individual's natural right does not, however, properly cease in the social condition}\footnote{On this point Spinoza departs from \emph{Hobbes}. Ep. 50: Quod ad politicam attinet, discrimen inter me et Hobbesium --- --- in hoc consistit, quod ego naturale jus semper sartum tectum conservo, quodque supremo magistratui in qualibet urbe non plus in subditos juris quam juxta mensuram potestatis, qua subditum superat, competere statuo, quod in statu naturali semper locum habet. cf. Bruder's edition of Spinoza's works. Vol. III. pag. 213. Note.}; for in it, just as fully as in the state of nature, man acts according to the laws of his nature and provides for his own benefit. In both conditions he is led by fear and hope; but the essential difference is that in society all fear the same thing, and that the ground of security, as well as the manner of living, is one for all, which does not abolish the individual's faculty of judgement. For he who resolves to obey all the commands of the state, whether because he fears its power or because he loves security, provides according to his own mind (ex suo ingenio) for his security and his advantage. Nor does the transfer of right, the submission to another's judgement, conflict with the command of reason; for it was shown, indeed, that in the state of nature it was impossible to make one's right valid, whereas the political community and the observance of the common laws pave the way to it, so that the rational man, i.e.\ the free man, thus puts up with \emph{the} lesser evil, in a single point
% --- p. 140 ---
\opage{140}\setcounter{footnote}{0}of having to subordinate himself to the irrational, since it is far outweighed by the goods that follow from society; for reason commands us to choose the lesser evil when it leads to a greater good.

Moreover, with respect to what is apparently dangerous and contrary to reason in the transfer of right, it must not be overlooked that this transfer has its natural limits. For just as in the state of nature that man was most powerful (potens) and most sui juris who was led by reason, so that state will be the most powerful, and most sui juris, which is founded on reason and governed by it. The right of the state is determined, namely, by the power of the multitude that is led as if by one mind; but it is reason alone in which men agree, and unity will take place only when the state strives after what sound reason declares to be useful for all men\footnote{Concord founded on fear is founded only on something negative (Eth. IV. App. c. 16. tract. polit. X, 8) and is sine fide. Its positive ground is justice, equity and honesty.}. The power of the state is thus conditioned by its rationality, and its power, and thus also its right, is diminished the more, the more people are provoked to revolt by irrational commands from it and by their resistance give it occasion for fear. Herein, then, lies a safeguard for the individual, and moreover there always remains a domain that withdraws itself from the power of the state, namely everything that cannot be influenced by rewards or threats, such as the verdicts of the faculty of judgement, conviction\footnote{In this domain the individual comes essentially to relate himself merely to the divine. The relation to other individuals and to the state vanishes here. cf. tract. polit. III, 10: mens, quatenus ratione utitur, non summarum potestatum, sed sui juris est.}, love, hatred and the like; likewise that which is so much against human nature that man prefers any evil to it, e.g.\ to bear witness against oneself, to kill one's nearest, not to strive to avoid death. To this the state has no right to compel; for a right to which no one was
% --- p. 141 ---
\opage{141}bound would be a self-contradiction. Against unreason and prejudices, religious ones included, the state on the other hand has an unconditional right (tract. polit. III, 7--9. cf. tract. theol.-polit. c. XVII).

Since jus summæ potestatis is nothing other than the very right of nature, it follows of itself \emph{that two different states stand to one another as two individuals in the state of nature}. They are thus by nature enemies and have the right to make war on one another and subjugate one another. If they place themselves in a relation of peace to one another, they are allies (confoederatæ), and their alliance lasts as long as its cause, namely fear of loss or hope of advantage. If this cause is removed, each state has the right to dissolve the alliance, and one cannot say that the state which breaks it acts faithlessly; for the right was equal for both, and the alliance presupposed the determinate circumstances. By this breach the state merely provides for its own security over against the other, and it even becomes an outright duty for the sovereign power to break its promise when it has promised another state something which afterwards either time or reason shows to be contrary to the welfare of the subjects (tract. polit. III, 11--14).

\emph{The supreme power of the state} is as it were the soul of society (imperii mens), which is to lead all; it therefore alone represents the state's right to determine what is good and evil, equitable and inequitable, i.e.\ what all, as individuals and in union, are to do and leave undone; it must therefore give laws, interpret them and judge, declare war and conclude peace, in short administer all public affairs without exception. If one takes the concepts of law and offence in the stricter sense developed above, one must say that the sovereign power cannot be bound by laws, since it itself determines them, and that it thus cannot offend. If, on the other hand, one takes the concepts in a wider sense, then the state can be said to offend when it does what must become its destruction, when it acts against reason, which makes it sui juris, when it does what of necessity---and the necessity of human nature the sovereign power cannot alter---
% --- p. 142 ---
\opage{142}\setcounter{footnote}{0}annihilates the reverence and fear felt for it, and, by turning these into exasperation, turns the social condition into a state of nature. The contract or the laws by which individuals have transferred their right to the sovereign power can be broken in the interest of the common weal; but what serves the common weal can, by the right of the state (jure civile), be decided only by him who holds dominion, and since he has the right of interpretation, no one can, by the right of the state, allege a transgression of the laws against him. If, however, the laws are violated in such a way that the fear of the majority is turned into exasperation, then the state is thereby dissolved and the contract abolished, and right is then asserted not jure civili but jure belli. The ground for keeping the contract thus becomes, for the head of state, the same as the ground, in the state of nature, for the individual not to kill himself (tract. pol. IV, 1--6).

If one compares the different forms of society with one another, then that one must be regarded as \emph{the best} which best realizes what was intended by the union. This was immediately the assertion of right; but since the power for this, both for the state and for the individual, was conditioned by their being led by reason and acting according to its laws, the \emph{best} state becomes the most rational\footnote{Tract. polit. II, 21: Optimi imperii jura ex rationis dictamine institui debent.}; for since reason expresses right and power most completely, it expresses the end most completely. Since the essence of reason is \emph{the true knowledge} of the divine and of our unity with and in it, and since this knowledge is our highest good, the end of the state can also be determined by it\footnote{De intell. emend. pag. 361.}. It can likewise be determined as \emph{peace} and \emph{security} (pax et vitæ securitas; tract. polit. V, 2 f. III, 6)\footnote{As the end of the state is set down (tract. polit. III, 6): pax, securitas, propulsio miseriæ; cf. tract. theol.-polit. pag. 34: finis universæ societatis et imperii est, \emph{secure et commode} vivere (in the same work: pag. 177: secure et \emph{optime} vivere); cf. also Eth. IV. App. c. 14.}, so that
% --- p. 143 ---
\opage{143}\setcounter{footnote}{0}that state becomes the best where men live in concord and the right of each is preserved inviolate, while the most turbulent state, since it comes nearest to the state of nature, corresponds least to its concept. For by peace in the state is understood not the negative: the absence of war, or: sluggish quiet, but something positive. It is a virtue which arises only from strength of soul (animi fortitudo), and is the same as the citizens' steadfast will to obey the laws. The peaceful (concordant) life in the state is thus to be conceived as the true human life, and this is essentially determined by reason, vera mentis virtute et vita, since reason expresses human nature as unity. The determination of the end of the state as peace and concord and as reason and knowledge thus coincide, and from reason \emph{freedom}\footnote{Tract. theol-polit. pag. 226 f.: Ex fundamentis reipublicæ sequitur, finem eius ultimum non esse dominari, nec homines metu retinere et alterius juris facere, sed contra unumquemque metu liberare, ut secure, quoad eius fieri potest, vivat, hoc est, ut jus suum naturale ad existendum absque suo et alterius damno optime retineat. Non, inquam, finis reipublicæ est, homines ex rationalibus bestias vel automata facere, sed contra ut eorum mens et corpus tuto suis functionibus fungantur, et ipsi libera ratione utantur, et ne odio, ira vel dolo certent, nec animo iniquo invicem ferantur. \emph{Finis ergo reipublicæ revera libertas est.}} is likewise inseparable, since only he is free who acts rationally. --- Whether the end of the state is reached or not, Spinoza regards more as the merit or fault of society than of the individual; homines enim civiles non nascuntur sed fiunt (tract. pol. V, 2). % print: uadskilllelig (misprint, three l's)

Spinoza distinguishes the particular \emph{forms of state} by the number of those to whom the right of the state is entrusted. If dominion lies with a council formed of the whole multitude, it is \emph{democracy}; if the council consists only of some chosen few, \emph{aristocracy}; if the whole power is in one man's hands: \emph{monarchy}. He places these 3 forms side by side, and seeks to show the possibility that in each of them the end of the state can be reached; yet he does not regard them as equally
% --- p. 144 ---
\opage{144}\setcounter{footnote}{0}good, but gives democracy the preference\footnote{tract. polit. XI, 1: tertium et omnino absolutum imperium, quod democraticum appellamus. VIII, 3: si quod imperium absolutum datur, illud revera est, quod integra multitudo tenet. VIII, 7: quo jus summæ potestatis majus est, eo imperii forma cum ratione magis convenit; cf. tract. theol.-polit. pag. 216: eodem modo, ac in imperio democratico fieri ostendimus, ubi omnes communi consensu deliberant, ex solo rationis dictamine vivere. Yet Spinoza concedes (tract. polit. XI, 2) that aristocracy would be the most excellent if the election of the patricians were led by reason alone.}, as much because it is the strongest (omnino \emph{absolutum} imperium) as also because there is less danger of irrational decisions when so many take part in the governing power (tract. theol. polit. pag: 180, tract. polit. VIII, 6), and finally because it gives the greatest freedom and equality (tract. polit. X, 8). In developing the particular forms of state, his chief task is to show how they must be arranged in order that dangerous upheavals may be avoided and the form once established be preserved (tract. polit. VI, 2). On the presupposition that reason governed all, everyone would voluntarily subordinate himself to society; for reason, the positively uniting element, would bind all together. But since the presupposition is that the mass of men is governed by the affects, the affects must be balanced by affects, in the interest of reason; there must be given in the institutions an attachment for the interests of all, a counterweight against the irrational passion of the rulers. Only by this equilibrium is quiet brought about, and through it the rational can assert itself. The necessity of finding a counterweight against the passions of the rulers and a guarantee for the continuance of the form of state appears most strongly in monarchy, which by its nature is most exposed to upheaval. For even if despotism seems to bring stability with it, this stability is only \emph{negative}, bondage and sluggishness, not positive, security and concord. On the contrary, there lies in monarchy as absolute not only something hateful to the subjects, and something that conflicts with the
% --- p. 145 ---
\opage{145}\setcounter{footnote}{0}laws of human nature\footnote{Tract. polit. VII, 14: imperio absoluto uti principi admodum periculosum, subditis admodum odiosum, et institutis tam divinis quam humanis adversum, ut innumera ostendunt exempla.}, but there lies in it also a contradiction, since the personality of the single man, regardless of all the contingencies to which it is exposed (tract. pol. VIII, 3), cannot contain the power that can secure his right, for which reason he must always seek help for himself, so that monarchy in reality becomes a hidden aristocracy, and the condition in it violent, unnatural and wretched\footnote{Tract. polit. VI, 8: sequitur, regem eo minus sui juris, et subditorum conditionem eo miseriorem esse, quo magis absolute civitatis jus in eundem transfertur.}. Only when the monarch's power is restricted, so that the fundamental law itself is made independent of the king's will (VII, 1), and the king's power ``is determined solely by the very power of the people and preserved by the very protection of the people'' (VII, 31), can the security of the monarchy and the peace of the subjects be secured\footnote{This is justified and does not conflict with the king's dignity: nam fundamenta imperii veluti regis æterna decreta habenda sunt, adeo ut eius ministri ei omnino obediunt, si, quando aliquid imperat, quod imperii fundamentis repugnat, mandata exsequi velle negent. --- Monarchy is to be so arranged, ut omne jus sit regis explicata voluntas, at non ut omne regis voluntas jus sit. (tract. polit. VII, 1, cf. the comparison with aristocracy VIII, 3).}.

The more precise determinations for the arrangement of the different forms of state\footnote{tract. polit. Cap. 6---11.} have no philosophical interest. This domain is, in Spinoza's view, not that of thought philosophical in the stricter sense, but that of probability (cf. Ep. 60). The determinations he gives therefore have significance only as practical calculations, which refer to something that by its essence has no subsistence, and which constantly lets the truer ground, the proper essence into which it goes back, shine through it, even if only as in the distance. But toward this deeper ground the direction runs, and when it shows itself (cf. the next chapter), the relation on which the doctrine of the state rested vanishes.

% --- p. 146 ---
\opage{146}\setcounter{footnote}{0}In the working out of this doctrine in detail, however, it becomes of interest, as was already intimated above, to see, with respect to Spinoza's conception of the fundamental concepts, how finite reality involuntarily asserts a significance for thought, and, with respect to the development of the special science, how Spinoza, in his sympathy for rational freedom, touches at many points on the political order of a more recent age and asserts the claims of humanity (with respect to the abolition of torture, to the right to freedom of thought, of speech and of religion, and the like).

In the exposition of the doctrine of the right of the state, the point has been shown at which the relation of the state to \emph{religion} comes forward. This relation can first be completely determined when the essence of religion has been developed, and therewith the proper relation of \emph{science} to both will then become clear.

It is in the Tractatus theologico-politicus that Spinoza develops his conception of the essence of \emph{religion} and its relation to the state and to philosophy. In order not to misunderstand the developments he gives, it is necessary here, more than anywhere else, to remember the significance he grants to accommodation\footnote{Tract. de intell. emend. pag. 361 he cites among the provisional rules of life: ad captum vulgi loqui et illa omnia operari, quæ nihil impedimenti adferunt, quo minus nostrum scopum (i.e.\ to reach the highest good) attingamus. Nam non parum emolumenti ab eo possumus acquirere, modo ipsius captui, quantum fieri potest, concedamus, adde, quod tali modo amicas præbebunt aures ad veritatem audiendam.}. The purpose of the treatise was, as was already stated in chapter 1, to make room for free philosophizing, and, in order to achieve this, \emph{on the one hand} philosophy had to be separated from faith (religion), and the different character of the two demonstrated; \emph{on the other hand} the inseparability of religious freedom from the freedom of the state, and the authority of the sovereign power in relation to religious matters, had to be established (cf. tract. theol.-polit. præf. and pag. 30, 160. 166). But thereby all the prevailing representations of the essence of religion
% --- p. 147 ---
\opage{147}\setcounter{footnote}{0}and of the significance of Scripture had to be affected, many prejudices had to be removed, and given the nature of these prejudices and the character of the estate\footnote{``Summum odium, quale theologicum esse solet.''} that defended them, and in them its own dominion, this had to be done with caution. Spinoza therefore seeks not only to make use of the points of attachment he could find in his countrymen's republican feeling for freedom and in their historical experiences, but in his developments he accommodates himself as much as possible to the prevailing representation, uses its forms and its weapons in his argumentation, appropriates its expressions, and takes care only to put another content into these, so that those who were familiar with his thoughts could find them again in the foreign forms, while those who did not know them would be prepared to grasp them. If one does not constantly remember this, one will not be able to avoid misunderstanding Spinoza's writing, as almost all his contemporaries did, and one will then find contradictions in it with his fundamental philosophical view which cannot be removed.

From what was developed in the 5th and 6th chapters it has become clear which form of knowledge Spinoza sets down as the true one, and which goal he sets for this knowledge. \emph{Reason} is the true, adequate knowledge, \emph{the knowledge of God} its goal, and since it completes knowledge, it completes action too. \emph{True} religion\footnote{Tract. theol.-polit. pag. 53: Deum sapienter timere, hoc est, vera religione colere. Eth. IV, 73 schol.: quæ ad veram vitam et religionem spectant.}, as distinct from religion in the ordinary meaning, is then the acting that proceeds from the knowledge of God (cf. the preceding chapter and Eth. IV, 37. schol. 1). A higher knowledge than that of reason cannot be given; for its knowledge of God is adequate, its ``natural light''\footnote{Lumen naturale (lumen naturæ), spiritus, intellectus (int. naturalis) and ratio are used as synonyms. Reason is determined as sufficient for the true knowledge of God (tract. theol.-polit. pag. 54), as a more excellent means to such knowledge than the historical element in Scripture (pag. 64), as agreeing with cognitio supranaturalis (pag. 139), as \emph{necessary} for the knowledge of God (so that faith alone becomes insufficient; pag. 156), as God's verum syngraphum (pag. 168), as uncorrupted and incorruptible (præf.; pag. 168: profecto non satis mirari possum, quod rationem, donum maximum et lucem divinam, mortuis literis et quæ humana malitia depravari potuerunt, submittere velint, et quod nullum existimetur scelus, contra mentem, verum Dei syngraphum, indigne loqui, eamque corruptam, cæcam et perditam statuere, at maximum habeatur scelus, talia de litera et verbi Dei idolo cogitare. Ep. 74: Ratione --- vero Dei verbo, quod in mente est, quodque nunquam depravari nec corrumpi potest. cf. Ep. 36).} is a
% --- p. 148 ---
\opage{148}\setcounter{footnote}{0}divine light, since it is an eternal, immediate, undisturbable, infallible, certain revelation of God's idea and natura, not in words, ``but in a far more excellent way''\footnote{Tract. theol.-polit. pag. 2, cf. pag. 144: ``Dei æternum verbum et pactum'' and ``vera religio'' is ``divinitus inscripta hominum cordibus i.e.\ humanæ menti'', and it is ``verum Dei syngraphum, quod ipse suo sigillo, \emph{nempe sua idea}, tanquam imagine suæ divinitatis consignavit.''}. % print: last quotation in this note opens with » and closes with “ (mixed pair)
Its spirit is a spiritus Dei, spiritus sanctus, spiritus divinus; through it a lex divina universalis expresses itself, and it comes to agree with the lex catholica\footnote{For by the expression catholica Spinoza designates what is common to all men, which as such must agree with reason and the natural light. Thus religio catholica (which is sharply distinguished from rel. \emph{Romana}) is determined as lex divina, per prophetas et apostolos universo humano generi revelata (præfat.); it is set down as synonymous with religio \emph{toti humano generi universalis} (pag. 148), and as \emph{maxime naturalis} (pag: 149); there is talk of a rel. catholica \emph{lumine naturali et prophetia revelata} (Ep. 49), of catholica documenta rationis (tract. theol.-pol. pag. 217), of a \emph{lex} catholica, which becomes equal to lex divina universalis (pag. 151 f.), and of a \emph{fides} catholica (pag. 160), whose object becomes the essential content of religion, the content that agrees with reason.}, religio catholica, fides catholica.

If, therefore, there is talk of a \emph{supernatural} knowledge\footnote{Spinoza himself, where he criticizes the representation of miracles, rejects the distinction between what is \emph{above} nature and what is against nature, since the former is placed within the limits of nature and must thus make a breach in its eternal order.} (cognitio supranaturalis), it is to be expected that Spinoza must
% --- p. 149 ---
\opage{149}\setcounter{footnote}{0}have used this expression improperly, and that one will find this intimated in him. This is also confirmed at once by his development of the fundamental religious representations: \emph{revelation} and \emph{prophecy}. Prophecy he first defines in general as: rei alicuius certa cognitio a Deo hominibus \emph{revelata}. But since ``revelation'' thus comes to comprise natural knowledge as well, since this depends on the knowledge of God and on his eternal decretum\footnote{Tract. theol.-polit. pag. 216: Religio eadem est, \emph{et a Deo æque revelata}, sive hoc sive illo modo (i.e.\ sive religionem lumine naturali sive prophetico \emph{revelatam} concipiamus) hominibus innotuisse supponatur. Cf. Cog. metaph. II, c. 12: decreta Dei, \emph{lumine naturali revelata}.}, natural knowledge is now distinguished more precisely from supernatural knowledge by this, that the latter (``quæcunque illa sit'', pag. 1) ``extends beyond the limits of the former''\footnote{Tract. theol.-polit. pag. 14: Non dubium est, eos multa \emph{extra intellectus limites} percipere potuisse.}, and that ``the laws of human nature considered by itself cannot be its causes'' (cf. the definition of the affects as passive states). The character of this specific revelation Scripture must teach us to know, since it is an empirical object that cannot be determined by intellectus. As its \emph{means} Scripture now gives words, signs, or both together, and these either real or imaginary. But as Spinoza conducts an ironical argumentation from the narratives of Scripture, he raises numerous difficulties, which he does not resolve, against the reality of the words and signs; and since words and signs, even as real, belong to the imagination, he can conclude with the result: that ``apart from Christ no one has received God's revelation in any other way than \emph{by the help of the imagination}, namely through words or images'', so that revelation and prophecy are put together with the imagination as proceeding from it (cf. tract. theol.-polit. pag. 13 f.). Therewith
% --- p. 150 ---
\opage{150}\setcounter{footnote}{0}they are separated from intellectus, ``for then things are comprehended when they are grasped by the mind alone (pura mente), independently of words and images'' (pag. 50). In \emph{Christ}\footnote{Where Spinoza speaks of \emph{Christianity}, one must remember that he uses ``christiana religio'' as synonymous with ``amor, gaudium, pax, continentia et erga omnes fides'' (tract. theol.-polit. præf.), thus with what he himself understands by religio (Eth. IV. 37. schol. 1), and regards it, in its simplicity (pag. 143 f.), as agreeing with reason. \emph{The spirit of Christ} therefore designates justice and love of one's neighbour (Ep. 49. 74), and he who has this spirit of love, even if he professes another religion than the Christian (Ep. 49), knows Christ ``according to the spirit'', and Christ is in him (pag. 164). This spirit becomes equal to idea Dei, ``on which alone it depends that man is free'' (Eth. IV, 68. schol.); it contains blessedness within itself (tract. theol.-polit. pag. 65); through it one grasps ``God's laws'' as ``eternal truths.'' Thus it becomes equal to spiritus sanctus. \emph{Christ} himself is conceived symbolically as ``God's eternal Son, i.e.\ God's eternal wisdom, which has revealed itself in all things, but most in the human mind, and most of all (omnium maxime) in Jesus Christ'' (Ep. 21). And as a historical, \emph{human} person Christ is ascribed an adequate knowledge, so that he becomes as it were the image of God's wisdom (tract. theol.-polit. pag. 7. 50); he is called os Dei, a teacher who was sent ad solam legem universalem docendum (pag. 56). But the dogma of the Incarnation in its ecclesiastical form Spinoza regards as absurd: that the infinite should have taken on a finite nature he regards as just as self-contradictory as if the circle had taken on the nature of the square (Ep. 21, cf. Ep. 23. 74 and de intell. emend. pag. 374.). And with the dogma of the Incarnation he rejects all the other dogmas that conflict with reason, ``the true word of God, which is in our mind and can never be corrupted.'' He calls them ``horribilia secreta'' (Ep. 74), ``absurda arcana'' (tract. theol.-polit. præf.). The narrative of Christ's Resurrection and Ascension he explains symbolically (Ep. 23 and 25), and in general he declares that knowing Christ according to the flesh (secundum carnem) is not necessary to beatitudo (Ep. 21).} revelation is not through the imagination and by its means, but ``immediately, mind to mind'' (de mente ad mentem); but he does not thereby become specifically different from other
% --- p. 151 ---
\opage{151}\setcounter{footnote}{0}men; for in their mind too ``God's mind and his eternal pronouncements (sententiæ) are inscribed.''

That the prophets had God's spirit therefore signifies, according to the usage of the Old Testament, which Spinoza develops in detail, not that they had a higher insight, but only that they were in possession of a \emph{peculiar} and \emph{uncommon} virtue (capacity), that they cultivated virtue with steadfastness, and that they perceived God's word or thought, and finally, that the cause of prophetic knowledge was not comprehended. According to what laws of nature this knowledge took place, Spinoza leaves undecided.

From the significance of the imagination for prophecy it now follows (by the essence of the imagination; cf. p. 88 ff.) that the prophets grasped almost everything in parables and riddles, and everything spiritual in bodily forms, and that prophecy, like the imagination, must be changeable, short-lived, and cannot, as reason can, contain certainty (tract. theol.-polit. cap. 1). Therefore the prophets themselves had to have an external guarantee. This they got in the \emph{signs}; but the certainty these could give was only a moral one, since false prophets too performed signs (and besides, the sign in general, by its nature as an object of the imagination, cannot give real certainty). The greatest moral certainty must be sought in \emph{the moral character of the prophets}; yet this too in turn remains imperfect, since no one is completely virtuous. With respect to the signs, it now follows again from their relation to the imagination that they had to conform to the prophet's representations and power of comprehension, and therefore be different for different prophets. For the same reason the content and form of revelation had to vary according to the prophets' different temperaments, circumstances of life and opinions, and according to the peculiar disposition of the imagination. This Spinoza carries out circumstantially, relying on statements of Scripture, but in such a way that his doctrine of the imagination and the affects constantly shines through the biblical demonstration. The result thus becomes that the prophets held different, indeed opposed, opinions, and many
% --- p. 152 ---
\opage{152}\setcounter{footnote}{0}errors (with respect to the merely speculative, not with respect to what concerns uprightness and good morals), and that prophecy never made the prophets more knowing (doctiores), but left them in their preconceived opinions, and that we are therefore not at all obliged to have faith in them\footnote{Tract. theol.-polit. pag. 28: Concludimus itaque, nos prophetis nihil aliud teneri credere præter id, quod finis et substantia est revelationis (i.e.\ vera vita).} with respect to the purely speculative (tract. theol.-polit. pag. 21). Likewise one must suppose that \emph{the prophets were ignorant of many things}, which already lies in their having been caught up in the common representations of their time, as one can see from manifold examples from the Old Testament. The reason why they are commended, then, becomes not their insight\footnote{Tract. theol.-polit. pag. 23: Nihil enim singulare de Dei attributis docuerunt, sed admodum vulgares de Deo habuerunt opiniones, ad quas etiam eorum revelationes accommodatæ sunt. cf. Ep. 32.} or superiority of mind, but their piety and constancy; their task was to teach true virtue, and from this it follows, on the one hand, that they must have been found among \emph{all} nations (chap. 3), on the other, that they are authorities for us only with respect to what concerns it.

The relation between natural knowledge and so-called supernatural revelation thus proves to be this: that the former comprises the eternal, unchangeable, universally valid truth, certain by itself\footnote{Ibid. pag. 173 f.: De veritate --- --- et certitudine rerum, quæ solius sunt speculationis, nullum spiritus testimonium dat præter rationem, quæ sola, ut jam ostendimus, veritatis regnum sibi vindicavit. --- Spiritus sanctus ``bears witness only to \emph{good works},'' and this witness is one with ``the peace of mind that arises in the mind from good works.''}, while the latter, as determined by the form of the imagination, is affected by it in its content too, and can therefore indeed contain ``more'' than the former, but not more truth\footnote{Cognitio naturalis and revelatio could therefore be set down as opposites (tract. theol.-polit. pag. 141, cf. princ. phil. Cart. II, 13. schol. Cog. metaph. II. c. 12), and the relation between cognitio supranaturalis and \emph{legitime argumentari} is determined as an inverse relation (tract. theol.-polit. pag. 139). The ``more'' in revelation becomes ``res imperceptibiles, quas tantum imaginari possumus'' (pag. 97).}.

% --- p. 153 ---
\opage{153}\setcounter{footnote}{0}Imagination and ``revelation'' thus become inseparable. The great multitude of men is led not by rational insight, but by the affects. They must therefore be influenced through these, and since the affects are bound up with the inadequate representations of the imagination, these become the means. But when the truths of reason are taken up into the form of these representations, they change their nature, and the fundamental representations of \emph{religion} become inadequate. Thus, while reason conceives existence as governed by the law of necessity of the divine nature, and sees the effect bound to the cause by the bond of necessity, representation conceives law\footnote{In this meaning law is opposed to æterna veritas (pag. 51), law to moral teachings (pag. 87 f.), lawgiver to teacher (pag. 89). Law in the vulgar meaning man can transgress, but not God's æternum decretum, which is ``in universa natura inscriptum'' (tract. polit. II, 22).} as an arbitrary one, given by God in the same way as human laws are given, and arbitrarily connected with reward and punishment, which through hope and fear can move men to do and leave undone what is commanded or forbidden. The universal, natural \emph{lex divina}, which according to the nature of my essence\footnote{Precisely for this reason it is universal, not tied to promises and threats, and contains beatitudo (chap. 4---5). --- The virtuous man will do the good and refrain from the bad because it agrees with his nature, not because a law sets rewards or punishments for it (Ep. 34).} determines its highest perfection as the knowledge of God and the love of God that follows from it, changes into a special law, which cannot do without peculiar \emph{ceremonies} affecting the imagination, and must be supported by something historical and by \emph{experience}, which cannot give clear insight, but only call forth obedience and devotion. Reason's view of the inviolable order of nature is transformed into the representation of
% --- p. 154 ---
\opage{154}\setcounter{footnote}{0}\emph{miracles}\footnote{On miracles cf. Cog. met. II. c. 9. 12. Ep. 23. 74.}, of a supernatural, immediate divine intervention that suspends the eternal laws of nature, the very laws of God's essence; and one now believes one has in miracles a surer proof of God's essence and existence than in those laws themselves, while nevertheless such a deviation from the universal laws of nature would bring discord into the order of nature and set the most original concepts of our mind, which stand and fall with its validity, in conflict with themselves, overthrow all certainty and precisely make God's existence impossible\footnote{Tract. theol.-polit. pag. 77: Quidquid contra naturam est, id contra rationem est, et quod contra rationem, id absurdum est, et proinde etiam refutandum. The miracle becomes a relative concept and designates only: opus, cuius causam naturalem exemplo alterius rei solitæ explicare non possumus, vel saltem ipse non potest, qui miraculum scribit aut narrat. --- Spinoza's critique of miracles finds indirect application to the concept of revelation, to which he himself also points the way (pag. 81).}.

From this fundamental difference between the conception of reason and that of religious representation, it follows that \emph{Scripture}, which contains revelation, is a peculiar domain, with a distinctive character conditioned by the power of comprehension of the multitude, and \emph{that Scripture and philosophy must therefore be set down as mutually independent}. Scripture must be regarded as something historically given, which must be grasped as such, and thus interpreted by itself. Interpretation must proceed purely objectively, not prove its truth, but establish what is factual\footnote{Tract. theol.-polit. pag. 87: tota bibliorum cognitio ab iisdem solis est petenda. --- pag. 84: dico methodum interpretandi scripturam haud differre a methodo interpretandi naturam, sed cum ea prorsus convenire. --- pag. 86: de solo sensu orationum, non autem de ejus veritate laboramus. Yet once the factual has been found out, reason must necessarily judge the truth of revelation (of prophecy) (Ep. 34. Cog. metaph. II. c. 8. Tract. theol.-polit. pag. 167 f.: Verum quidem est, scripturam per scripturam interpretandam esse, quamdiu de solo orationum sensu et mente prophetarum laboramus, sed postquam verum sensum eruimus, necessario judicio et ratione utendum, ut ipsi assensum præbeamus).}. The means to this is an unbiased \emph{theory of Scripture},
% --- p. 155 ---
\opage{155}\setcounter{footnote}{0}built on knowledge of the language of Scripture, of the history of the books, of the personal lives, circumstances of their times and fortunes of the individual authors.

Scripture, regarded in this way, will prove to consist of very heterogeneous books, from different times, as an incomplete collection of historical documents. It will prove that there is found in it no agreement with respect to the speculative, which by the nature of prophecy must have been conceived differently by the different prophets, but only with respect to the practical\footnote{With the New Testament the relation is somewhat different than with the Old, since in the apostles the element of revelation recedes before natural knowledge; yet discrepancies are found in them too.}. Scripture can be called \emph{holy} and \emph{divine} only because it has regard to piety and religion; these designations are conditioned by the use that is made of the thing. ``Nihil extra mentem absolute, sed tantum respective ad ipsam sacrum aut profanum aut impurum est.'' \emph{``The word of God''} must acquire a far smaller extent than Scripture, which is called the word of God only because it contains what is essential to religion. While Scripture has undergone manifold changes, \emph{the word of God contained in it is uncorrupted.} The extent of ``the word of God'', of the essential teaching of Scripture, must be determined by the purpose of Scripture: to make itself intelligible to all men. It can therefore not comprise subtle speculations, which are accessible only to few, but only the very simplest things, which even the most simple-minded can grasp. Since it does not wish to teach science, it requires not insight, but only \emph{obedience} (pag. 154), and since obedience to God shows itself \emph{only} in love of one's neighbour, the essential content of the word of God becomes: to love God above all things and one's neighbour as oneself (pag. 151), together with the propositions that follow from this with necessary consequence: ``that God exists, takes care of all, that he is almighty, that by his decision it goes well with the pious, and
% --- p. 156 ---
\opage{156}\setcounter{footnote}{0}ill with the ungodly, and that our salvation depends solely on his grace'', and finally the general moral precepts, which are more determinate expressions of this\footnote{Tract. theol.-polit. pag. 151. cf. pag. 63 and pag. 163 f.: The essential dogmas are contained in the proposition: dari ens supremum, quod justitiam et caritatem amat, cuique omnes, ut salvi sint, obedire tenentur, eumque cultu justitiæ et caritate erga proximum adorare. Therein lies: 1) Deum, hoc est ens supremum, summe justum et misericordem, sive veræ vitæ exemplar existere; 2) Eum esse unicum; 3) Eum ubique esse præsentem, vel omnia ipsi patere; 4) Ipsum in omnia supremum habere jus et dominium, nec aliquid jure coactum, sed ex absoluto beneplacito et singulari gratia facere; 5) Cultum Dei eiusque obedientiam in sola justitia et caritate sive amore erga proximum consistere; 6) Omnes, qui hoc vivendi ratione Deo obediunt, salvos tantum esse, reliquos autem, qui sub imperio voluptatum vivunt, perditos; 7) Denique Deum poenitentibus peccata condonare (for otherwise all would despair, whereas this faith strengthens the love of God). The different possible conceptions of these propositions have no influence on faith, provided only that they do not act harmfully on action. --- It is easily seen that among the determinations given there occur several which Spinoza elsewhere rejects as anthropomorphic (e.g.\ God's justitia and misericordia, cf. pag. 50, 51, 158), and that the expressions are so framed that his own view can be put into them, as e.g.\ in the doctrine of the forgiveness of sins.}. Scripture thus requires only the knowledge that is necessary for obedience. For since man at this standpoint is regarded as led only by affects, not by insight, only that knowledge acquires significance which in the form of affects can influence the affects, thus only such representations of God as can directly influence men's way of acting, e.g.\ the representation of divine justice and love, whereas men, with respect to the knowledge of God's essence as it is in itself, without direct regard to human action, can ``sine scelere toto coelo errare'' (pag. 157). Only the relation to actions gives opinions the character of pious or impious: ``so that, if anyone, by believing what is right, becomes refractory (contumax), then in reality he has an impious faith,
% --- p. 157 ---
\opage{157}\setcounter{footnote}{0}whereas if anyone, by believing what is false, becomes obedient, he has a pious faith''. What is required in Scripture is simple faith: ``To believe in God, to have reverence for him, or, what is the same, to obey God'' (pag. 160). Faith as objective therefore becomes not a sum of dogmas, but is nothing other than ``to have such a representation (talia sentire) of God that it cannot be dispensed with for obedience to him, and that conversely, when obedience is given, this representation must be an inseparable consequence.'' (pag. 161). It is thus not saving ``per se'', but merely ``ratione obedientiæ'', and he who is truly obedient necessarily has a true and saving faith. Faith is judged \emph{solely} by works, and requires not so much true as pious (pia) dogmas, i.e.\ such as move the mind to obedience (yet he who professes them must not see that they are false, for then accepting them becomes an impossibility). Everyone therefore also has the right to interpret the dogmas in the way they become most suited to prompt him to obedience. % print: quotation opens „ (p. 156) and closes « (p. 157), mixed pair

Once the essence of faith has thus been determined, the relation between it and philosophy can now be definitively stated: ``The aim of \emph{philosophy} is nothing but truth, that of \emph{faith} nothing but obedience and piety\footnote{Tract. theol.-polit. pag. 170: Ratio regnum veritatis et sapientiæ [est], theologia autem pietatis et obedientiæ; cf. princ. phil. Cart. II. 13. schol.}. The foundation of philosophy is the common notions (notiones communes), and it has its source in nature alone; the foundation of faith, on the other hand, is history and language, and its source is Scripture and revelation alone'' (tract. theol.-polit. pag. 165). There is thus no community and no kinship between the two; they are as far apart as heaven and earth. Faith makes no encroachment on the freedom to philosophize; it grants everyone the right to hold any opinions whatsoever, and condemns as heretics and schismatics only those who preach insubordination, hatred, strife and anger, while it counts as pious only those
% --- p. 158 ---
\opage{158}\setcounter{footnote}{0}who, according to the capacities of their intellect, prompt to justice and love. Just as Scripture is not to be reinterpreted according to reason, so reason is not to subordinate itself to Scripture. The proposition on which the whole of theology\footnote{Tract. theol.-polit. pag. 170: Theology is equal to revelatio, quatenus indicat scopum, quem diximus Scripturam intendere (nempe rationem et modum obediendi, sive veræ pietatis et fidei dogmata), hoc est, id quod proprie vocatur Dei verbum.} rests, that men can be blessed by obedience alone without the knowledge of things\footnote{Hoc est, quod ad salutem sive beatitudinem satis sit, divina decreta tanquam jura sive mandata amplecti, nec opus sit, easdem ut æternas veritates concipere, non ratio, sed revelatio docere potest. --- Spinozæ Opera ed. Paulus I, p. 442).}, cannot indeed be proved by reason and the natural light, which must set down blessedness as one with the knowledge of God. Consistently, then, it would have to be rejected; for reason alone decides truth (pag. 170), also with respect to the understanding of the dogmas. But this consequence Spinoza does not state; he contents himself with giving the premises for it. Instead he adduces what can give this doctrine a ``moral certainty.'' This becomes essentially \emph{the authority of the prophets}, which is grounded on the moral content of their teaching, agreeing with reason, and which brings us to assume: ``that they taught, not deceitfully, but with full conviction (ex vero animo), that men become blessed through obedience and faith, and since they have moreover confirmed this by signs, we are persuaded to believe that they did not say it lightly (temere), nor were they out of their minds when they prophesied.'' Moreover, this doctrine, which can be accepted without danger and harm, also contains a great consolation for those who are not very strong in reason (qui ratione non ita pollent). ``For since absolutely all can obey, whereas, in comparison with the whole human race, there are only few who acquire a virtue led by reason alone, we would, if we did not have this testimony of Scripture, despair of the salvation of almost all'' (cap. XV. the conclusion). One must therefore not % print: closing ) in the second note without an opening (
% --- p. 159 ---
\opage{159}\setcounter{footnote}{0}reject this doctrine out of hand because it cannot be proved mathematically\footnote{How Spinoza's utterances about blessedness through obedience are to be understood lies immediately in his doctrine of the affects and of the imagination. In one of the notes added later to the tract. theol.-polit. (cf. Bruder's edition III, p. 217) he sets ``the love of God'', which is a virtue that with necessity inheres in the man who knows God, against ``obedience'', which merely has regard to the will of the one who commands, not to the necessity and truth of the thing. Obedience is thus conditioned by the imagination's conception of the eternal truths as arbitrary laws: a conception which reason corrects, and as it is corrected, obedience passes over at once into love, which follows from true knowledge with the same necessity as light from the sun. ``\emph{Ex rationis igitur ductu Deum quidem amare, sed non obedire ei possumus}'' cf. tract. polit. II, 22.}.

From what was developed at the beginning of the chapter concerning the right of the state and of the sovereign power, it follows, when it is applied to religious matters, that \emph{jus circa sacra} must lie with the head of state. For since justice and love are not the ruling powers in the state of nature, but only become so in the state, they can first acquire legal force through the right of the state, and this is the same as: through the decision of those who have the right to rule. But since the kingdom of God consists solely in the right (law) of justice and love, the consequence again becomes \emph{that God has no kingdom except through those who hold dominion.} This also follows from the fact that reason does not comprehend that God as prince and lawgiver gives men laws (see above), since, then, the divine teachings cannot acquire the force of law immediately from God, but must acquire it mediately through those who have the right to rule and give laws, and God only through them becomes ruler and lord over human affairs. Divine right must thus depend on the decision of the sovereign power\footnote{Concerning the extent of the right of the sovereign power over against the Church cf. tract. theol.-polit. pag. 221.}, and the
% --- p. 160 ---
\opage{160}\setcounter{footnote}{0}interpretation of it must likewise belong to that power. The outward worship of religion (cultus externus) and all exercise of piety must be directed according to the peace and concord of the state. For without the state nothing good can subsist; therefore the welfare of the state must be the highest law; but what the welfare of the state is, only the rulers of the state can decide; it is therefore their right to determine in what way everyone is to show piety (pietas) toward his neighbour, that is to say: in what way everyone is obliged to obey God\footnote{We are obliged to show pietas \emph{toward all}; regard for the general therefore becomes determinative, but this regard is taken by the head of state.}. It is the duty of the subjects in this respect to obey unconditionally the commands of the authorities, and the right of the authorities is limited only by impossibility and contrariety to reason. ``Piety itself and the \emph{inner} worship of God,'' (Ipsa pietas et Dei \emph{internus} cultus), becomes a domain which the sovereign power cannot reach\footnote{The individual's jus libere sentiendi de religione, libere judicandi, eandemque sibi explicandi et interpretandi, is inalienable (tract. theol.-polit. pag. 103; cf. pag. 215. 225). Religion as consisting in animi simplicitas and veracitas is nullius juris neque autoritatis publicæ. pag. 102. --- The freedom of religious opinions lies, moreover, already in Spinoza's conception of the relation of religion to the imagination.}, but neither can it wish to encroach upon it, for the end of the state and of piety is the same: love (caritas) and concord (tract. theol.-polit. c. 16. 19. cf. tract. polit. III, 10. VI, 40. VII, 26, VIII, 46 ff.). % print: closing ) after "225." missing in the second note

From this right belonging to the sovereign power, and from the essence of religion, it now follows again that it will be dangerous for the state and for religion itself when the clergy is permitted to decree laws and dogmas, to refer the purely speculative to divine right, and to give laws about disputable opinions. (This is shown in chapter 18 by examples from the history of the Jews.) In this direction, therefore, \emph{science} can demand the protection of the state against clerical encroachments. And the state must in general respect \emph{freedom of}
% --- p. 161 ---
\opage{161}\setcounter{footnote}{0}\emph{thought}, since the right to think is inalienable\footnote{Tract. theol.-polit. pag. 231: Omnes homines non possunt eadem sentire. --- pag. 227: Fieri non potest, ut omnes æque eadem sentiant et uno ore loquantur.}, and likewise \emph{the freedom to utter one's thoughts,} which is so deeply grounded in human nature that it cannot be taken from men. Yet freedom of speech can be limited in so far as it immediately contains an action or has a hostile tendency against the sovereign power (pag. 227. 228). But every unjustified attempt to suppress freedom of thought and of speech must prove fruitless and at the same time dangerous for the sovereign power itself and harmful to morality and public quiet, since it will call forth faithlessness, hypocrisy, insubordination, discord and the decline of the state in spiritual respects. The end of the state is indeed precisely this: not to cow, but to free from fear, to make rational development, or freedom, possible: finis reipublicæ revera \emph{libertas} est. The best state must therefore grant everyone the same freedom to philosophize as faith grants everyone (pag. 229). And that the freedom of thought which is demanded is possible without harm to the common weal, and that the standing of the authorities is sufficient to prevent abuse, Spinoza shows by an example from his native land, by the example of the city of Amsterdam, where the different sects and nations live together in concord and peace. He can then conclude the treatise with the result: ``Nothing is more conducive to the security of the state than that piety and religion be placed \emph{solely} in the exercise of love and equity, and that the right of the sovereign power, with respect to both the sacred and the secular, be extended only to actions, so that it is otherwise granted to everyone to think what he will and to utter what he thinks'' (pag. 233, cf. the Preface).

\bigskip
\begin{center}\rule{2.5cm}{0.4pt}\end{center}
\bigskip

\newpage

\chapter*{Chapter Eight.}
\addcontentsline{toc}{chapter}{Chapter Eight. Man's Freedom and Eternity}
\markboth{Man's Freedom and Eternity}{Man's Freedom and Eternity}

\begin{center}
\textbf{Man's Freedom and Eternity.}
\end{center}

% --- p. 162 ---
\opage{162} In Chapter 6, through the doctrine of the affects and
their power over man, man's condition of bondage was
developed. Through the determination of the relation of reason to
the affects, and of the conception of the affects as good and evil
conditioned by this relation, we were led to the opposition between
the concept of the free man and that of the slave (of the wise man and of the ignorant);
but this opposition was still held as relative, within
a more and a less. The same relativity remained
in the doctrine of the social relation and of religion, even though
its removal lay as a simple consequence in the
hints given. What now remains is to determine decisively
what was earlier only hinted at: the concept of human
\emph{freedom}; but for the time being the development of this concept
again takes the form of a doctrine of \emph{liberation}; precisely because the point of departure
was originally taken from man represented as a single, finite modus,
dominated by the affects. From
that very thing which is conceived as finite, the possibility of overcoming finitude,
limitation, is to be shown; but this
overcoming can only be partial: the finite's approximation to
the infinite, going on without end. But through
the consideration of reason, of the essence of the mind, the
true conception then breaks forth: limitation
and time vanish; eternity as being swallows up the becoming of time into itself.

On this point it will be necessary to attend closely
to the hints toward the right conception that are given
in Spinoza. Precisely here, where he collides most strongly
with the common consciousness, he must express himself accommodatingly,
and as it were veil his view.

From the development of the essence of the affects as passive it followed
that they had their origin in the confused and unclear
ideas. Through them the mind was passive, whereas it is
% --- p. 163 ---
\opage{163}\setcounter{footnote}{0}active as producing clear and distinct ones. These latter
must therefore become the most essential remedy against the affects\footnote{Just as, in the development of human bondage, the starting point
was taken from the body, so, in the development of freedom, the starting point is taken
from the mind. A dualism involuntarily betrays itself.}
in that they themselves act as \emph{active} affects. The affect, which is
a passion, ceases to be one as soon as we form a
clear and distinct idea of it; for it is a passion only as a
confused idea, and the idea we form of it is not
different from the affect itself insofar as this is referred to
the mind alone (Eth. V. 3). But since in all things there is something
common, and what is common can only be conceived adequately, there is
no affection of the body of which we cannot
form some clear and distinct idea, so that
it is therefore in everyone's power, at least in part, to gain
mastery over the affects. When the affect becomes the object
of a clear and distinct knowledge, the mind,
when it is affected by it, is determined to think that which
it conceives clearly and distinctly, and in which it fully
rests; and the affect itself, so conceived, is separated from
the thought of the external cause (since through the thought of what is
common the separation of limitation is removed) and is joined
with true thoughts. Through this separation not
only do passionate love and passionate
hatred cease (cf. their definitions); but the appetites that usually
arise from these affects are prevented, by becoming rational,
from rising to excess. For it is the same
appetite by which men are said to act and to be passive;
thus ambitio and pietas are the same appetite, that
men should live according to our disposition; but the former is irrational,
the latter rational, since the former has arisen from inadequate,
the latter from adequate ideas (Eth. V, 4 schol.).

The adequate knowledge of the affect is thus the
most essential remedy against it as passion. In conceiving
what is common, it removes, as we saw, the representation
of the external, single cause, and frees us from the
% --- p. 164 ---
\opage{164}\setcounter{footnote}{0}passionate relation to it; but at the same time, since
through it we conceive things as \emph{necessary}, it brings it about that we are affected
less by them than if we conceived them as free
(Eth. V, 5); and because it conceives what is always present
and always equally common, the active affects that proceed
from it become, with respect to duration in time, stronger than those
that refer to single things, which can be regarded as
not present; the causes that call them forth become
more manifold, and stronger through their combined action. And since
the mind, as long as it does not struggle with affects that are
contrary to our nature, i.e. that hinder the mind from understanding,
has the power, as an active mode of thought, to form clear and
distinct ideas, and to deduce others from them, --- whence it,
on account of the parallelism of mind and body, follows that
the affections of the body are at the same time connected according to the order
of reason\footnote{Spinoza is compelled by the assumption of this parallelism, in positing
a bodily affection that corresponds to the mind's adequate idea,
to posit an imagination that has the same object as the intellect
(Eth. V, 1. 12. dem.), and that thus can only improperly keep
that name (cf. Eth. II, 17 schol.).} --- it thereby gains a greater power of resistance
against the passive states, which cannot so easily
disturb what is rationally linked together as what is uncertain
and vague.

Since everything that we understand clearly and distinctly is referred back
to God, and since in every affect there is something that can be understood
thus, the mind can refer every affect back to
him. But when man understands himself and his affects
clearly and distinctly, \emph{he loves God}, and that the
more highly, the more comprehensive the knowledge is (Eth. V,
14 f.). Love arises when the idea of God accompanies
the distinct idea (cf. the definition of love). It
must, above all else, occupy the mind, since it is called forth
by every affection. No one can hate God; for since the idea of God
in us is clear and distinct, we \emph{act} in having
% --- p. 165 ---
\opage{165}\setcounter{footnote}{0}it, and consequently no affect of sorrow can arise from it\footnote{God cannot even be said, insofar as he is the cause of things, to
be the cause of sorrow when they affect us; for insofar as we understand
the cause of sorrow, it ceases to be a passive state, i.e. it ceases
to be sorrow, and insofar, therefore, as we understand God as the cause of sorrow,
we feel joy (Eth. V, 18. schol.).}. Nor can it, as the highest good
common to all men,
be tainted by the affect of envy or of jealousy,
but it is strengthened ever more, the more men we
imagine as united with God by the same bond
of love. There is no affect that is directly
contrary to it and that can destroy it; \emph{it is therefore the most constant of all the affects},
and insofar as it
is referred to the body (quatenus ad corpus refertur), it can
not be destroyed except with the body itself (Eth. V, 20.
schol.). Since its object is unchangeable and eternal, and
truly in our power, it itself acquires the unchangeableness of its object\footnote{Cf. de intell. emend. pag. 357 f.}, can grow ever stronger, more and more
displace the passive affects, and fill the greatest
part of the mind.

Even at this point, then, where Spinoza has brought
forward the love of God as the strongest and most unchangeable
affect, against which no contrary affect can
hold its ground, man is nevertheless still posited as not completely
freed from the affects\footnote{Herein lies the same contradiction that appeared in the doctrine of
knowledge and of the affects through the conception of man as
abstractly finite.}, but only as confining
them to a small region; and the difference between the free man
and the slave still remains quantitatively determinable by the
different proportion between the adequate and inadequate ideas
of which each one's mind consists. This halting at an approximation
has, as already touched upon, its ground in the point of departure
taken; but in the concept of the knowledge of God and
the love of God a standpoint is given from which
this point of departure vanishes, and the conception consistent with Spinoza's
% --- p. 166 ---
\opage{166}\setcounter{footnote}{0}fundamental view comes forth. Spinoza
presents this at the end of his Ethics; but partly
in such a way that in his expressions he accommodates himself to the common
representation, partly in such a way that at the end
he again adds determinations that essentially concern the relative,
merely quantitative opposition.

He makes the transition from the doctrine of the remedies against
the affects to the new developments with the words (Eth. V, 20.
schol.) that he has now finished with what concerns \emph{the present
life} (vitam præsentem), and that it is now time
to pass on to what pertains to the duration of the mind
(durationem) without relation to the body\footnote{Cf. Eth. V, 40. schol.: de mente, quatenus sine relatione ad corporis existentiam consideratur.}. By the expression
``the present life'' the relation to the imagination is indicated,
since the determinations of time have significance only for
it. By the expression \emph{duratio} it is indicated that the talk
of the mind without relation to the body\footnote{Eth. V, 23: Mens humana non potest cum corpore absolute destrui,
sed ejus aliquid remanet, quod æternum est.} is accommodating,
for we do not ascribe to the human mind a duration
(duratio) that can be determined by time except insofar as it
expresses the actual existence of the body, i.e. ipsi durationem
non tribuimus nisi durante corpore (Eth. V, 23. demonst.);
and according to Spinoza's fundamental conception it is just as impossible
to think a mind without a body as a body without
a mind (cf. Ep. 58). What is therefore developed in connection
with the representation of the immortality of the soul is a
development of \emph{the mind's true, eternal essence as the idea
of the body conceived under the form of eternity.}

The mind's highest \emph{virtue}, i.e. its essential \emph{nature} (V,
25. dem.), and therefore its highest \emph{striving}, is to understand
things by the ``third kind of knowledge'', \emph{intuitive}
\emph{knowledge}, which proceeds from the adequate knowledge
of God's attributes to the adequate knowledge of the essence of
things. This knowledge of things is \emph{direct}
% --- p. 167 ---
\opage{167}\setcounter{footnote}{0}knowledge of God; for through it they are known in their unity
with him (Eth. V, 24, cf. Eth. II, 45. schol., tract. theol.-polit. pag. 46 f.). It expresses most perfectly the human
essence in its activity; it is thus the source of the
highest acquiescentia, i.e. the highest joy, accompanied by
the idea of oneself and one's virtue. But since this kind of knowledge
understands things in God, the joy is necessarily accompanied by
the idea of God as cause, and thereby arises \emph{amor Dei
intellectualis}, i.e. a love of God, not insofar as
we imagine him as present through the imagination,
but insofar as we understand through reason (intellectus) that he
is eternal (Eth. V, 32. schol.).

But at this point, through the essence of the highest
knowledge and through the intellectual love
of the eternal, a light falls back on the preceding development.
The third kind of knowledge can indeed arise
from the second, whose ideas are also true and adequate,
but not from the first, the imagination, with its incomplete
and confused ideas\footnote{The denial of this transition from inadequate to adequate
knowledge, i.e. from the being of the mind as finite to its unity with
the infinite, is the necessary parallel to the denial of the transition
from the infinity of substance to existing finite things
(cf. chapter 2).}. But only for the first
kind of knowledge does time have significance; adequate knowledge
is in the form of eternity; for it time becomes something
unreal, vanishing. The third kind of knowledge has its
adequate (formal) ground in \emph{the mind as eternal}, i.e. in the mind
insofar as it expresses the essence of the body under the form of
eternity, not insofar as it expresses the present,
actual existence of the body, which is the condition for the imagination and
for the conception of things in relation to time and duration.
But since the cause, the mind, is eternal, the effect, the
intuitive knowledge, must necessarily be so too; it is
an eternal logical consequence of the mind's nature (Eth. V, 33. schol.).

That the mind is \emph{eternal} follows from the fact that in God, as the
cause not only of the existence of the human body,
% --- p. 168 ---
\opage{168}\setcounter{footnote}{0}but also of its essence, --- which therefore must be known with eternal necessity
through the very essence of God, --- there must
be given an idea that expresses the essence of the human body
under the form of eternity. But since the human
mind is the idea of the human body, that idea must
belong to the essence of the mind, which thus becomes eternal. This
eternity of the mind cannot, since eternity has no relation
to time, be the object of a memory; but we nevertheless have
consciousness of it; for ``the eyes of the mind, by which it
sees and observes things, are the proofs themselves'' (Eth. V,
23. schol., cf. 34. schol.). Memory and all the determinations
of the imagination are tied to the duration of the body
(duratio); but the mind expresses the body as actual not only
insofar as it conceives it in relation to time and place,
but also insofar as it understands it as contained in
God and following from the necessity of the divine nature\footnote{Cf. Eth. V, 36. schol.: quia nostræ mentis essentia in sola cognitione consistit, cuius principium et fundamentum Deus est, hinc
perspicuum nobis fit, quomodo et qua ratione mens nostra secundum essentiam et existentiam ex natura divina sequatur, et continuo a Deo pendeat.}, thus insofar as it understands it as real, as
having existence through God's eternal essence, i.e. \emph{as eternal}
(Eth. V, 29. schol. 30. dem.). So determined, the mind is
``an eternal mode of thought (æternus cogitandi modus), which is determined
by another eternal mode of thought, while this
in turn is determined by another, and so on to infinity,
so that all in unity (simul) constitute God's eternal
and infinite intellectus'' (Eth. V, 40. schol.). % print: opening » missing

Since the mind and the highest form of knowledge is eternal,
\emph{does not come to be} (V, 31. schol.), \emph{the intellectual
love} of God must also be eternal, and it must
be \emph{the only} eternal love. It is God's very
love, with which God loves himself (God loves
himself in that he, as infinite, is in possession of an infinite
perfection, accompanied by the idea of himself as
% --- p. 169 ---
\opage{169}\setcounter{footnote}{0}causa sui, V, 35. dem.), not as infinite, but insofar as
he is determined by the essence of the human mind considered
under the form of eternity, or, in other words:
the mind's intellectual love of God is a part of the
infinite love with which God loves himself (V,
36). \emph{Insofar, then, as God loves himself}, he loves
men, and God's love of men and the
mind's intellectual love of God are one and the
same\footnote{When in prop. 17 it is denied that God loves in the proper sense,
this denial applies to love as passio, and when it is
understood thus, it becomes untrue to demand love in return from God
(V, 19). But in intellectual love the mind is in unity
with God, does not relate itself to him as to another; it is therefore
one with God's immediate certainty of his infinite active being,
and the expression love designates something qualitatively other than in
that place.}. Since it follows with necessity from the essence of the
mind, insofar as this is considered as an eternal consequence of God's
nature, there can be nothing that is contrary to it,
or that can abolish it; for anything of the kind would be
an opposite of the true, and, in gaining power over
it, would make it the false\footnote{The relation here is thus other than with regard to single
things, which are considered in relation to a certain time and a certain place; for
of these there holds, by the nature of quantity, the axiom that there is in
nature no single thing without there being another
that is more powerful and stronger, and by which the former can be destroyed.
But the conception of things as single is at this point
abandoned.}.

Intellectual love is, as the mind's highest
expression, our \emph{freedom}\footnote{Freedom presupposes existence through itself (Eth. I, defin. 7, cf. I,
17. Coroll.), thus unity of essence and existence, the fundamental determination
of substance. When freedom is ascribed to the mind, it is thus
seen in unity with substance, so that the latter's determination can be
transferred immediately to it.} (V, 36. schol.) and our \emph{blessedness}
(beatitudo). Blessedness expresses not a transition to a
greater perfection, but an eternal possession of perfection
(Eth. V. 33. schol.). It is one with
% --- p. 170 ---
\opage{170} virtue itself, not a reward for it; only in the very active
power that blessedness expresses does our virtue, i.e. our active
capacity, consist (Eth. V, 42).

At this point the Spinozist Ethics and the Spinozist
system are concluded. Through the apparent
becoming in time, through the representation of the single finite
things, we have returned to the thought of the one,
infinite, eternal being, which \emph{is}, and does not \emph{become}, which, as
its own cause, posits itself in its inner fullness of eternal
modi. The same view that formed itself for the mind
in the conception of infinite extension, which, through
the laws of motion and rest infinitely modified, was eternally
like itself, comes forth here too in the domain of mind.
The eternal thought, which is the active essence of substance, is,
as infinitely modified by the ideas, eternally like itself.
What is limitation and negation has eternally vanished;
the incomplete, inadequate ideas are, seen from this standpoint,
true and adequate, and governed by the same eternal
law of necessity. Only the eternal is, and in the light of the eternal
what appeared in the form of time vanishes, and it
vanishes, not as something that \emph{has} been; but its being
is its vanishing. The human mind is, according to its
positive essence, intellectus, in unity with substance, enclosed
by God's eternal idea, by the intellectual love
with which he loves himself.

From this standpoint must be understood the propositions that
Spinoza, again approaching the conception of representation (of the imagination),
adds at the end of his Ethics (V, 38. ff.).
In these the mind is still determined as a \emph{composite}
of adequate and inadequate ideas; reason takes
its place in it \emph{alongside} the imagination; the eternal in
the mind becomes a greater or lesser \emph{part}; that which death
cannot touch is compared, as the \emph{more} perfect,
with the perishable, and the opposition between the wise man
and the ignorant (prop. 42. schol.) loses its unconditionedness
through limiting determinations. But in Spinoza's proposition
that the eternal has no relation to the finite,
% --- p. 171 ---
\opage{171} the corrective is given for that division of the mind; what is called
the mind's \emph{better part}, intelligence (Eth. IV, App. c. 32),
becomes the one real and active thing in it (V, 40. Coroll.),
and the quantitative comparison is rejected (ibid.).
And when the wise man is posited as ``sui et Dei et rerum \emph{æterna
quadam necessitate} conscius,'' the transition and the difference of degree
are again removed, eternal being affirmed, and the end of the system closed
together with its beginning.

\bigskip
\begin{center}\rule{2.5cm}{0.4pt}\end{center}
\bigskip

\begin{center}
{\large\textbf{Chapter Nine.}}\\[4pt]
\textbf{Characterization and Critique of Spinozism.}\\[6pt]
\rule{2cm}{0.4pt}
\end{center}
\addcontentsline{toc}{chapter}{Chapter Nine. Characterization and Critique of Spinozism}
\markboth{Characterization and Critique of Spinozism}{Characterization and Critique of Spinozism}
\bigskip

Through what was developed in the preceding sections the
moments are given for the complete determination of the peculiar
character of Spinozism. To establish this, one
must start from what was the system's essential
task, its inner mainspring, which never ceases to act
even where heterogeneous elements force their way in in the execution.
What is deviation from the fundamental thought becomes the object
not of the characterization of the system, but of the critique
of it.

We saw that the fundamental thought was the concept of substance in its
infinity, that the task of thought was to see it as \emph{all}
reality, that its inner impelling desire was energetically
to hold fast the thought of infinity as encompassing
everything. All reality is \textit{in} substance, substance is all
reality: that is the system's fundamental thought, and by it
the system is determined as \emph{pantheistic}. But this determination
does not yet distinguish it from many other systems of entirely
different origin; it becomes only a common designation,
which must receive its specific difference. This is given by what is
posited in the relation between the two attributes: the
% --- p. 172 ---
\opage{172}concept of necessity as the form of the being of substance. We saw that
the concept of substance had its origin in the concept of the
thinking I; it was, according to its provenance, the objectified
infinite essence of thought. But we saw at the same time that thought,
which had separated nature from itself as its opposite,
and which, in determining it abstractly as its opposite,
had no other determination for it than the abstract
concept of extension, had to seek its content in
nature; and since in the unity of substance it was again set in
unity with nature, but only in the unity of \emph{parallelism},
the law of nature abstractly conceived, the law of
necessity, had to become the form of thought, and, as the form of thought,
also the form of thought's objectified infinite essence.
Even though, therefore, the point of departure is originally taken from
self-consciousness, the form of the being of substance becomes
\emph{not the synthesis of self-consciousness, but the immediate
self-certainty of causality} in its necessity. Since substance,
causa sui, is conceived by thought in the whole
wealth of determinations which follow timelessly from its essence with eternal necessity,
in infinite modifications, the
form of this eternal unfolding, which is not separated from its
source, is the law of necessity. Substance is the necessary
\emph{cause} of necessary \emph{effects}, not the free originator of
freely posited \emph{ends}, and since in representation the effects
are separated from their cause, their infinite series comes to be governed
by the compelling law of external necessity.

The system's pantheistic fundamental character thereby acquires
a closer individualization. \emph{That which infinitely is
becomes the necessary power of nature, immediately self-certain
in its unity with thought}. ``God or nature'' expresses
this concept.

But with this determination of substance the
relation of human thought to it is given, and through
the character of thinking conditioned thereby, the character of the system
is again more closely determined. Thought is in substance
with an immediate certainty that finds its own essence again in substance.
Since limitation is only a negative, every real
% --- p. 173 ---
\opage{173} separation is removed, every scepticism excluded with regard
to the infinite. The infinite is the first thought of human
thought, its original certainty of itself.
And in the infinite it is in unity with its opposite:
extension. As expressions of the same substance the
modes of the two attributes run parallel; a scepticism with regard
to nature is thereby removed. In this immediate certainty
and security lies the system's character as \emph{dogmatism}.
Thought cannot turn critically against itself;
for it has no self subsisting for itself: it \emph{is} only in the
infinite and in it ``in unity with the whole of nature''. And since
it is \emph{thought}, as this determinate form of the human
essence, to the exclusion of other forms of the mind,
such as feeling and the like, that finds itself again in substance, and
is therefore determined as the only adequate expression of the
human essence, thought, reason, alone determines
the truth; it is the imperturbable revelation
of the divine, which includes certainty and is the sole
judge of the false and the true. Thereby the system acquires
the character of \emph{rationalism}: whatever does not congrue
with the universally valid determinations of thought, it rejects as
belonging to untrue and incomplete representation.

Finally, in the twofold determination of the concept of substance
and of thought's relation to it, the determination
of the form of the development of thought is given. The mind's form of existence is, like
that of substance, causality in its self-certainty: from the mind
as an active mode of thought there develop the adequate ideas linked together
with mathematical necessity, which form
a parallel to the series of things linked together by the law of natural
necessity. The unfolding of thought is thus in the form
of mathematical necessity; but since true
knowledge is in the element of eternity, the development
as development vanishes in the unity of substance, which in its eternal fullness
shows itself to \emph{intuitive knowledge}, which with one glance
sees substance in its infinity together with that which infinitely,
with eternal necessity, follows from its essence.

% --- p. 174 ---
\opage{174} The determinations \emph{pantheism} and \emph{naturalism} thus designate
the system from the objective side, the determinations
\emph{dogmatism} and \emph{rationalism} from the subjective; but
the two sides cannot be separated; the character of dogmatism
is, just as much as that of rationalism, given in the objective determinations.

With the character of the system thus established, it becomes
the basis for the task of critique. This task receives its limitation
from what is already historically given: philosophy
in its progressive development marks the one-sidedness of
the Spinozist system in integrating it. In
the presentation of Spinoza's relation to later philosophy
(see the next chapter) the results of this objective,
indirect, historical critique will be indicated. In this
chapter the task of critique will then be to judge Spinozism
from \emph{its own} standpoint, to investigate whether
it has reached the goal it has set itself, whether
it has been able to carry through its fundamental thought, whether inconsistencies
have not forced their way in, and whether these have not then been called forth with
necessity by the deficiency of the fundamental thought,
and therefore partly demonstrate its unfitness to
encompass what it wanted to encompass, partly raise the demand
for other fundamental determinations, whereby a transition is made to
the historical progressive development.

From the domain of this critique, then, every problem
is provisionally excluded that has not arisen within the fundamental view
given by the system. As long as thought consistently
unfolds without touching such a problem, the demand
on it to solve it is unfounded. If it succeeds
in closing itself into an inner whole without reaching
the problem, the problem is by that very fact rejected as unjustified.
Only through inconsistencies does thought lose this right of rejection.

Next, it follows of itself that such contradictions
in the system as are only \emph{apparent} do not become
objects of critique. Of such there are, precisely in
Spinoza, on account of his peculiar conception of
% --- p. 175 ---
\opage{175}\setcounter{footnote}{0}the finite, a great many, and it is precisely on them that,
e.g., the critique of Bayle and of most of Spinoza's contemporary opponents
relies. In the foregoing several such
have been pointed out, and the contradiction has been shown to be removed when
it is seen from Spinoza's standpoint; they will therefore be
passed over here.

Spinozism has set itself its task through its fundamental concept,
and this, we saw, was substance, conceived in its
strict infinity, and in the parallelism of its mutually
absolutely independent attributes. The task, then, is to think
the thought of infinity in such a way that finitude at no
point establishes itself over against it, while infinity nevertheless
acquires a fullness; and next: to think this fullness
as expressed in the mutually independent attributes, each of which
expresses the infinite wholly.

\emph{First}, then, the thought of infinity is to be thought as
\emph{rich in content}. As such it can only be thought when
substance posits in itself an other, a really subsisting
opposite, which is nevertheless taken up into its unity, bound by
it. But here the difficulty at once comes forth. The
infinite substance cannot posit an opposite in
itself; the form of its being is infinite causality, which,
as infinite, posits itself only as infinite, and immediately
is one as cause and effect. Limitation would
only be negation, but the infinite cannot negate
itself, and therefore cannot limit itself; by its very thought
the thought of finitude is excluded\footnote{Could thought derive
a subsisting finite from the infinite, then, according to
Spinoza's conception of the relation between thought and
its object, a corresponding development, and thus a relation of
\emph{time}, would have to be posited in the infinite being determined only by the form of \emph{eternity}.}. But the drive of human
thought to self-affirmation, to thinking, to philosophizing,
rests on the separation in existence between the finite and the
infinite: if the finite is absolutely removed in thought, the
possibility of human thinking is annihilated; to philosophize
becomes a self-contradiction. This forces out the
% --- p. 176 ---
\opage{176}\setcounter{footnote}{0}doubleness in the conception of what becomes the effect of substance.
Seen from the concept of substance, holding fast to the
proposition that limitation is negation, one gets only from
infinite to infinite. Substance as causa and as
causatum is in eternal oppositionless unity with itself.
But since separation is eternally removed, the concept sinks together
into itself without difference; human thought expires,
as it were, in its airless space. Therefore the need of thought
forces forth a fuller view: the infinite that
follows with eternal necessity from the nature of substance (Eth.
I, 16) becomes something infinitely modified in its unity with itself.
In the infinity of extension the
\emph{infinitely many} modi proceed according to the unchangeable
laws of motion and rest; in the infinity of thought the ideas proceed in
their infinity from the eternal law of reason; in both attributes
there forms an image of an eternally unchanged life moved within itself,
which expresses the same infinity and
remotely suggests an \emph{organic} conception. But the laws that
bring the multiplicity bound by the unity into infinity
remain uncomprehended\footnote{In the fact that the multiplicity
of the modifications cannot be deduced, and therefore stands as
uncomprehended, it immediately lies that finite consciousness,
inadequate knowledge, is not deduced either.
It becomes, like the subsisting finite, a presupposition
that dissolves itself in contradictions, and vanishes as it
goes back into the infinite.}, and since the thought of finitude
everywhere dissolves itself in contradictions (cf. p. 79, Note 2, p. 89,
121 and 136), and limitation becomes only something vanishing
--- and only thus can it be conceived where causality, not
end, is the fundamental concept, --- the fuller
image again vanishes in the indifference of abstraction. But even if it
constantly vanishes, it is nevertheless constantly posited anew,
for only through it does thought acquire movement, only through it does it avoid
denying itself. The general inconsistencies to which
this must lead have been brought out at the end
of the second chapter (p. 63 f.). In the more special developments
their consequences have been shown in detail. It was the
% --- p. 177 ---
\opage{177}\setcounter{footnote}{0}doubleness in the conception of the finite that disturbed the limit
between imagination and reason, between ratio and
scientia intuitiva, that gave universal knowledge a doubtful
truth, that made the concept of the common hover
between meaning the property common to really separate single things,
and the positive one that is the real in
the things only represented as single; it was this that let
the absolutely different\footnote{They relate to each other as
ens to non-ens, as infinite to finite.} adequate and
inadequate ideas be united in the human mind in an
uncertain neighborliness; that held thought fast in domains which,
strictly speaking, belong only to the imagination; and it was finally
this that, even after the single had been conceived under
the form of eternity and necessity, the form of being of the \emph{infinite},
held it fast as res \emph{particulares} æternæ. But
at this last point the contradiction breaks forth decisively.
Since the human mind in intuitive knowledge, which is
the eternal determination of its essence (cf. chapter 8), is posited as an
æternus cogitandi modus that is a part of God's infinite
intellectus (Eth. V, 40. schol.), either this modus,
which, even if it is eternal, is nevertheless single, must express the infinity
of substance, which is impossible, or it must vanish
as single in the infinite, and in its vanishing deny
itself and make substance empty of content. At this
point, therefore, there is a decisive pointing toward another determination
of the fundamental concepts. The conception of limitation
merely as negation made the infinite abstract and
brought the drive of thinking into contradiction with itself. This
points toward the demand to conceive limitation as
the self-limitation of the being, as its self-determination in
its self-active separation from others, --- toward the concept of the \emph{monad},
and with this concept Leibnitz integrates Spinoza's
fundamental thought (see the following chapter).

The \emph{second} task Spinozism had set itself
was to conceive the attributes of substance, thought and
% --- p. 178 ---
\opage{178}extension, as independent, in that they, without including each other's
concept, each by itself express substance, and thus to
derive everything that followed with necessity from the nature of the single
attribute from it alone, without regard to the other.
Under the same form of causality, thought and
extension should then have been conceived with the parallel series
of their modifications. But where the concept of imagination
and the passive states of the mind are developed, Spinoza starts
from the relations of the body; only through these can we form
a view of the essence of the confused, incomplete ideas;
in particular, the nature of memory is made perspicuous only by
the body's capacity to retain the impressions received, and Spinoza
here even comes to speak of the body's being acted upon
by external causes as the \emph{ground} of the confused ideas
(de emend. intell. p. 384), and of influences that hinder
the mind from thinking (Eth. V, 10). And on the other side,
where Spinoza develops man's freedom, his essence as
acting, as determined by reason and its adequate
ideas, the point of departure is taken from the mind alone; only from
it is the view formed of the ``imago'' that corresponds
to what is clearly and distinctly conceived (Eth. V, 12), and the mind
is now ascribed the power to order the affections of the body
(Eth. V, 10 schol.). The parallelism is thus not carried through;
there is wavering, as everywhere where the modes of action of the one absolute
are posited as constituting absolutely different
forms of existence, between a fusion and a hidden
dualism. In Spinoza, thought is on the one side determined in its
form by the form of the activity of nature, and in this
form it expresses, since in substance it is in unity with
extension, extension's immediate certainty of itself.
It is thus determined by extension both in form and in content;
for the determination for thought itself becomes to
be \emph{the immediately self-certain thought of extension}.
On the other side, by including in its
own concept the determination of its opposite, it reaches out
over that opposite, which was also originally determined only by the
abstract form of thought and has therefore received determinations
% --- p. 179 ---
\opage{179}\setcounter{footnote}{0}that conflict with its peculiar essence (e.g. the determination
as indivisible). This unconscious mutual determining
also sometimes brings Spinoza, as if involuntarily,
to apply concepts that are excluded by the fundamental determination,
thus the concept of \emph{end} (e.g. Eth. IV.
App. C. 4. Tract. polit. V, 2. De emend. intell. p. 360);
and the dualism that remains at bottom calls forth
the insoluble problem of the innumerable attributes and their
relation to thinking (cf. p. 85 f.).

But since Spinoza has thus not been able, at the two decisive points,
to solve the task he himself set,
critique acquires not only the right to declare the fundamental task
unsolved, but it gets a wider scope, since what is hinted at through
the inconsistencies becomes justified demands on
the system. It acquires the right to bring out \emph{against} Spinoza
\emph{the abstract character of the system}, which makes everything
sink together into the undifferentiated unity of substance, and,
under the semblance of existence that is left to the finite, lets
it be determined only by a \emph{quantitative} difference, in which
the energy of all oppositions is paralysed\footnote{Since the difference
becomes only quantitative, the determination
in Spinoza that an external influence is determined both by the nature
of what acts and of what is acted upon also loses its significance. The persistence
of individuality under quantitative change, under
the affections, becomes inexplicable.}, even the opposition
between good and evil, and all fullness vanishes,
since individuality becomes a relative, floating concept, which
is posited not by an inner peculiarity of essence, but by
the externally effected combination of the modifications of extension
and of the ideas corresponding to them (cf. chapter 5). It can
bring out --- what lies partly in this abstraction, partly in
the unconscious mutual determining of the attributes --- that no
peculiar determination can be given for the \emph{human
mind} as qualitatively different from the other modi,
though a demand for this is made in the system's point of departure,
and that there is therefore just as little room for a free
% --- p. 180 ---
\opage{180}\setcounter{footnote}{0}point of departure for self-determination, for an individual freedom,
for ends, as for the ethical conditioned thereby\footnote{Just as
physics became metaphysics, ethics becomes physics.}. It
can bring out that everything in the mind that is not determined by
thinking is excluded from its essence, and that therefore whole
spheres of the life of the mind, such as art, are wholly ignored,
others, such as religion, conceived only according to their relation
to the thought of the intellect. It can bring out that in the domain
of nature the dynamic element given in the principle is, within the limits
of the finite, overwhelmed by the mechanical, and that
the concept of the organic is made impossible by the oppositionlessness of substance.
It can finally bring out that the system
has not demonstrated the possibility of the concept with which
it concludes: the intellectual love of God proceeding from intuitive
knowledge; that the problem which is
given in the single mind's intuitive knowing --- a problem
which, thought through, would turn dogmatism into \emph{criticism}
--- is not solved, and that the system has only granted itself
a dispensation, on grounds of impossibility, from solving it. And critique must,
even if it concedingly holds fast to this concluding concept,
bring out its poor content, which is conditioned by the
mind's being determined as the idea of the body (of extension),
and extension only by quantity, motion and rest.
Intuitive knowledge, which gets only this content, even if in
the form of infinity, becomes abstract and poor, and the intellectual
love of God is touched by this poverty.
But on this point it is nevertheless important not to overlook what
so often happens with abstract systems: that the human spirit,
and still more perhaps the human heart, puts its richer content
into the abstract forms, and now through an illusion
sees them as full through that content. What strikes
Spinozism as a system does not strike Spinoza as a thinker.
In the exalted calm of his thought, in the energetic seriousness of his life,
in the love with which, through the purification of resignation,
he sank his thought and his will into the infinite
One, there speaks a deeper content, the drive of a richer being,
% --- p. 181 ---
\opage{181} than what he has been able to fix in the forms of thought, and their
deficiency must not lead us to overlook the former.

If we now, having let critique bring out what is
deficient in Spinoza's philosophy, wish to indicate his lasting
significance, we must not restrict it to what is independent of those
deficiencies: that he made the demand that
the substrate of nature too should be conceived as active, just as
he, in his theory of knowledge resting on the principles of existence,
energetically held fast to the activity of thinking, --- that
he in his doctrine of the affects has given a lasting foundation
for what one might call the natural history of the human mind,
--- and that he has asserted the autonomy of philosophy
and its independence of the religious, so that he, alone in his
time, wanted to be \emph{only} a thinker, and in his thinking had
his religion; but we must place it in this: that he, suggesting
the idea of an all-organism (cf. especially Eth. II, 13. schol. and
Ep. 15), against dualism, against the splitting up of the unity
of existence, wanted to think the all-encompassing thought of infinity,
to seek the deeper ground in which the infinite essence of the human
mind is in active unity with the infinite essence of nature.
And although he has determined this ground of unity
only as consuming the oppositions, yet
the thought that was the inner impelling power of his philosophy
is the ground of all thoughts, and, even though it takes up other
determinations into itself, it nevertheless remains that which bears these, as
the infinite existing thought, encompassing them, bears nature and
the human spirit, which have their ground in it.

\bigskip
\begin{center}\rule{2.5cm}{0.4pt}\end{center}
\bigskip

\begin{center}
{\large\textbf{Chapter Ten.}}\\[4pt]
\textbf{Spinoza's Relation to Later Philosophy.}\\[6pt]
\rule{2cm}{0.4pt}
\end{center}
\addcontentsline{toc}{chapter}{Chapter Ten. Spinoza's Relation to Later Philosophy}
\markboth{Spinoza's Relation to Later Philosophy}{Spinoza's Relation to Later Philosophy}
\bigskip

Spinoza's system was in its whole character far too alien
to the common consciousness, and far too sharply opposed to the
% --- p. 182 ---
\opage{182} general direction of its time, for it to have been able to gain
an immediate, thoroughgoing influence. The disappearance of individuality
in the infinite, the annihilation of every anthropomorphic
representation of God, the transference of the law of natural necessity
to his essence and to his effects,
the removal, conditioned by the pantheistic fundamental view,
of the opposition between an existing good and evil, the conception
of the religious fundamental determinations as mere forms of representation,
gave too strong offence to the religious consciousness,
which Protestantism had taught to accentuate individuality
and, in the struggle for salvation, to sharpen the opposition
between good and evil. Spinoza's thoughts terrified,
repelled, were not understood. In mysticism's sinking into the
divine the heart had found its satisfaction;
it was the God of the \emph{heart} to whom the mystic gave himself up,
and with whom, in this surrender, he still preserved
himself. But Spinoza's infinite is the infinite of \emph{thought}, which
gives the feelings no compensation for the surrender.

It is therefore not to be wondered at that Spinoza found
only few contemporary disciples. Of only one do we have a
writing in Spinoza's spirit, namely \emph{Abrah. Joh. Cuffeler's}:
specimen artis ratiocinandi naturalis et artificialis ad pantosophiæ
principia manuducens (1684) --- a kind of logic. Another
writing, by the clergyman \emph{van Leenhoff}: Hemel op aarden
(Heaven on Earth), which on its appearance (1703) aroused
much controversy, contains, so far as one can judge from
the reports in Bayle and Brucker, only a doctrine of the independent
value of virtue, and of resignation under the law of comprehended
necessity, which, even if influenced by Spinozism,
does not go back to its first principles.

But even if Spinoza found only few adherents,
Spinozism nevertheless soon became, as the name for something not understood and
abhorred, a heretic's name with which people branded
one another, and much free-thinking has probably also given
itself a more distinguished appearance by appropriating this name.
Therefore an approximately contemporary author, \emph{Roellius} (de
religione naturali, § 151, pag. 166, cited by Brucker),
% --- p. 183 ---
\opage{183} can say: Spinozam tota armenta in Belgio sequi ducem. But
from several of the polemical writings against Spinoza still preserved,
which, although attacking him, brought upon their authors the accusation
of being his secret adherents, one sees
how vague and indeterminate the conception of the character of Spinozism was,
just as the whole polemical literature
without exception shows that Spinoza was constantly misunderstood
precisely in the most essential points: in the conception of
the relation of finitude to God and of the form of being of the divine.
From the representation of the finite as
really existing, and of God as acting arbitrarily, ---
representations that could not be given up without giving up the
whole foundation on which one stood, people fought against
Spinoza, and drew absurd consequences from his doctrine.

It therefore had to be the fate of Spinozism to take effect only
late: it has been like a kernel enclosed in
hard shells, which must lie long in the earth before it can
sprout, but which then also unfolds in all the more vigorous
growth.

On only one of Spinoza's contemporaries has his thinking
had an influence that has acquired significance for history;
but this influence has been through opposition,
as an incitement to an independent development of thought that
forms an integrating opposite to Spinoza. This
contemporary is \emph{Leibnitz}. It is true that Leibnitz has only seldom expressed
himself on his relation to Spinoza, of whom he speaks on the whole
with diplomatic caution, and still less put
specific points of his doctrine in a direct connection with
Spinoza's; but there are nevertheless points of comparison enough present
that can both form a parallelism illuminating for both systems
and at the same time, while they demonstrate Leibnitz's peculiarity,
show a similarity to his predecessor that can
hardly everywhere be explained without a definite influence.
Where the similarity lies in the whole character of the philosophy of the age,
nothing speaks for a dependence; thus when
Leibnitz meets Spinoza in conceiving ideas as the mind's
original property; when he lets mechanism,
% --- p. 184 ---
\opage{184} to which he otherwise assigns a definite domain, assert itself
also within the limit of the mind and of the organic,
and when he asserts the significance of philosophy in relation
to the religious. But in a more decisive way
Leibnitz's determination of the essence of substance as
force recalls Spinoza's conception of it as causality; ---
Leibnitz's fundamental concept, the \emph{monad}, recalls the \emph{Corpora simplicissima}
hypothetically assumed in Spinoza's physics
(cf. Eth. II, 13), particularly when, with the determination of
their relation to the universal individuum, one compares
the determinations about the necessary harmony between the parts
and the whole (Ep. 15), which form a parallel to the harmonia præstabilita
given in the concept of the monad; --- the concept decisive for Leibnitz's
whole doctrine, \emph{the confused representations},
which always involve infinity, recalls
Spinoza's conception of the mind as composed of an infinite
number of ideas, corresponding to the body's infinite composition,
and of these ideas as existing only inadequately in
the single mind (as ideæ confusæ); --- the juxtaposition
of the clear and distinct ideas with acting, of the confused
with passion, and the referring back of acting to the mind, of
passion to the body, is the same in both systems; likewise
the parallelism between soul and body, the significance of the body's perfection
for the soul, the conception of the whole universe
as containing life, of difference as difference of \emph{degree}, % print: Forkjel (misprint for Forskjel)
and of death and coming-to-be as a continuous transition.

But through this similarity, which points to a kinship,
the fundamental difference comes forth. When Leibnitz conceives
the essence of substance as activity, this
activity is immediately determined as \emph{self-}activity, in the sense
of \emph{individual} determining and separating from others. The
one infinite substance becomes the \emph{infinitely many}
monads, which, limited \emph{by themselves}, in their independence
express by their limit the ideal relation
to all the others, and in an individual way include the whole
universe within themselves. But Leibnitz's conception of activity
as \emph{self-activity} again stands in connection with the fact that
% --- p. 185 ---
\opage{185} force is conceived as a \emph{spiritual} principle, that monad
becomes synonymous with \emph{soul}, and that \emph{the form of the mind,
in decisive distinction from Spinoza, becomes the synthesis of self-consciousness},
not the immediate self-certainty of causality.
And this conception of the form of the mind, which is
a logical consequence of conceiving it as originally
producing ideas that are in it, and constitutes a specific
difference between its mode of being and that of matter, now conditions
the conception of God's essence, of the finite, of the mind's
acting, and of the relation between thought and extension.
In God's essence it is not the law of mathematical necessity
that becomes determining, but a hypothetical, moral necessity:
a necessity conditioned by regard for the best
determines the will. The concept of end again acquires its validity
for God and the finite; it becomes the explanatory
presupposition for the efficient causes, for mechanism.
The coming-to-be of the finite becomes not an effect of necessity,
but of God's freedom, on which the existence of things
depends, while their essence is given with eternal unchangeable
necessity in God's. The finite acquires a
centre for self-activity, a freedom in dependence,
a self-determination that also makes evil possible.
The human mind becomes no longer a modus cogitandi,
but an imperishable principium cogitationis. And
extension now becomes no longer an attribute, not something
real that in its infinity expresses the divine;
it becomes merely the phenomenon of the composition of the single substances.
What is lasting, real, active, is only the
single substances, the monads, whose essence is spiritual;
body, matter, extension become the passive, the non-real,
which has its subsistence only in confused representation. In
one word, then: individuality, which gives the infinite
total organism fullness, breaks forth with the form of self-consciousness,
and what is not expressed by this form, the material,
loses its rank as the parallel of the mind, and becomes
phenomenon, even though a parallelism is mirrored in the phenomenon.

% --- p. 186 ---
\opage{186} Thus Leibnitz himself expresses, in relation to Spinoza,
the peculiar nature of his system's fundamental concept, the monad.
Influence is converted into independent production; what is given
from without becomes an incitement for creative peculiarity.
And by this peculiarity he integrates
Spinoza.

In the period between Leibnitz and the last decades
of the previous century Spinoza was almost forgotten. \emph{Wolf}
polemicized against him (in his theologia naturalis), but with
half concepts and without understanding him, and the direction of the age
was unfavourable to a deeper thinking. Toward the end of the century
\emph{French materialism} invoked
Spinoza as an authority. In his concept of God it saw
only the concept of nature, in his doctrine of necessity
it found its determinism, in his doctrine of the mind's relation
to the body its sensualism. The idealistic foundation
in him it overlooked; it did not bring those determinations
into connection with the whole of his views, but appropriated
all the more eagerly the details of his polemic
against the theological representations, and of his critique.

In Germany Spinoza at the same time found a kindred spirit
\emph{in Lessing}, who could find no rest in the spiritless
philosophy of his time, and, since he did not himself have the power to produce
one richer in content, sought back to the past. The
controversy that was aroused when \emph{Jacobi} revealed the secret of
Lessing's Spinozism is well known. It became apparent in this
controversy that Jacobi, who, like Leibnitz, was by his peculiarity
opposed to Spinoza, but in his genius had the
sympathetic understanding of the opposite, was the only one
who had penetrated more deeply into Spinoza's thoughts, while \emph{Mendelssohn}
and \emph{Herder}, the former even without the external
knowledge of Spinozism, sought, through a misunderstanding of
it, to bring it together with the trivialized theistic concept
of the Leibnitzian theodicy.

This controversy about Spinoza called the study of him back
to life, and the new direction of philosophy soon gave it a decisive
significance. When \emph{Kant} had made the
% --- p. 187 ---
\opage{187}\setcounter{footnote}{0}unproved presupposition of dogmatism, the possibility of knowledge, into a problem,
and in the attempt at its solution had given the premisses
from which \emph{Fichte} drew the consequences by placing
all reality in self-consciousness, the dialectic of the latter
again forced forth the need of an objective supplement. Fichte
himself had, at the later stage of development of his philosophy,
to seek such a one, and could then seek it only in the philosophy
that was the sharpest opposite of the system of self-consciousness:
in Spinozism with its absolute \emph{being}. The
same need drove \emph{Novalis} and \emph{Schleiermacher} to an
enthusiastic attachment to Spinoza, which nevertheless did not annihilate
their peculiarity. It finally found its most adequate
expression in \emph{Schelling}, when, in place of the
subjective principle of the Fichtean philosophy, self-consciousness,
he posited the subjective-objective, \emph{reason}, the absolute, which
is the indifference of the subjective and the objective, the absolute
identity, which is never sublated as identity, so that
everything that is, is this absolute identity itself, that there
is nothing finite, but everything is infinity itself, which only
in its form of being is touched by the quantitative difference
between preponderant subjectivity and objectivity, but in its
essence is and remains absolute totality. Through the fundamental concept
of the system of absolute identity the fundamental concept of Spinozism
thus acts, despite the different historical presuppositions,
and Schelling rejects Spinozism as ``one-sidedly
realistic''\footnote{\emph{Schelling}: philos. Schriften I, p. 417 ff.} only when, through his ``positive'' philosophy, he makes an
attempt to put theism in the place of pantheism.

\emph{Hegel's} philosophy too takes Spinozism as its foundation
and expressly acknowledges it as such\footnote{With regard to Hegel's relation to Spinoza, cf. \emph{Hegels}
sämmtl. Werke. Bd. I, p. 60---79. III, 117 f. 178. 396 ff. 466.
IV, 194 f. V, 9 ff. VI, Preface S. XVII. p. 301. 303. 310 f.
313. 391. 403. XII, 442---449. XV, 336 ff. 341 ff. 359. 362 f.
365 ff. XVII, 7 ff.}. But
the Spinozist concept of substance, in which Hegel does not find
% --- p. 188 ---
\opage{188}\setcounter{footnote}{0}the determination of activity, he demands be completed by
the form of subjectivity, so that through this it can be torn out
of its abstraction and immobility, and be conceived as \emph{spirit}
and as \emph{freedom}, as ``absolute \emph{person}'', not merely as
the absolute \emph{object} (``die absolute Sache'').

Thus Spinoza's influence still reaches into the
system that stands as the last historically concluded one. Spinozism
acts here through homogeneity, as something kindred to
the ruling philosophical direction of the age, and
as something thus still active it becomes the object of passionate
and frequently misunderstanding polemic from the thinker
who, like the whole of the most recent philosophy, acknowledges Kant,
though only not the idealistic centre of his philosophy, as
the point of departure of his thinking, --- from \emph{Herbart}\footnote{Cf. e.g. \emph{Herbart}: Allgemeine Metaphysik, I, p. 128---190 and
many places in the second part. Analytische Behandlung des Naturrechts
und der Moral p. 32---43. p. 242 and elsewhere.}.

Through this influence on the philosophy of our time
the significance is also justified that has been ascribed in this work to
Spinoza as representative of the starting point of
modern idealism. Spinoza and Hegel mark the two end points;
Hegel's philosophy has been possible only through Spinoza;
the deeper motives of Spinoza's philosophy are first understood
fully through Hegel.
\end{document}
